\documentclass[a4paper,11pt,reqno]{amsart}
\usepackage[utf8]{inputenc}
\usepackage{enumerate}
\usepackage{amsmath, amssymb,amsthm}
\usepackage{mathtools}
\usepackage{leftidx}
\usepackage{setspace}
\numberwithin{equation}{section}
%\usepackage{mathbbol}
\setcounter{tocdepth}{1} %2=show subsections, 1= show only sections
\usepackage[left=3cm, right=3cm]{geometry}
\usepackage{hyperref}
\usepackage{xcolor}
%\definecolor{name}{model}{color-spec}
\definecolor{my-black}{rgb}{0,0,0}
\definecolor{my-blue}{rgb}{0,0,0.8}
\definecolor{my-red}{rgb}{0.8,0,0} 
\definecolor{my-green}{rgb}{0,0.5,0}
\hypersetup{colorlinks,
    linkcolor	= {my-blue},
    citecolor	= {my-red},
    urlcolor		= {my-black}
}
\usepackage{graphicx}
\usepackage{tikz-cd}
\onehalfspacing
%Javier's suggestion to avoid overfull boxes
%\overfullrule=10cm better not to use because ArXiv might mess up things
%\usepackage{microtype}

\theoremstyle{plain} %italic

\newtheorem{lemma}{Lemma}[section]
\newtheorem{theorem}{Theorem}
\newtheorem{corollary}{Corollary}[section]
\newtheorem{proposition}{Proposition}[section]
\newtheorem{assumption}{Assumption}[section]
\newtheorem{problem}{Problem}[section]
\newtheorem{consequence}{Consequence}[section]
\newtheorem*{theorem*}{Theorem} %to avoid numbering


\theoremstyle{definition} %non-italic
\newtheorem{remark}{Remark}[section]

\newtheorem{definition}{Definition}[section]
\newtheorem{notation}{Notation}[section]
\newtheorem{conjecture}{Conjecture}[section]
\newtheorem{example}{Example}[section]
\newtheorem{claim}{Claim}[section]

\theoremstyle{remark}
\newtheorem*{remark-non}{Remark}

\DeclareMathOperator{\id}{id}
\DeclareMathOperator{\real}{Re}
\DeclareMathOperator{\imag}{Im}
\DeclareMathOperator{\image}{im}
\DeclareMathOperator{\supp}{supp}
\DeclareMathOperator{\vol}{vol}
\DeclareMathOperator{\inv}{inv}
\DeclareMathOperator{\Tr}{Tr}
\DeclareMathOperator{\ord}{ord}
\DeclareMathOperator{\Hom}{Hom}
\DeclareMathOperator{\Aut}{Aut}
\DeclareMathOperator{\End}{End}
\DeclareMathOperator{\Hol}{Hol}
\DeclareMathOperator{\GL}{GL}
\DeclareMathOperator{\SL}{SL}
\DeclareMathOperator{\PGL}{PGL}
\DeclareMathOperator{\PSL}{PSL}
\DeclareMathOperator{\Spur}{Sp}
\DeclareMathOperator{\PSO}{PSO}
\DeclareMathOperator{\SO}{SO}
\DeclareMathOperator{\Or}{O}
\DeclareMathOperator{\Res}{Res}
\DeclareMathOperator{\diam}{diam}
\DeclareMathOperator{\sgn}{sgn}
\DeclareMathOperator{\sinc}{sinc}
\DeclareMathOperator{\disc}{disc}


\newcommand{\R}{\mathbb{R}}
\renewcommand{\P}{\mathbb{P}}
\newcommand{\C}{\mathbb{C}}
\newcommand{\N}{\mathbb{N}}
\newcommand{\Q}{\mathbb{Q}}
\newcommand{\Z}{\mathbb{Z}}
\newcommand{\F}{\mathbb{F}}
\newcommand{\g}{\mathfrak{g}}
\renewcommand{\o}{\mathfrak{o}}
\renewcommand{\d}{\mathfrak{d}}
\renewcommand{\a}{\mathfrak{a}}
\renewcommand{\b}{\mathfrak{b}}
\newcommand{\p}{\mathfrak{p}}
\newcommand{\I}{\mathbf{I}}

\renewcommand{\sl}{\mathfrak{sl}}
\newcommand{\h}{\mathfrak{h}}
\renewcommand{\H}{\mathbb{H}}
\textwidth 16cm
\oddsidemargin  0cm \evensidemargin 0cm \parindent0.7em
\textheight 23cm
\allowdisplaybreaks


\title[Time-frequency localization]{Eigenfunctions, Free Boundaries, and Time-Frequency Localization}
\author{Jo{\~a}o P.G. Ramos}



\begin{document}
\maketitle 
\begin{abstract}
We study Gaussian time--frequency localization operators through their Bargmann--Fock representation and develop an inverse and free-boundary approach to their spectral theory.

We develop a perturbative inverse theory showing that functions sufficiently close to the ground-state Hermite function can be realized as eigenfunctions of localization operators. As a byproduct, we recover the classical disk characterization of Abreu and D\"orfler in the simply connected setting by the same free-boundary method. As an application we prove that the stability estimate of Gomez-Guerra, Ramos and Tilli is sharp.

In the second part of the paper we apply these inverse results to Faber--Krahn-type problems for Gaussian time--frequency concentration. We prove that local maximizers of the first localization eigenvalue under a measure constraint are necessarily disks, and we establish existence of global maximizers via a concentration--compactness argument in the Bargmann--Fock space, recovering the Nicola--Tilli theorem from this perspective.
\end{abstract}


% ==============================================================
% Section 1: Introduction
% ==============================================================

\section{Introduction}

\subsection{Time--frequency localization and the Gaussian model}

The study of time--frequency concentration lies at the heart of modern harmonic analysis, with deep connections to signal processing, mathematical physics, and phase--space methods; see, for instance, \cite{Grochenig,Mallat,Folland}. A central object in this theory is the \emph{short-time Fourier transform} (STFT), defined for a function $f \in L^2(\mathbb R)$ and a window $\varphi \in L^2(\mathbb R)$ by
\[
V_\varphi f(x,\omega) = \int_{\mathbb R} f(t)\,\overline{\varphi(t-x)}\,e^{-2\pi i t \omega}\,dt.
\]
This transform encodes the joint time--frequency content of $f$, and its squared modulus $|V_\varphi f|^2$ provides a phase--space energy density. Understanding how this energy can be concentrated in regions of the time--frequency plane is a fundamental problem, with both theoretical and applied significance.

In applications ranging from signal recovery to uncertainty principles, one is often interested in maximizing the quantity
\[
\int_\Omega |V_\varphi f(x,\omega)|^2\,dx\,d\omega
\]
over functions $f$ of unit norm and over sets $\Omega$ of prescribed measure. This leads naturally to the study of \emph{time--frequency localization operators}, introduced by Daubechies in \cite{Daubechies} and further developed in \cite{Daubechies-Paul,Ramanathan-Topiwala-Weyl,Ramanathan-Topiwala-Spectrogram,Cordero-Grochenig,Boggiatto-Cordero-Grochenig,Bayer-Grochenig,Kisil-CrossToeplitz,Nicola-Tilli-2,Ramos-time-frequency}. For general windows, sharp concentration bounds are notoriously elusive, and Lieb's uncertainty principle \cite{Lieb} remains a fundamental benchmark; see also \cite{abreu2021donoho} for related large-sieve phenomena in modulation and polyanalytic Fock spaces. In the Gaussian setting, by contrast, these operators acquire a particularly rigid structure: via the Bargmann transform \cite{Bargmann}, they can be realized as Toeplitz-type integral operators on the Bargmann--Fock space, and the problem becomes one of complex analysis in weighted spaces; compare \cite{Berezin,Berger-Coburn-Symbol,Berger-Coburn,Berger-Coburn-Heat,Zhu,Grochenig,Folland}. In this picture, the STFT is intertwined with an entire function whose weighted modulus encodes the concentration properties of the signal, and geometric questions translate into analytic ones.

\subsection{Spectral and geometric interplay}

A central theme in this area is the interplay between the \emph{spectral properties} of localization operators and the \emph{geometry} of the underlying domains. In the Gaussian case, Nicola and Tilli proved the sharp Faber--Krahn inequality for the STFT, showing that disks maximize concentration among all measurable sets of a fixed measure \cite{Nicola-Tilli}; see also \cite{Nicola-Tilli-2,Bastianoni-Cordero-Nicola,Bastianoni-Teofanov} for complementary spectral and regularity results on localization operators. The stability theory for this variational problem was subsequently developed in \cite{Gomez-Guerra-Ramos-Tilli}, where quantitative control of both the first eigenfunction and the Fraenkel asymmetry is obtained from the spectral deficit. Related concentration questions for other transforms, including wavelet transforms, have been investigated in \cite{Ramos-Tilli}.

From a complementary perspective, one may pose inverse questions: to what extent does spectral information determine the geometry of the localization domain? Conversely, given a function, can one construct a domain for which it is an eigenfunction of the associated localization operator? The first systematic result in this direction is due to Abreu and D\"orfler \cite{Abreu-Doerfler}, who showed that if a Hermite function is an eigenfunction of a localization operator on a simply connected domain, then the domain must be a disk. This inverse viewpoint sits naturally within the broader circle of ideas surrounding weighted quadrature identities, free-boundary problems, and inverse potential theory: classical results of Aharonov, Schiffer, and Zalcman \cite{Aharonov-Schiffer-Zalcman} characterize disks via mean-value identities for harmonic functions; Karp \cite{Karp} and Cherednichenko \cite{Cherednichenko} address inverse problems for Newtonian potentials; Gustafsson and collaborators \cite{Gustafsson,Gustafsson-Shahgholian,Shahgholian} develop the theory of quadrature domains and obstacle problems; and Isakov's monograph \cite{Isakov} provides a perturbative framework for related inverse source problems. The inverse viewpoint is also tied to the special role of Gaussian decay and Hermite structures in the Bargmann model; see \cite{Kulikov-Oliveira-Ramos,Goncalves-Oliveira-Steinerberger,Thangavelu}.

\subsection{Eigenfunctions as geometric data}

The guiding principle of the present paper is that, for Gaussian localization operators, eigenfunctions should be regarded as \emph{geometric data}. More precisely, in rigid situations the eigenfunction should determine the domain, and the ability to prescribe a near-Gaussian eigenfunction should in turn have strong consequences for nearby variational problems. The three main threads of the paper are organized around this idea:
\begin{enumerate}
\item rigid eigenfunctions force rigid domains;
\item near the Gaussian, one can prescribe the eigenfunction and construct the corresponding localization domain;
\item once such an eigenfunction-to-domain mechanism is available, it can be used to analyze sharpness and extremizers in concentration inequalities.
\end{enumerate}

The present work contributes to this line of research by developing a unified inverse and free-boundary framework for Gaussian time--frequency localization operators. Our approach builds on, and extends, the rigidity result of Abreu and D\"orfler \cite{Abreu-Doerfler}, the variational theory of Nicola and Tilli \cite{Nicola-Tilli,Nicola-Tilli-2}, and the quantitative stability theorems of G\'omez, Guerra, Ramos, and Tilli \cite{Gomez-Guerra-Ramos-Tilli}. At the technical level, our arguments also draw on ideas from inverse potential problems and quadrature-domain theory, in the spirit of \cite{Isakov,Cherednichenko,Gustafsson,Gustafsson-Shahgholian,Shahgholian}. In this way, the paper places localization operators squarely within the broader free-boundary and inverse-potential framework suggested by the Gaussian Bargmann model.

Despite these advances, two fundamental questions have remained open. First, no general construction of localization domains associated with prescribed near-Gaussian functions has been available, beyond a handful of highly symmetric examples. Second, the role of eigenfunction rigidity in variational concentration problems has not been clarified, leaving the sharpness of recent stability inequalities and the structure of local maximizers out of reach. The results of the present paper address these questions in a unified manner.

\subsection{An inverse construction near the Gaussian}

Our first main result is the constructive half of this picture. It shows that, near the Gaussian, one can indeed prescribe the eigenfunction and realize it through a suitable localization domain.

\begin{theorem}\label{thm:perturb}
Let $\lambda\in (0,1)$ and $N\in\mathbb N$. There exists $\varepsilon_0(\lambda,N)>0$ such that if
\[
f_0 = h_0 + \sum_{k=1}^N a_k h_k
\]
satisfies $\max_k |a_k|<\varepsilon_0$, then there exists a domain $U_\lambda$ such that
\[
\mathcal L_{U_\lambda} f_0 = \lambda f_0 .
\]
\end{theorem}

This theorem provides, to the best of our knowledge, the first systematic construction of localization domains associated with nontrivial linear combinations of Hermite functions. In particular, it supplies the inverse counterpart to the Abreu--D\"orfler rigidity phenomenon: rather than deducing the domain from a rigid eigenfunction, we construct domains starting from prescribed near-Gaussian eigenfunctions. The dependence of $\varepsilon_0$ on $\lambda$ and $N$ is explicit in our argument, and reflects the size of the spectral gap of the disk localization operator on the relevant Hermite subspace.

The proof is a direct perturbation argument at the Daubechies disk. After
applying the coefficientwise conjugate differential operator, we reformulate
the problem as an infinite system of weighted moment equations for an analytic
radial graph over the disk. At the disk the linearized moment map is diagonal
in Fourier modes, with an explicit inverse. An implicit-function theorem in
analytic boundary spaces produces a nearby domain satisfying the symmetrized
identity exactly, and a final elementary ODE argument recovers the original
eigenvalue equation.

A noteworthy aspect of this construction is that it reaches well beyond monomial-type eigenfunctions, the only case previously accessible by direct symmetry arguments. Our theorem shows that a much broader class of near-Gaussian functions can be realized spectrally, opening the door to a richer inverse theory and revealing new connections with nonlinear potential theory and weighted quadrature domains in the spirit of \cite{Gustafsson-Shahgholian,Shahgholian}.

\subsection{A new proof of the Abreu--D\"orfler rigidity theorem}

The same free-boundary framework also yields the rigid half of the correspondence, namely a short recovery of the classical theorem of Abreu and D\"orfler.

More precisely, we provide two new proofs of the Abreu--D\"orfler rigidity result based on our methods: one through the local radiality mechanism for the logarithmic potential, and a second through a holomorphic auxiliary function attached to the associated Poisson problem.

\begin{theorem}\label{thm:eig-loc-U}
Let $U$ be a simply connected bounded domain. If a Hermite function is an eigenfunction of the localization operator associated with $U$, then $U$ is a disk.
\end{theorem}

The proof transports the eigenvalue equation to the Bargmann--Fock space, where Hermite functions correspond to monomials and the localization operator becomes an integral operator with reproducing-kernel structure. In this setting, the eigenvalue equation translates into a functional identity involving exponential generating functions. Differentiating appropriately and invoking a Paley--Wiener type argument, we reduce the problem to the vanishing of a Fourier transform along complex lines.

% TODO: add citation for Ehrenpreis--Malgrange division theorem (no bib entry exists yet)
This allows us to apply a divisibility lemma of Ehrenpreis--Malgrange type, which produces a logarithmic potential solving a Poisson equation with source supported on the domain. The decisive insight is that the spectral assumption imposes strong constraints on this potential: it must exhibit a form of radial behavior on the boundary. In the simply connected setting, this constraint forces the domain to be a disk, recovering the Abreu--D\"orfler rigidity theorem from the same inverse/free-boundary perspective.

\subsection{Sharpness of the GGRT stability inequality}

We now turn to consequences of this inverse point of view. Our first application concerns the quantitative stability theorem of G\'omez, Guerra, Ramos, and Tilli for the Faber--Krahn inequality for the Gaussian STFT. In the formulation most relevant here, if $U\subset \R^2$ is a measurable set of finite positive measure, $\lambda_1(U)$ denotes the first eigenvalue of the associated localization operator, and $f_U$ is a corresponding $L^2$-normalized first eigenfunction, then Corollary~1.4 of \cite{Gomez-Guerra-Ramos-Tilli} gives the bounds
\[
\min_{z_0\in\C,\ |c|=1}\|f_U-c\varphi_{z_0}\|_{L^2(\R)}
\leq
C e^{|U|/2}
\left(
1-\frac{\lambda_1(U)}{1-e^{-|U|}}
\right)^{1/2},
\]
and
\[
\mathcal A(U)
\leq
K(|U|)
\left(
1-\frac{\lambda_1(U)}{1-e^{-|U|}}
\right)^{1/2},
\]
where $\varphi_{z_0}$ ranges over the Gaussian optimizers and $\mathcal A(U)$ denotes the Fraenkel asymmetry of $U$.

Our perturbative inverse theorem provides the missing input needed to test this stability estimate on genuine spectral data. Starting from a prescribed polynomial perturbation of the Gaussian, Theorem~\ref{thm:perturb} produces an exact eigenpair $(U_\varepsilon,f_\varepsilon)$ for the localization operator, with $U_\varepsilon$ a smooth perturbation of the disk and $f_\varepsilon$ a corresponding first eigenfunction. The constructive direction of the paper thus yields explicit near-extremizers, against which one can compare the size of the spectral deficit with the actual geometric displacement of the domain from the class of disks. The resulting one-parameter family shows that the square-root dependence in Corollary~1.4 of \cite{Gomez-Guerra-Ramos-Tilli} is the best possible.

\begin{theorem} \label{thm:GGRT-sharp}
The exponent $1/2$ in Corollary~1.4 of \cite{Gomez-Guerra-Ramos-Tilli} is optimal. More precisely, the powers of the deficit
\[
1-\frac{\lambda_1(U)}{1-e^{-|U|}}
\]
appearing in both the eigenfunction estimate and the Fraenkel-asymmetry estimate there cannot be replaced by $\beta$ for any $\beta>1/2$.
\end{theorem}

The argument rests on constructing a carefully chosen family of perturbations of the disk via the inverse theorem above. Starting from a polynomial perturbation of the ground-state Hermite function, we obtain a family of domains whose geometry deviates from that of a disk at first order. A detailed asymptotic analysis of the eigenvalue equation, including a first-variation computation by a transport formula, then identifies the leading Fourier modes of the boundary deformation.

This analysis shows that the deformation is governed by second-order harmonics, which cannot be eliminated by translations of the disk. Comparing the perturbed domains with arbitrary disks then yields lower bounds on their symmetric difference that match the upper bounds in the stability estimate of \cite{Gomez-Guerra-Ramos-Tilli}. The proof combines spectral perturbation theory, geometric measure estimates, and harmonic analysis on the circle, illustrating how inverse constructions can be used to probe the sharpness of variational inequalities.

\subsection{Local maximizers are disks}

We next connect this eigenfunction-to-domain philosophy to variational problems. Our fourth main result concerns the structure of local maximizers for Gaussian time--frequency concentration.

\begin{theorem}
Any local maximizer of $\lambda_{1,\varphi}$ under a measure constraint is, up to translation, a disk.
\end{theorem}

The proof relies on combining variational arguments with the inverse spectral perspective developed earlier. We first show that any local maximizer must coincide with a superlevel set of the density associated with its first eigenfunction in the Bargmann--Fock space. This reduces the geometric problem to a functional one.

We then analyze the spectral properties of the localization operator at a maximizer, showing that the first eigenvalue is simple and that derivatives of the eigenfunction remain eigenfunctions. This strong rigidity property forces the maximizing eigenfunction into a Gaussian form. Once this is known, the corresponding domain is determined as well, and one concludes that the only admissible configuration is a disk. The argument thus shows that extremizers are governed by the same rigidity mechanism that underlies the inverse theory.

\subsection{A new proof of the Nicola--Tilli theorem}

Finally, we address the global variational problem and recover the Nicola--Tilli theorem from the same perspective.

\begin{theorem}[New proof of Nicola--Tilli]
Disks maximize Gaussian time--frequency concentration among all sets of fixed measure.
\end{theorem}

Our proof is based on a concentration--compactness argument in the Bargmann--Fock space. We consider maximizing sequences and analyze their possible limiting behavior, ruling out vanishing and dichotomy scenarios using the structure of the Fock space and the invariance properties of the STFT. This yields the existence of a maximizer. Related concentration--compactness methods have recently been used in nearby time--frequency concentration problems, notably by Nicola, Romero, and Trapasso for ambiguity-function type concentration functionals \cite{Nicola-Romero-Trapasso}; see also \cite{Marceca-Romero-Speckbacher,Samuelsen} for recent work on Fourier concentration operators, and \cite{Fulsche-Hagger} for closely related developments in the polyanalytic Fock-space setting.

We then combine this existence result with the rigidity of local maximizers established earlier. Since any global maximizer is, in particular, a local maximizer, its eigenfunction is rigid and therefore determines the maximizing domain. This provides a new conceptual route to the Nicola--Tilli theorem, avoiding rearrangement inequalities and instead relying on inverse and free-boundary methods. The argument highlights the power of treating eigenfunctions as geometric data.

\subsection{Organization of the paper}

The paper is organized as follows. In Section~\ref{sec:preliminaries} we review the Bargmann--Fock representation and establish several auxiliary results, including the bridge from time--frequency analysis to free boundary problems via the key divisibility lemma. Section~\ref{sec:perturb-proof} proves Theorem~\ref{thm:perturb}; the self-contained disk inverse theorem is established there as part of the construction. Section~\ref{sec:classical-rigidity} explains how the same machinery recovers the classical theorem of Abreu and D\"orfler in the simply connected case (Theorem~\ref{thm:eig-loc-U}). Section~\ref{sec:GGRT-sharp} contains the sharpness result, Theorem~\ref{thm:GGRT-sharp}. Section~\ref{sec:local-max} establishes the local rigidity statement, showing that local maximizers of the Faber--Krahn problem in the Gaussian STFT setting are disks. Finally, Section~\ref{sec:global-max} proves existence of global maximizers via concentration--compactness in the Bargmann--Fock space and combines this with the local rigidity to give a new proof of the Nicola--Tilli theorem.

\medskip

In summary, the results of this paper establish a bridge between time--frequency analysis, inverse spectral theory, and free-boundary problems in potential theory. This perspective not only resolves several open problems, but also supplies a flexible and robust framework for studying a wide range of questions in phase--space analysis.

% ==============================================================
% Section 2: Preliminaries
% ==============================================================

\section{Preliminaries}\label{sec:preliminaries}

\subsection{Properties of the Bargmann-Fock space} We briefly recall the connection between the Gaussian short-time Fourier transform
(STFT) and the Bargmann--Fock space, which will be used systematically in what
follows. Excellent references for the material in this subsection are
\cite{Bargmann,Folland,Grochenig,Zhu}; see also
\cite{Berezin,Berger-Coburn,Daubechies}.

\medskip

Let $\varphi(x) = 2^{1/4} e^{-\pi x^2}$ be the normalized Gaussian. For $f \in L^2(\mathbb{R})$, its Gaussian STFT is
\begin{equation}\label{eq:STFT}
V_\varphi f(x,\omega)
= \int_{\mathbb{R}} f(t)\, \varphi(x-t)\, e^{-2\pi i t \omega}\, dt.
\end{equation}
The Bargmann transform
$\mathcal{B} : L^2(\mathbb{R}) \to \mathcal{F}^2(\mathbb{C})$ is defined by
\begin{equation}\label{eq:Bargmann}
(\mathcal{B}f)(z)
= 2^{1/4} \int_{\mathbb{R}} f(t)\, e^{2\pi t z - \pi t^2 - \frac{\pi}{2} z^2}\, dt,
\qquad z \in \mathbb{C}.
\end{equation}
It is well known that $\mathcal{B}$ is a unitary isomorphism onto the Bargmann--Fock
space $\mathcal{F}^2(\mathbb{C})$, endowed with the norm
\[
\|F\|_{\mathcal{F}^2}^2
= \int_{\mathbb{C}} |F(z)|^2 e^{-\pi |z|^2} \, dA(z).
\]
A direct computation shows that, for $z = x + i\omega$,
\begin{equation}\label{eq:STFT-B}
V_\varphi f(x,-\omega)
= e^{\pi i x \omega}\, (\mathcal{B}f)(z)\, e^{-\frac{\pi}{2}|z|^2}.
\end{equation}
In particular, setting $F = \mathcal{B}f$, one obtains
\begin{equation}\label{eq:energy}
\int_{\mathbb{R}^2} |V_\varphi f(x,\omega)|^2 \, dx\, d\omega
= \int_{\mathbb{C}} |F(z)|^2 e^{-\pi |z|^2} \, dA(z)
= \|F\|_{\mathcal{F}^2}^2.
\end{equation}

\medskip

\noindent
\subsubsection*{Weyl operators and invariance.}
For $a \in \mathbb{C}$, define the operator
\begin{equation}\label{eq:Ta}
(T_a F)(z) := e^{\pi \overline{a} z - \frac{\pi}{2}|a|^2} F(z - a).
\end{equation}
Then $T_a$ is unitary on $\mathcal{F}^2(\mathbb{C})$. Indeed,
\begin{align*}
\|T_aF\|_{\mathcal F^2}^2
&=
\int_{\C} e^{2\pi \Re(\overline{a}z)-\pi |a|^2}\,|F(z-a)|^2 e^{-\pi |z|^2}\,dA(z) \\
&=
\int_{\C} |F(z-a)|^2 e^{-\pi |z-a|^2}\,dA(z)
=
\|F\|_{\mathcal F^2}^2,
\end{align*}
because
\[
2\Re(\overline{a}z)-|a|^2-|z|^2=-|z-a|^2.
\]
Moreover,
\begin{equation}\label{eq:uF-translate}
u_{T_a F}(z) = u_F(z - a).
\end{equation}
Indeed, since
\(2\Re(\overline{a}z)-|a|^2-|z|^2 = -|z-a|^2\),
\begin{equation}\label{eq:u-translation-global}
u_{T_aF}(z)
=
|T_aF(z)|^2e^{-\pi |z|^2}
=
e^{2\pi\Re(\overline{a}z)-\pi|a|^2}|F(z-a)|^2e^{-\pi |z|^2}
=
|F(z-a)|^2e^{-\pi |z-a|^2}=u_F(z-a).
\end{equation}
In particular, the functionals $I_F$ defined by,
for \(s>0\),
\[
I_F(s):=\sup_{|\Omega|=s}\int_\Omega u_F\,dA,
\]
satisfy, by the translation-invariance of the Lebesgue measure on \(\C\),
\begin{equation}\label{eq:I-translation-global}
I_{T_a F}(s) = I_F(s), \qquad a \in \mathbb{C},
\end{equation}
reflecting the invariance of the Gaussian STFT under time--frequency shifts;
see \cite{Daubechies,Grochenig,Kulikov-Oliveira-Ramos,Goncalves-Oliveira-Steinerberger}.

\medskip

\noindent
\subsubsection*{Connection with localization.}
If $F = \mathcal{B}f$, then for every measurable $\Omega \subset \mathbb{C}$,
\[
\int_\Omega u_F \, dA
= \int_{\widetilde{\Omega}} |V_\varphi f(x,\omega)|^2 \, dx\, d\omega,
\]
where $\widetilde{\Omega} = \{(x,\omega) : (x,-\omega) \in \Omega\}$.
Thus, concentration properties of the STFT are naturally encoded by the function
\begin{equation}\label{eq:uF}
u_F(z) := |F(z)|^2 e^{-\pi |z|^2}, \qquad z \in \mathbb{C},
\end{equation}
which is the main object of interest in the Bargmann--Fock setting, in terms of the
quantities $I_F(s)$.

\medskip

We next collect a few elementary facts that will be used repeatedly.

\begin{lemma}\label{lem:Fock-basic}
Let $F,G \in \mathcal{F}^2(\mathbb{C})$. Then:
\begin{enumerate}
\item[(i)] The space $\mathcal{F}^2(\mathbb{C})$ is a reproducing kernel Hilbert space with kernel
\[
K_w(z) = e^{\pi z \overline{w}}, \qquad z,w \in \mathbb{C}.
\]

\item[(ii)] (Sharp pointwise bound) For every $z \in \mathbb{C}$,
\[
u_F(z) \le \|F\|_{\mathcal{F}^2}^2,
\]
and equality holds at some point $w \in \mathbb{C}$ if and only if $F$ is
proportional to the reproducing kernel $K_w$.

\item[(iii)] $u_F(z) \to 0$ as $|z| \to \infty$.

\item[(iv)] One has the estimate
\[
\|u_F - u_G\|_{L^1(\mathbb{C})}
\le \big( \|F\|_{\mathcal{F}^2} + \|G\|_{\mathcal{F}^2} \big)
\|F - G\|_{\mathcal{F}^2}.
\]
\end{enumerate}
\end{lemma}

\begin{proof}
(i) For \(w\in \C\), the function \(K_w(z)=e^{\pi z\overline w}\) is entire and
satisfies, after completing the square via
\(2\Re(z\overline w)-|z|^2=-|z-w|^2+|w|^2\),
\[
\|K_w\|_{\mathcal F^2}^2
=
\int_{\C} e^{2\pi \Re(z\overline w)-\pi |z|^2}\,dA(z)
=
e^{\pi |w|^2}\int_{\C}e^{-\pi |z-w|^2}\,dA(z)
=
e^{\pi |w|^2},
\]
since the Gaussian \(e^{-\pi |\cdot|^2}\) integrates to \(1\) over \(\C\) (with
the Lebesgue measure \(dA\) on \(\C\simeq \mathbb R^2\)). In particular
\(K_w\in \mathcal F^2(\C)\). Now, expanding \(F\in \mathcal F^2(\C)\) along the
orthonormal basis \(\{e_k\}_{k\geq 0}\) of Lemma \ref{lem:monomials} and using
the Taylor expansion
\(K_w(z)=\sum_{k\geq 0}\frac{\pi^k\overline w^k}{k!}z^k=\sum_{k\geq 0}\overline{e_k(w)}\,e_k(z)\),
one obtains the reproducing identity
\[
\langle F,K_w\rangle_{\mathcal F^2}
=
\sum_{k\geq 0} a_k\,e_k(w)
=
F(w),
\qquad F(z)=\sum_{k\geq 0}a_k e_k(z).
\]
Point evaluation is therefore continuous, and \(F\) is entire as a uniform
limit on compact sets of polynomials. This is the standard description of the
Bargmann--Fock space, see \cite{Bargmann,Folland,Grochenig,Zhu}.

(ii) Applying Cauchy--Schwarz to the reproducing identity from (i),
\begin{equation}\label{eq:pointwise}
|F(z)|
=
|\langle F,K_z\rangle_{\mathcal F^2}|
\le
\|F\|_{\mathcal F^2}\,\|K_z\|_{\mathcal F^2}
=
\|F\|_{\mathcal F^2}\,e^{\frac{\pi}{2}|z|^2},
\end{equation}
which immediately gives \(u_F(z)=|F(z)|^2 e^{-\pi |z|^2}\le \|F\|_{\mathcal F^2}^2\).
Equality at a point \(z=w\) corresponds to equality in Cauchy--Schwarz, that is,
\(F=c\,K_w\) for some \(c\in \C\); conversely, for \(F=c K_w\) one computes
\(u_F(w)=|c|^2 e^{2\pi |w|^2}e^{-\pi |w|^2}=|c|^2 e^{\pi |w|^2}=|c|^2\|K_w\|_{\mathcal F^2}^2
=\|F\|_{\mathcal F^2}^2\), proving the rigidity statement. The sharpness of this
pointwise bound goes back to \cite{Lieb}.

(iii) Fix \(\varepsilon>0\), and choose a polynomial \(P\) such that
\[
\|F-P\|_{\mathcal F^2}<\varepsilon
\]
using the density of polynomials in \(\mathcal F^2(\C)\). By the pointwise bound from (ii),
\[
|F(z)-P(z)|e^{-\frac{\pi}{2}|z|^2}\le \varepsilon
\qquad\text{for every }z\in \C.
\]
Since \(P\) is a polynomial,
\[
P(z)e^{-\frac{\pi}{2}|z|^2}\to 0
\qquad\text{as }|z|\to\infty.
\]
Hence
\[
\limsup_{|z|\to\infty}|F(z)|e^{-\frac{\pi}{2}|z|^2}\le \varepsilon.
\]
As \(\varepsilon>0\) is arbitrary, we conclude that
\[
|F(z)|e^{-\frac{\pi}{2}|z|^2}\to 0,
\]
that is, \(u_F(z)\to 0\) as \(|z|\to\infty\).

(iv) Writing
\[
|u_F - u_G|
= \big||F|^2 - |G|^2\big| e^{-\pi|z|^2}
\le |F-G|(|F|+|G|) e^{-\pi|z|^2},
\]
the claim follows by Cauchy--Schwarz.
\end{proof}

\medskip

We now note the following crucial fact about the Bargmann transform: it
transforms the basis of normalized Hermite functions \(\{h_k\}_{k\geq 0}\) on
\(L^2(\R)\) (see \cite{Thangavelu}) onto normalized monomials, yielding a
natural orthonormal basis on $\mathcal{F}^2(\C).$ More precisely, one has
\(\mathcal B h_k=e_k\) for every \(k\ge 0\), see \cite{Bargmann,Folland,Grochenig}.

\begin{lemma}\label{lem:monomials}
The functions
\[
e_k(z) := b_k z^k, \qquad k \ge 0,
\]
form an orthonormal basis of $\mathcal{F}^2(\mathbb{C})$, where
\[
b_k = \left( \frac{\pi^k}{k!} \right)^{1/2}.
\]
In particular,
\[
\langle e_k, e_m \rangle_{\mathcal{F}^2} = \delta_{k,m}.
\]
\end{lemma}

\begin{proof}
Writing \(z=re^{i\theta}\) and using \(dA(z)=r\,dr\,d\theta\) on \(\C\),
a direct computation in polar coordinates gives
\[
\int_{\mathbb{C}} |z|^{2k} e^{-\pi |z|^2}\,dA(z)
= \int_0^{2\pi}\!\!\int_0^\infty r^{2k+1} e^{-\pi r^2}\,dr\,d\theta
= 2\pi \int_0^\infty r^{2k+1} e^{-\pi r^2}\,dr
= \frac{k!}{\pi^k},
\]
where the last identity follows from the substitution \(s=\pi r^2\), giving
\(2\pi\int_0^\infty r^{2k+1}e^{-\pi r^2}\,dr=\pi^{-k}\int_0^\infty s^k e^{-s}\,ds
=\pi^{-k}\Gamma(k+1)\). Thus
\(\|z^k\|_{\mathcal{F}^2}^2 = \frac{k!}{\pi^k}\), which yields the normalization.
For orthogonality, again writing \(z=re^{i\theta}\),
\[
\langle z^k,z^m\rangle_{\mathcal F^2}
=
\int_0^\infty r^{k+m+1}e^{-\pi r^2}\,dr\int_0^{2\pi}e^{i(k-m)\theta}\,d\theta
=
0
\qquad\text{for }k\neq m.
\]
For completeness, since \(\mathcal B\colon L^2(\R)\to \mathcal F^2(\C)\) is a
unitary isomorphism and \(\mathcal Bh_k=e_k\) for every \(k\ge 0\), where
\(\{h_k\}_{k\ge 0}\) is the orthonormal basis of Hermite functions on
\(L^2(\R)\) (see \cite{Thangavelu}), the family \(\{e_k\}_{k\ge 0}\) is the
image of an orthonormal basis under a unitary, hence itself an orthonormal
basis of \(\mathcal F^2(\C)\). Equivalently, polynomials are dense in
\(\mathcal F^2(\C)\); see \cite{Bargmann,Zhu}.
\end{proof}

\medskip

\begin{lemma}\label{lem:growth}
Let $F \in \mathcal{F}^2(\mathbb{C})$. Then:
\begin{enumerate}
\item[(i)] $F$ admits the expansion
\[
F(z) = \sum_{k=0}^\infty a_k e_k(z),
\qquad \sum_{k=0}^\infty |a_k|^2 = \|F\|_{\mathcal{F}^2}^2.
\]

\item[(ii)] (Growth bound) For every $z \in \mathbb{C}$,
\[
|F(z)| \le \|F\|_{\mathcal{F}^2} e^{\frac{\pi}{2}|z|^2}.
\]

\item[(iii)] (Derivative bounds) For every $z \in \mathbb{C}$,
\[
|F'(z)| \le C(1+|z|)\, \|F\|_{\mathcal{F}^2} e^{\frac{\pi}{2}|z|^2}.
\]
\end{enumerate}
\end{lemma}

\begin{proof}
(i) follows from the orthonormal basis in Lemma \ref{lem:monomials}.

(ii) follows from the reproducing kernel estimate:
\[
|F(z)| = |\langle F, K_z \rangle|
\le \|F\|_{\mathcal{F}^2} \|K_z\|_{\mathcal{F}^2}
= \|F\|_{\mathcal{F}^2} e^{\frac{\pi}{2}|z|^2}.
\]

(iii) Starting from the reproducing identity
\(F(z)=\langle F,K_z\rangle_{\mathcal F^2}=\int_{\C}F(w)e^{\pi \overline{w}z}
e^{-\pi |w|^2}\,dA(w)\) and differentiating in \(z\) under the integral sign
(justified by the absolute convergence on compacta of \(z\), via the pointwise
bound from (ii)), we obtain
\[
F'(z)=\pi \int_{\C}F(w)\,\overline{w}\,e^{\pi\overline{w}z}e^{-\pi |w|^2}\,dA(w)
=\pi \langle F, w\,K_z\rangle_{\mathcal F^2},
\]
since \(\overline{wK_z(w)}=\overline{w}\,e^{\pi\overline{w}z}\). Cauchy--Schwarz
then gives
\[
|F'(z)|
\le
\pi \|F\|_{\mathcal F^2}\,\|wK_z\|_{\mathcal F^2}.
\]
Now, by completing the square as in (i),
\[
\|wK_z\|_{\mathcal F^2}^2
=
\int_{\C} |w|^2 e^{2\pi\Re(w\overline z)-\pi |w|^2}\,dA(w)
=
e^{\pi |z|^2}\int_{\C} |w+z|^2 e^{-\pi |w|^2}\,dA(w),
\]
after the change of variables \(w\mapsto w+z\). Expanding
\(|w+z|^2=|w|^2+2\Re(w\overline z)+|z|^2\), and using that
\(\int_\C e^{-\pi |w|^2}\,dA(w)=1\),
\(\int_\C |w|^2 e^{-\pi |w|^2}\,dA(w)=\frac{1}{\pi}\), and
\(\int_\C w\,e^{-\pi |w|^2}\,dA(w)=0\) (by rotational invariance), we conclude
\[
\int_{\C} |w+z|^2 e^{-\pi |w|^2}\,dA(w)
=
\frac{1}{\pi}+|z|^2
\le
C(1+|z|^2),
\]
for some absolute constant \(C>0\). Therefore
\[
|F'(z)| \le C(1+|z|)\,\|F\|_{\mathcal F^2}e^{\frac{\pi}{2}|z|^2}.
\]
\end{proof}

\medskip

\begin{lemma}\label{lem:max}
Let $F \in \mathcal{F}^2(\mathbb{C})$ and set
\[
T := \sup_{z \in \mathbb{C}} u_F(z).
\]
Then $0 \le T \le \|F\|_{\mathcal{F}^2}^2$, and $T = \|F\|_{\mathcal{F}^2}^2$
is attained if and only if $F$ is proportional to the reproducing kernel $K_w$
for some $w \in \mathbb{C}$.
\end{lemma}

\begin{proof}
The bound and the characterization of the equality case are immediate from
Lemma~\ref{lem:Fock-basic}(ii).
    \end{proof}

\medskip

    We proceed with a few elementary facts.

\begin{lemma}\label{lem:basic-properties-global}
Let \(F,G\in \mathcal F^2(\C)\). Then:
\begin{enumerate}
    \item[(i)] \(I_F(s)\) is attained by a superlevel set of \(u_F\), and
    \[
    0\le I_F(s)\le \|F\|_{\mathcal F^2}^2.
    \]
    \item[(ii)] For every \(c\in \C\),
    \[
    I_{cF}(s)=|c|^2 I_F(s).
    \]
    \item[(iii)] One has the Lipschitz bound
    \[
    |I_F(s)-I_G(s)|
    \le
    \|u_F-u_G\|_{L^1(\C)}
    \le
    \big(\|F\|_{\mathcal F^2}+\|G\|_{\mathcal F^2}\big)\,
    \|F-G\|_{\mathcal F^2}.
    \]
\end{enumerate}
\end{lemma}

\begin{proof}
The identity
\[
\int_\C u_F\,dA=\int_\C |F(z)|^2e^{-\pi |z|^2}\,dA(z)=\|F\|_{\mathcal F^2}^2
\]
is immediate from the definition of the Fock norm. By
Lemma~\ref{lem:Fock-basic}(ii)--(iii), \(u_F\) is bounded, continuous, and
satisfies \(u_F(z)\to 0\) as \(|z|\to\infty\); thus
\(u_F\in C_0(\C)\cap L^1(\C)\).

The existence of a maximizing set for \(I_F(s)\), of the form
\(\{u_F>t\}\) (possibly together with part of \(\{u_F=t\}\)) with \(t\geq 0\)
chosen so that the resulting set has measure \(s\), is the classical bathtub
principle, which applies since \(u_F\in C_0(\C)\cap L^1(\C)\). The bound
\[
0\le I_F(s)\le \int_\C u_F\,dA=\|F\|_{\mathcal F^2}^2
\]
is immediate. The scaling identity in (ii) is obvious.

For (iii), for every measurable \(\Omega\subset \C\) with \(|\Omega|=s\),
\[
\left|\int_\Omega u_F\,dA-\int_\Omega u_G\,dA\right|
\le \int_\C |u_F-u_G|\,dA.
\]
Taking suprema over \(|\Omega|=s\) gives
\[
|I_F(s)-I_G(s)|\le \|u_F-u_G\|_{L^1(\C)}.
\]
Finally,
\begin{align*}
\|u_F-u_G\|_{L^1(\C)}
&=
\int_\C \big||F|^2-|G|^2\big|\,e^{-\pi |z|^2}\,dA \\
&\le
\int_\C |F-G|(|F|+|G|)e^{-\pi |z|^2}\,dA \\
&\le
\|F-G\|_{\mathcal F^2}\big(\|F\|_{\mathcal F^2}+\|G\|_{\mathcal F^2}\big)
\end{align*}
by Cauchy--Schwarz.
\end{proof}


\subsection{From time-frequency analysis to free boundary problems}

Having translated, in the previous subsection, the time--frequency
localization problem into a question about the vanishing of a Fourier
transform along the two complex lines $\eta=\pm i\xi$, we now bridge the gap
between this Fourier-analytic statement and the free boundary formulation
that will dominate the remainder of the paper. The bridge is provided by a
divisibility lemma of Ehrenpreis--Malgrange flavour: any compactly supported
distribution whose Fourier transform vanishes on the characteristic variety
$\{\eta=\pm i\xi\}$ of the Laplacian must be the Laplacian of a compactly
supported $L^2$ function which, in our setting, is moreover continuous off a
finite set of points.

In order to state the lemma cleanly, we make precise the class of
distributions with which we work. We say that a compactly supported
distribution $\omega$ is a \emph{singular-continuous free measure} if it
defines a (finite, signed) Radon measure on $\R^2$ whose Lebesgue
decomposition with respect to planar Lebesgue measure $\mathrm{d}x\,
\mathrm{d}y$ takes the form
\[
\omega = \omega_{\mathrm{ac}} + \omega_{\mathrm{pp}}
       = f\,\mathrm{d}x\,\mathrm{d}y + \sum_{j=1}^N c_j\,\delta_{p_j},
\]
with \(f\in L^\infty(\R^2)\) compactly supported, \(c_j\in \R\setminus\{0\}\), and
$p_1,\ldots,p_N\in \R^2$ pairwise distinct. In other words, the continuous
singular part $\omega_{\mathrm{cs}}$ in the Lebesgue--Radon--Nikodym
decomposition $\omega=\omega_{\mathrm{ac}}+\omega_{\mathrm{cs}}+
\omega_{\mathrm{pp}}$ is required to vanish identically, and the pure
point part is finite. The terminology \emph{free measure} is suggested by
the role of $\omega$ as the right-hand side of a free boundary problem
arising in the next subsection.

\begin{lemma}\label{lemma:fourier-to-fboundary}
Let $\omega\in \mathcal{S}'(\R^2)$ be a compactly supported singular-continuous
free distribution such that
\[
\widehat{\omega}(w,\pm iw)=0,\qquad \forall w\in\C.
\]
Then there exist a finite set $P\subset \R^2$ and a compactly supported
function $u\in L^2(\R^2)$ such that
\[
\Delta u=\omega
\]
in the sense of distributions, and
\[
u\in C(\R^2\setminus P).
\]
\end{lemma}

\begin{proof}
By the Paley--Wiener--Schwartz theorem (see, e.g.,
\cite{Hormander,Folland}), the Fourier
% TODO(bib): the keys \texttt{Hormander} (Ark. Mat. 1991) and
% \texttt{Folland} (Harmonic Analysis in Phase Space) do not contain
% Paley--Wiener--Schwartz at the cited locations; a textbook entry
% (Hormander Vol.~I or Rudin's Functional Analysis) should be added.
transform $\widehat\omega(\xi,\eta)$ extends to an entire function on $\C^2$
of exponential type, satisfying the bound
\[
|\widehat\omega(\xi,\eta)|\le C(1+|\xi|+|\eta|)^{N}\,e^{R(|\Im\xi|+|\Im\eta|)}
\]
for some $C,R>0$ and some integer $N\ge 0$ depending on the order of
$\omega$, where $R$ is the radius of any disk containing
$\operatorname{supp}\omega$. In particular, $\widehat\omega$ admits an
everywhere convergent power series expansion
\[
\widehat{\omega}(\xi,\eta)=\sum_{k,l\ge 0} a_{k,l}\,\xi^k\eta^l,
\qquad (\xi,\eta)\in\C^2.
\]
Using the vanishing hypothesis on the complex lines $\eta=\pm i\xi$, we may
group terms by total homogeneity: setting $\eta=\pm iw$, $\xi=w$, and
collecting powers of $w$,
\[
0=\widehat{\omega}(w,\pm iw)
   =\sum_{n\ge 0}\Bigl(\sum_{k=0}^n a_{k,n-k}(\pm i)^{\,n-k}\Bigr)w^n,
\qquad \forall w\in\C.
\]
Since this holds for all $w\in\C$, comparing coefficients yields
\[
\sum_{k=0}^n a_{k,n-k}(\pm i)^{\,n-k}=0,\qquad \forall n\ge 0.
\]
For each $n\ge 0$, define the auxiliary one-variable polynomial
\[
p_n(t):=\sum_{k=0}^n a_{k,n-k}\, t^{\,n-k}.
\]
The two scalar equations above precisely say that $p_n(\pm i)=0$, so
$(t-i)(t+i)=t^2+1$ divides $p_n(t)$. Hence we may write
\[
p_n(t)=(t^2+1)\,q_n(t),
\]
with $q_n$ a polynomial of degree at most $n-2$ (and $q_n\equiv 0$ when
$n<2$).

We now reassemble the homogeneous pieces. Let
\[
H_n(\xi,\eta):=\sum_{k=0}^n a_{k,n-k}\,\xi^k\eta^{\,n-k},
\qquad n\ge 0,
\]
so that $\widehat\omega=\sum_{n\ge 0}H_n$ is the decomposition of
$\widehat\omega$ into homogeneous components (of degree $n$). The argument
above shows that each $H_n$ vanishes on the two complex lines
$\eta=\pm i\xi$, hence is divisible, as a homogeneous polynomial in two
variables, by $\xi^2+\eta^2$. Concretely,
\[
H_n(\xi,\eta)=(\xi^2+\eta^2)\,G_{n-2}(\xi,\eta)
\]
for a (uniquely determined) homogeneous polynomial $G_{n-2}$ of degree
$n-2$, with $G_{n-2}\equiv 0$ for $n<2$. Summing the homogeneous pieces, the
formal series
\[
F(\xi,\eta):=\sum_{n\ge 2} G_{n-2}(\xi,\eta)
\]
converges absolutely on every polydisc on which $\widehat\omega$ does, and
defines an entire function on $\C^2$ satisfying
\[
\widehat{\omega}(\xi,\eta)=(\xi^2+\eta^2)\,F(\xi,\eta).
\]
Since $\widehat\omega$ has a zero of order at least two at the origin
(coming from the vanishing of $H_0$ and $H_1$, which is forced by
$p_0(\pm i)=p_1(\pm i)=0$), this factorization is consistent at $0$ and $F$
is entire there as well.

We next quantify the size of $F$ on $\R^2$. Since $\omega$ is a finite
(signed) Radon measure with compact support, its Fourier transform is in
fact bounded on $\R^2$:
\[
|\widehat{\omega}(\xi,\eta)|\le \|\omega\|_{\mathrm{TV}},
\qquad (\xi,\eta)\in\R^2,
\]
where $\|\omega\|_{\mathrm{TV}}$ is the total variation; equivalently,
the Paley--Wiener bound holds with $N=0$. On $\R^2$, $\xi^2+\eta^2=
|\xi|^2+|\eta|^2$, so for $|\xi|^2+|\eta|^2\ge 1$ we have $\xi^2+\eta^2
\ge \tfrac{1}{2}(1+|\xi|^2+|\eta|^2)$, whence
\[
|F(\xi,\eta)|=\frac{|\widehat\omega(\xi,\eta)|}{\xi^2+\eta^2}
\le \frac{2\|\omega\|_{\mathrm{TV}}}{1+|\xi|^2+|\eta|^2},
\qquad |\xi|^2+|\eta|^2\ge 1.
\]
Near the origin, $F$ is entire (as established above) and hence bounded;
absorbing this region into the constants we obtain
\[
|F(\xi,\eta)|\le \frac{C''}{1+|\xi|^2+|\eta|^2},
\qquad (\xi,\eta)\in\R^2,
\]
for some $C''>0$. Therefore $F|_{\R^2}\in L^2(\R^2)\cap L^1(\R^2)$. On
$\C^2$, $F$ is entire (by the formal-series construction above) and
satisfies an exponential-type bound of the same rate $R$ as
$\widehat\omega$. Indeed, on the open set
$\Omega:=\{(\xi,\eta)\in\C^2:|\xi^2+\eta^2|\ge 1\}$ one has the pointwise
estimate
\[
|F(\xi,\eta)|=\frac{|\widehat\omega(\xi,\eta)|}{|\xi^2+\eta^2|}
\le C(1+|\xi|+|\eta|)^{N}\,e^{R(|\Im\xi|+|\Im\eta|)},
\qquad (\xi,\eta)\in\Omega,
\]
by the Paley--Wiener bound on $\widehat\omega$. On the complement
$\Omega^c=\{|\xi^2+\eta^2|<1\}$, the function $F$ is entire and the
maximum modulus principle, applied on each complex line through a given
point in $\Omega^c$ in a direction transverse to the variety
$\{\xi^2+\eta^2=0\}$, bounds $|F|$ at points of $\Omega^c$ by a uniform
constant times the supremum of $|F|$ on the boundary of a fixed thin
tubular neighbourhood of that variety, where the previous estimate
applies. Combining these two regimes and absorbing the resulting
bookkeeping constant yields, after enlarging $C'$ if necessary,
\[
|F(\xi,\eta)|\le C'\,(1+|\xi|+|\eta|)^{N}\,
e^{R(|\Im \xi|+|\Im \eta|)},
\qquad (\xi,\eta)\in\C^2,
\]
which is the Paley--Wiener--Schwartz bound for a compactly supported
distribution of order $\le N$ in $\overline{B(0,R)}$.

By the Paley--Wiener--Schwartz theorem applied in the reverse direction
(see again \cite{Hormander}), the entire function $F$
satisfies an exponential-type bound on $\C^2$ with the same exponential
rate $R$ as $\widehat\omega$, and integrable decay along $\R^2$, so it is
the Fourier--Laplace transform of a distribution of order $\le N$
compactly supported in the closed disk $\overline{B(0,R)}$. The $L^2$
bound on the real locus shows that this distribution is, in fact, an
$L^2(\R^2)$ function. Let $u\in L^2(\R^2)$ denote this compactly supported
representative, characterised by $\widehat u = F$ on $\R^2$ and
$\operatorname{supp} u\subset \overline{B(0,R)}$.

Computing distributional Laplacians on the Fourier side, with the
convention $\widehat g(\xi,\eta)=\int g(x,y)e^{-2\pi i(x\xi+y\eta)}\,
\mathrm{d}x\,\mathrm{d}y$, we have
\[
\widehat{\Delta u}(\xi,\eta) = -4\pi^2(\xi^2+\eta^2)\,\widehat{u}(\xi,\eta)
= -4\pi^2(\xi^2+\eta^2)\,F(\xi,\eta) = -4\pi^2\,\widehat\omega(\xi,\eta);
\]
with any other standard normalisation, the same identity holds up to a
nonzero real multiplicative factor. In all cases, $\Delta u$ is a constant
nonzero real multiple of $\omega$. Replacing $u$ by $\lambda u$ for the
appropriate (nonzero, real) $\lambda$ -- this is the harmless rescaling
referred to above -- we may, and do, assume that
\[
\Delta u = \omega
\]
in the sense of distributions. The rescaling preserves all the previously
established properties of $u$: compact support, membership in $L^2(\R^2)$,
and the Fourier identity $\widehat u = \lambda F$.

To pin down $u$ pointwise, observe that since $\omega$ is a compactly
supported Radon measure on $\R^2$, the convolution
\[
v(x,y) := \frac{1}{2\pi}(\log|\cdot|*\omega)(x,y)
\]
is a well-defined locally integrable function on $\R^2$ which solves
$\Delta v = \omega$ in the sense of distributions; here $G(z)=\log|z|$ is
(up to the standard normalisation $\tfrac{1}{2\pi}$) the fundamental solution
of the Laplacian in $\R^2$, in the sense that
$\Delta\bigl(\tfrac{1}{2\pi}\log|\cdot|\bigr)=\delta_0$, see, e.g.,
\cite{Hormander} for the underlying distribution theory and
\cite{Aharonov-Schiffer-Zalcman} for related representations in the planar
quadrature-domain context.
% TODO(bib): \texttt{Garnett} (Bounded Analytic Functions, GTM 236) does
% not contain the previously-cited Ch.~II Thm~1.3 on continuity of
% logarithmic potentials of bounded densities; consider adding Ransford
% (Potential Theory in the Complex Plane) or Landkof (Foundations of
% Modern Potential Theory) to the bibliography.
The difference $u-v$ is then a tempered
distribution which is harmonic on $\R^2$ (since
$\Delta(u-v)=\omega-\omega=0$), and hence, being harmonic, real-analytic.
Outside the union of the supports of $u$ and $\omega$, the function $v$ is
in fact real-analytic and grows at most logarithmically at infinity, with
$v(x)=\frac{M}{2\pi}\log|x|+O(|x|^{-1})$ as $|x|\to\infty$ where
$M=\omega(\R^2)$, while $u$ vanishes there because of its compact support.
Hence $u-v$ is harmonic on $\R^2$ and grows at most logarithmically at
infinity; a standard refinement of the harmonic Liouville theorem (any
harmonic function on $\R^2$ growing slower than a polynomial of degree $1$
is constant) then forces $u-v$ to be a constant, and this constant must
vanish to be consistent with $u\in L^2(\R^2)$ and the local integrability
of $v$. Equivalently, and more transparently: by the compact support of
$\omega$ together with $\widehat u = F$ on $\R^2$, the function $u$ is
uniquely determined among compactly supported $L^2$ solutions of
$\Delta u = \omega$, since any two such would differ by an entire $L^2$
function which is harmonic, hence vanishes. Thus
\[
u(x,y)=\frac{1}{2\pi}\int_{\R^2}\log|(x,y)-(x',y')|\,\mathrm{d}\omega(x',y'),
\]
the logarithmic potential of $\omega$.

We finally analyse the regularity of $u$. Decomposing
$\omega = f\,\mathrm{d}x\,\mathrm{d}y+\sum_{j=1}^N c_j\,\delta_{p_j}$ as in
the definition of singular-continuous free measure, we obtain
correspondingly
\[
u(x,y)=\frac{1}{2\pi}\int_{\R^2}\log|(x,y)-(x',y')|\,f(x',y')\,
\mathrm{d}x'\mathrm{d}y'
+\frac{1}{2\pi}\sum_{j=1}^N c_j\,\log|(x,y)-p_j|.
\]
The first summand is the logarithmic potential of the bounded, compactly
supported density \(f\in L^\infty(\R^2)\) and is therefore not merely
continuous but in fact $C^{1,\alpha}(\R^2)$ for every $\alpha\in(0,1)$
(by the standard Calder\'on--Zygmund / Schauder estimates applied to
$\Delta=\partial_x^2+\partial_y^2$ with right-hand side in $L^\infty$
of compact support, see, e.g., the Newtonian-potential discussion in
\cite{Hormander}).
The continuity used here is in fact elementary: by dominated convergence,
$z\mapsto \int \log|z-w|\,f(w)\,\mathrm{d}A(w)$ is continuous since
$\log|\cdot-w|$ is locally integrable in $w$ uniformly on compacta in $z$,
and $f$ is bounded with compact support. It is here that the absence of a continuous
singular part in $\omega$ is essential: a generic continuous singular
measure (e.g.\ the Hausdorff measure on a Cantor-type set of positive
$1$-capacity) gives rise to a logarithmic potential which need not be
continuous, only quasi-continuous, and the conclusion of the lemma would
fail in that generality. Each term $\log|(x,y)-p_j|$ is continuous on
$\R^2\setminus\{p_j\}$, with a logarithmic singularity at $p_j$. Setting
$P:=\{p_1,\ldots,p_N\}$, we conclude that $u\in C(\R^2\setminus P)$, as
claimed. The fact that $u$ is compactly supported has already been
established in the construction above and is consistent with the
representation as a logarithmic potential, which differs from a compactly
supported function only by a globally harmonic correction that, as
explained, must vanish identically.
\end{proof}

We postpone the proof of Theorem~\ref{thm:eig-loc-U} until after the
perturbative inverse theorem. The same inverse/free-boundary mechanism used
to construct domains for prescribed eigenfunctions also recovers the
classical Abreu--D\"orfler disk theorem in the simply connected setting; see
Section~\ref{sec:classical-rigidity} below.


% ==============================================================
% Section 3
% ==============================================================

\section{Proof of Theorem \ref{thm:perturb}}\label{sec:perturb-proof}

\subsection{The disk inverse step}\label{subsec:disk-inverse-step}

The proof rests on a local inverse theorem at the Daubechies disk. The result
below is self-contained, but its point of view is inspired by perturbative
inverse-source and quadrature-domain methods: Isakov's local free-boundary
scheme for inverse source problems, Cherednichenko's local inverse logarithmic
potential theorem, and the quadrature-domain framework developed by
Gustafsson, Shahgholian, Karp, and Aharonov--Schiffer--Zalcman. In the present
disk-based setting the required inverse step can be proved directly by
diagonalizing the full moment map in Fourier modes.

Fix \(0<\lambda<1\), and let \(r_\lambda>0\) be determined by
\[
1-e^{-\pi r_\lambda^2}=\lambda .
\]
For a real-valued \(2\pi\)-periodic function \(\varphi\) with
\(\|\varphi\|_{C^0}<r_\lambda/2\), set
\[
U_\varphi
:=
\{re^{i\theta}:0\le r<r_\lambda+\varphi(\theta)\}.
\]
For \(\tau>0\), let \(\mathcal A_\tau^\R\) be the Banach space of real-valued
periodic functions
\[
\varphi(\theta)=\sum_{k\in\Z}\varphi_k e^{ik\theta},
\qquad
\varphi_{-k}=\overline{\varphi_k},
\]
with norm
\[
\|\varphi\|_{\mathcal A_\tau}
:=
\sum_{k\in\Z}|\varphi_k|e^{\tau |k|}<\infty .
\]
This is a Banach algebra, and \(\mathcal A_\tau^\R\hookrightarrow C^m\) for
every fixed \(m\). In particular, small elements of \(\mathcal A_\tau^\R\)
produce real-analytic \(C^2\) Jordan domains.

\begin{proposition}[Disk symmetrized inverse theorem]\label{prop:disk-inverse-tailored}
Let \(N\ge1\), \(0<\lambda<1\), and \(r_\lambda\) be as above. There are
\(\tau>0\), \(\varepsilon>0\), and a \(C^1\) map
\[
b=(b_1,\ldots,b_N)\longmapsto \varphi_b\in\mathcal A_\tau^\R,
\qquad
\varphi_0=0,
\]
defined for \(\max_j|b_j|<\varepsilon\), such that, for
\[
P_b(z):=1+\sum_{j=1}^N b_jz^j,
\qquad
P_b^\#(z):=\overline{P_b(\overline z)},
\qquad
Q_b(w):=P_b^\#(\partial_w/\pi)P_b(w),
\]
and
\[
U_b:=U_{\varphi_b},
\]
one has
\begin{equation}\label{eq:disk-inverse-identity}
\int_{U_b}|P_b(z)|^2e^{-\pi |z|^2}e^{\pi\overline z w}\,dA(z)
=
\lambda Q_b(w),
\qquad w\in\C .
\end{equation}
Moreover,
\[
\|\varphi_b\|_{C^2(\mathbb S^1)}
\le C_{\lambda,N}\max_{1\le j\le N}|b_j|,
\]
and the branch is locally unique among sufficiently small analytic radial
graphs in \(\mathcal A_\tau^\R\).
\end{proposition}

\begin{proof}
We solve the coefficient equations obtained by expanding the entire identity
\eqref{eq:disk-inverse-identity}. Write
\[
Q_b(w)=\sum_{j=0}^N q_j(b)w^j,
\]
and \(q_j(b)=0\) for \(j>N\). The constant coefficient is real:
\[
q_0(b)=1+\sum_{k=1}^N \frac{k!}{\pi^k}|b_k|^2.
\]
Since
\[
e^{\pi\overline z w}
=
\sum_{j=0}^\infty \frac{\pi^j}{j!}\overline z^{\,j}w^j,
\]
the desired identity is equivalent to the infinite system
\begin{equation}\label{eq:disk-moment-system}
\frac{\pi^j}{j!}\int_{U_\varphi}
|P_b(z)|^2\overline z^{\,j}e^{-\pi |z|^2}\,dA(z)
=
\lambda q_j(b),
\qquad j\ge0 .
\end{equation}

We introduce the weighted sequence space adapted to the linearization at the
circle. Put
\[
a_j:=2\pi e^{-\pi r_\lambda^2}r_\lambda^{j+1}\frac{\pi^j}{j!},
\qquad j\ge0,
\]
and let \(\mathcal Y_\tau\) be the space of sequences
\(y=(y_j)_{j\ge0}\), with \(y_0\in\R\), such that
\[
\|y\|_{\mathcal Y_\tau}
:=
\frac{|y_0|}{a_0}
+2\sum_{j=1}^\infty e^{\tau j}\frac{|y_j|}{a_j}
<\infty .
\]
Define
\[
\mathcal F_j(b,\varphi)
:=
\frac{\pi^j}{j!}\int_{U_\varphi}
|P_b(z)|^2\overline z^{\,j}e^{-\pi |z|^2}\,dA(z)
-\lambda q_j(b),
\qquad j\ge0 .
\]
Thus \(\mathcal F(b,\varphi)=0\) is precisely
\eqref{eq:disk-moment-system}. At \(b=0\), \(\varphi=0\), the domain is the
disk \(D_{r_\lambda}\), and radial symmetry gives
\[
\frac{\pi^j}{j!}\int_{D_{r_\lambda}}
\overline z^{\,j}e^{-\pi |z|^2}\,dA(z)
=
\begin{cases}
\lambda, & j=0,\\
0, & j\ge1.
\end{cases}
\]
Hence \(\mathcal F(0,0)=0\).

We next compute the derivative with respect to the boundary graph. Let
\(\psi(\theta)=\sum_{k\in\Z}\psi_k e^{ik\theta}\in\mathcal A_\tau^\R\).
Differentiating the polar-coordinate formula
\[
\int_{U_\varphi}H(z)\,dA(z)
=
\int_0^{2\pi}\int_0^{r_\lambda+\varphi(\theta)}
H(re^{i\theta})r\,dr\,d\theta
\]
at \(\varphi=0\), with
\(H(z)=\overline z^{\,j}e^{-\pi|z|^2}\), gives
\[
D_\varphi\mathcal F_j(0,0)[\psi]
=
\frac{\pi^j}{j!}
\int_0^{2\pi}
\psi(\theta)r_\lambda^{j+1}e^{-ij\theta}e^{-\pi r_\lambda^2}
\,d\theta
=
a_j\psi_j .
\]
Consequently
\[
D_\varphi\mathcal F(0,0)\colon
\mathcal A_\tau^\R\to\mathcal Y_\tau,
\qquad
\psi\mapsto(a_j\psi_j)_{j\ge0},
\]
is an isomorphism. Its inverse is explicit:
if \(y\in\mathcal Y_\tau\), then
\[
\psi_0=\frac{y_0}{a_0},
\qquad
\psi_j=\frac{y_j}{a_j}\quad(j\ge1),
\qquad
\psi_{-j}=\overline{\psi_j}.
\]
The definition of the two norms gives
\(\|\psi\|_{\mathcal A_\tau}=\|y\|_{\mathcal Y_\tau}\).

It remains to justify the use of the Banach-space implicit function theorem.
For \(\|\varphi\|_{\mathcal A_\tau}\) small, the function
\(\theta\mapsto r_\lambda+\varphi(\theta)\) extends holomorphically to the
strip \(|\Im\theta|<\tau\) and stays in a fixed annulus
\(r_\lambda/2<|r|<3r_\lambda/2\). In polar coordinates,
\[
\mathcal F_j(b,\varphi)+\lambda q_j(b)
=
\frac{\pi^j}{j!}\int_0^{2\pi}e^{-ij\theta}
\int_0^{r_\lambda+\varphi(\theta)}
P_b(re^{i\theta})P_b^\#(re^{-i\theta})r^{j+1}e^{-\pi r^2}\,dr\,d\theta .
\]
Expanding the inner integral at \(r_\lambda\) gives an absolutely convergent
power series in \(\varphi\), whose coefficients are analytic functions of
\(\theta\) and polynomial functions of \(b\). Since \(\mathcal A_\tau\) is a
Banach algebra and \(P_b(re^{i\theta})P_b^\#(re^{-i\theta})\) is a finite
trigonometric polynomial of degree at most \(N\), the resulting series and
its first derivatives in \((b,\varphi)\) converge in \(\mathcal Y_\tau\) after
possibly decreasing \(\tau\) and restricting to
\(\|b\|+\|\varphi\|_{\mathcal A_\tau}<\eta\). More explicitly, Cauchy's
estimate on a slightly wider strip gives, for \(0<\tau'<\tau\),
\[
\sum_{j\ge0}e^{\tau' j}\frac1{a_j}
\left|
\frac{\pi^j}{j!}\widehat{R_j}(j)
\right|
\le
C_{\lambda,N,\tau,\tau'}\|R\|_{\mathcal A_\tau},
\]
where \(\widehat{R_j}(j)\) denotes the \(j\)-th Fourier coefficient of the
analytic coefficient appearing in the \(j\)-th moment equation. Applying this
estimate term by term to the Taylor expansion in \(\varphi\) yields
\[
\|\mathcal F(b,\varphi)-\mathcal F(b,\psi)
-D_\varphi\mathcal F(b,\psi)[\varphi-\psi]\|_{\mathcal Y_{\tau'}}
\le
C(\|b\|+\|\varphi\|_{\mathcal A_\tau}
+\|\psi\|_{\mathcal A_\tau})
\|\varphi-\psi\|_{\mathcal A_\tau}^2,
\]
and analogous estimates for differentiation in \(b\). Shrinking the strip
once and then relabelling \(\tau'\) as \(\tau\), this proves that
\[
\mathcal F:
\{(b,\varphi):\|b\|+\|\varphi\|_{\mathcal A_\tau}<\eta\}
\longrightarrow \mathcal Y_\tau
\]
is \(C^1\).

The Banach implicit function theorem now applies at \((0,0)\), because
\(D_\varphi\mathcal F(0,0)\) is an isomorphism. Thus, for \(\|b\|\) small,
there is a unique \(\varphi_b\in\mathcal A_\tau^\R\), small in
\(\mathcal A_\tau\), such that \(\mathcal F(b,\varphi_b)=0\), and
\(b\mapsto\varphi_b\) is \(C^1\). The embedding
\(\mathcal A_\tau\hookrightarrow C^2\) gives the displayed \(C^2\) bound.
Finally, \(\mathcal F(b,\varphi_b)=0\) gives all coefficient identities
\eqref{eq:disk-moment-system}; summing the power series in \(w\) gives
\eqref{eq:disk-inverse-identity} for every \(w\in\C\), since both sides are
entire functions. This proves the proposition.
\end{proof}

\begin{remark}
The proof above is the only place where analyticity of the boundary
parametrization is used. The constructed domains are therefore \(C^\infty\),
hence \(C^2\), which is the regularity needed later. The argument does not
assert a general local quadrature-domain theorem near an arbitrary reference
domain; it proves exactly the disk-based inverse statement used below.
\end{remark}

\begin{proof}[Proof of Theorem \ref{thm:perturb}]
We again use the Bargmann transform. Let
\[
P(z)=1+\sum_{k=1}^N b_k z^k,
\]
where the coefficients $b_k$ are chosen so that the Bargmann transform of
\[
f_0=h_0+\sum_{k=1}^N a_k h_k
\]
is exactly $P$. We also write
\[
P^\#(z):=\overline{P(\overline z)}
=1+\sum_{k=1}^N \overline{b_k}z^k
\]
for the coefficientwise conjugate polynomial. Then the condition that $f_0$ be an eigenfunction of $\mathcal{L}_U$ with eigenvalue $\lambda$ is equivalent to
\begin{equation}\label{eq:eigenfunction-general-new}
\int_U P(z)e^{-\pi|z|^2}e^{\pi \overline z w}\,dA(z)=\lambda P(w),
\qquad w\in \C.
\end{equation}
Set
\[
Q(w):=P^\#(\partial_w/\pi)P(w).
\]
By Proposition~\ref{prop:disk-inverse-tailored}, after decreasing
\(\varepsilon_0(\lambda,N)\) if necessary, there exists a real-analytic radial
graph \(U=U_b\), \(C^2\)-close to the Daubechies disk of radius
\(r_\lambda\), for which
\begin{equation}\label{eq:almost-reduced-new}
\int_U |P(z)|^2e^{-\pi|z|^2}e^{\pi\overline z w}\,dA(z)
=
\lambda Q(w),
\qquad w\in\C .
\end{equation}
We now recover the unsymmetrized identity
\eqref{eq:eigenfunction-general-new}. Define
\[
G(w):=\int_U P(z)e^{-\pi|z|^2}e^{\pi\overline z w}\,dA(z).
\]
The function \(G\) is entire, and differentiating under the integral is
justified by boundedness of \(U\). Since
\[
\Big(\frac{\partial_w}{\pi}\Big)^j e^{\pi\overline z w}
=\overline z^{\,j}e^{\pi\overline z w},
\]
we have
\[
P^\#(\partial_w/\pi)G(w)
=
\int_U |P(z)|^2e^{-\pi|z|^2}e^{\pi\overline z w}\,dA(z)
=
\lambda Q(w)
=
P^\#(\partial_w/\pi)(\lambda P)(w).
\]
Thus \(H:=G-\lambda P\) is an entire solution of
\[
P^\#(\partial_w/\pi)H=0.
\]

Let \(\widetilde P^\#(\zeta):=P^\#(\zeta/\pi)\), and let
\(\zeta_1,\ldots,\zeta_\ell\) be its zeros, with multiplicities
\(m_1,\ldots,m_\ell\). The elementary structure theorem for
constant-coefficient ordinary differential equations gives
\[
H(w)=\sum_{j=1}^\ell R_j(w)e^{\zeta_j w},
\]
where \(\deg R_j<m_j\). On the other hand, if
\(M:=\sup_{z\in U}|z|\), then
\[
|G(w)|\le C(1+|w|)^N e^{\pi M|w|},
\qquad w\in\C,
\]
and therefore \(H\) has exponential type at most \(\pi M\).
Since \(U\) stays in a fixed neighborhood of \(D_{r_\lambda}\), we have
\(M\le C_\lambda\). The zeros of \(P^\#\), and hence of
\(\widetilde P^\#\), escape to infinity as \(b\to0\). Shrinking
\(\varepsilon_0(\lambda,N)\) once more, we may assume
\[
|\zeta_j|>10M,\qquad 1\le j\le \ell .
\]
If some \(R_j\) were nonzero, choose one of maximal \(|\zeta_j|\) for which
\(R_j\not\equiv0\) and relabel it \(\zeta_1\). Since the \(\zeta_j\) are
\emph{distinct} (they are the distinct roots of \(\widetilde P^\#\)),
the Cauchy--Schwarz inequality
\(\Re(\zeta_j\overline{\zeta_1})\le|\zeta_j||\zeta_1|\) is strict whenever
\(\zeta_j\ne\zeta_1\) and \(|\zeta_j|=|\zeta_1|\); when \(|\zeta_j|<|\zeta_1|\)
strictness is automatic. Hence for every \(\zeta_j\ne\zeta_1\),
\[
\Re(\zeta_j\overline{\zeta_1})<|\zeta_1|^2 .
\]
Along \(w=t\overline{\zeta_1}\), \(t\to+\infty\), the term
\(R_1(w)e^{\zeta_1w}\) has modulus \(|R_1(t\overline{\zeta_1})|e^{t|\zeta_1|^2}\),
which grows like a positive power of \(t\) times \(e^{t|\zeta_1|^2}\)
because \(R_1\not\equiv0\); every other term is bounded by a polynomial
in \(t\) times \(e^{t\Re(\zeta_j\overline{\zeta_1})}\), with strictly smaller
exponential rate. Therefore
\[
|H(t\overline{\zeta_1})|\ge c e^{t|\zeta_1|^2}
\]
for all large \(t\). The exponential-type bound gives, for every
\(\delta>0\),
\[
|H(t\overline{\zeta_1})|
\le C_\delta e^{(\pi+\delta)Mt|\zeta_1|}.
\]
Letting \(t\to\infty\) and then \(\delta\downarrow0\) yields
\(|\zeta_1|\le \pi M\), contradicting \(|\zeta_1|>10M\). Hence all
\(R_j\) vanish and \(H\equiv0\), which proves
\eqref{eq:eigenfunction-general-new}.

The Bargmann transform is unitary and intertwines the localization operator
with the integral operator in \eqref{eq:eigenfunction-general-new}. Therefore
the inverse Bargmann transform \(f_0\) of \(P\) satisfies
\[
\mathcal L_U f_0=\lambda f_0 .
\]
This proves the theorem.
\end{proof}


% ==============================================================
% Section 4
% ==============================================================

\section{A Classical Rigidity Corollary}\label{sec:classical-rigidity}

The localization operator $\mathcal{L}_U$ on the Bargmann--Fock space is
diagonalized by Hermite functions whenever $U$ is a disk centered at the
origin: this is an immediate consequence of the rotational symmetry of $U$
together with the radial structure of the Hermite weight. The converse
question, raised by Abreu and D\"orfler, asks whether eigenfunction status of
even a \emph{single} Hermite function already forces the underlying domain to
be a disk. The next theorem records a positive answer in the simply connected
$C^2$ regime.

For convenience we recall the statement of Theorem~\ref{thm:eig-loc-U} from the introduction:
\emph{let $U\subset \C$ be a simply connected bounded domain with $C^2$
boundary, and let $h_k$ be a Hermite function ($k\ge 0$). If $h_k$ is an
eigenfunction of the localization operator $\mathcal{L}_U$, then $U$ is a
disk.} (The original Abreu--D\"orfler theorem~\cite{Abreu-Doerfler} dispenses
with the boundary regularity assumption; we will recover this stronger form
in the second proof below.)

We will give two proofs of Theorem~\ref{thm:eig-loc-U}. The first, presented
below, leverages the Fourier-analytic machinery developed earlier in the
paper, reducing the question to a free boundary problem for a logarithmic
potential. The second, in the next subsection, proceeds along more
classical complex-analytic lines.

\subsection{First proof}

\begin{proof}[First proof under the additional assumption that $\partial U$ is $C^2$]
The Bargmann transform of the Hermite function $h_k$ is a nonzero scalar
multiple of the monomial $z^k$. Thus, arguing as in
\cite{Abreu-Doerfler}, the assumption that $h_k$ is an eigenfunction of
the localization operator $\mathcal{L}_U$ is equivalent to the identity
\[
\int_U z^k\, e^{\pi w\overline{z}}\, e^{-\pi |z|^2}\,dA(z)
=\lambda\, w^k,
\qquad w\in \C.
\]
Both sides are entire functions of $w$. Since $U$ is bounded and the
weight $e^{-\pi|z|^2}$ is integrable, the dominated convergence
theorem applied to difference quotients shows that for each
multi-index $\alpha$ in $w$, the partial derivative
$\partial_w^\alpha$ commutes with the integral on the left-hand side;
the absolute integrability of $|z|^{|\alpha|+k} e^{\pi
|w||z|}e^{-\pi|z|^2}\mathbf{1}_U$ for fixed $w$ provides the required
local majorant. Iterating $k$ times, we obtain
\[
\frac{1}{(\pi)^k}\,\partial_w^k
\Big[\int_U z^k\, e^{\pi w\overline{z}}\, e^{-\pi |z|^2}\,dA(z)\Big]
=\int_U z^k\overline{z}^k\, e^{\pi w\overline{z}}\, e^{-\pi |z|^2}\,dA(z)
=\int_U |z|^{2k}\, e^{\pi w\overline{z}}\, e^{-\pi |z|^2}\,dA(z).
\]
On the right-hand side, $\partial_w^k(\lambda w^k)=\lambda\,k!$.
Setting $c(\lambda,k):=\lambda\,k!/\pi^k\in\R$ and using
$\mathbf{1}_U$ to extend the integration to $\C$, we obtain
\begin{equation}\label{eq:fourier-delta}
\int_{\C} \mathbf{1}_U(z)\,|z|^{2k}\,e^{-\pi |z|^2}\, e^{\pi \overline{z}\, w}\,dA(z)
=c(\lambda,k),
\end{equation}
for some real constant $c(\lambda,k)$.

Consider now the compactly supported distribution
\[
\eta:=|z|^{2k}e^{-\pi |z|^2}\mathbf{1}_U(z)-c(\lambda,k)\,\delta_0,
\]
which has compact support in any closed ball containing $U\cup\{0\}$.
Writing $z=x+iy$ and $w\in \C$, we may rewrite the exponential factor in
\eqref{eq:fourier-delta} as
\[
e^{\pi \overline{z}\, w}=e^{\pi(x-iy)w}=e^{\pi xw}\, e^{-i\pi yw}.
\]
Hence \eqref{eq:fourier-delta} implies that the Fourier transform
$\widehat{\eta}$, viewed as an entire function on $\C^2$ via Paley--Wiener,
satisfies
\[
\widehat{\eta}(w,-iw)=0,\qquad \forall w\in\C.
\]
Since $c(\lambda,k)\in \R$, taking complex conjugates in
\eqref{eq:fourier-delta} after replacing $w$ by $\overline{w}$ also yields
\[
\widehat{\eta}(w,iw)=0,\qquad \forall w\in\C.
\]

Applying Lemma~\ref{lemma:fourier-to-fboundary} to the distribution
$\eta$, we obtain a function $u\in L^2(\R^2)$ such that
\begin{equation}\label{eq:potato}
\Delta u=f(|z|)\,\mathbf{1}_U-c(\lambda,k)\,\delta_0
\end{equation}
in the sense of distributions, where
\[
f(|z|)=|z|^{2k}e^{-\pi |z|^2}.
\]
Writing $u$ explicitly as a logarithmic potential of the right-hand side
of \eqref{eq:potato}, we conclude that $U$ satisfies
\begin{equation}\label{eq:potential}
\int_U \log|z-z'|\, f(|z'|)\,dA(z')
=c(\lambda,k)\,\log|z|,
\qquad z\in \C\setminus (U\cup\{0\}).
\end{equation}

\medskip

\noindent\emph{Reduction to $0\notin U$.}\quad
Suppose first that $0\in U$, and pick $\rho>0$ small enough that
$\overline{D_\rho(0)}\subset U$. The radial measure
$f(|z|)\,\mathbf{1}_{D_\rho(0)}\,dA$ generates a logarithmic potential
\[
\Phi_\rho(z):=\int_{D_\rho(0)} \log|z-z'|\, f(|z'|)\,dA(z'),
\]
which, by the planar Newton shell theorem for logarithmic potentials of
radial densities (see, e.g., \cite[Ch.~2]{Isakov}), is radial; outside
$D_\rho(0)$ it equals
\[
\Phi_\rho(z)=\Big(\int_{D_\rho(0)}f(|z'|)\,dA(z')\Big)\log|z|+\text{const}.
\]
Subtracting the closed disk $\overline{D_\rho(0)}$ from $U$ thus preserves
identity \eqref{eq:potential} on the complement of the new set, at the
cost of replacing $c(\lambda,k)$ by $c(\lambda,k)-M_\rho$, where
$M_\rho:=\int_{D_\rho(0)}f(|z'|)\,dA(z')>0$. The new domain
$U':=U\setminus\overline{D_\rho(0)}$ satisfies
$0\notin \overline{U'}$, so we may assume from now on that
$0\notin \overline{U}$ in \eqref{eq:potential}, with $c(\lambda,k)$
replaced (if needed) by the real constant $c(\lambda,k)-M_\rho$, which
we continue to denote by $c(\lambda,k)$. Furthermore, this constant must
be nonzero: comparing the two sides of \eqref{eq:potential} as
$|z|\to\infty$ shows that the left-hand side behaves like
$\big(\int_U f\,dA\big)\log|z|+O(|z|^{-1})$, while the right-hand side
is $c(\lambda,k)\log|z|$, forcing
$c(\lambda,k)=\int_U f(|z'|)\,dA(z')>0$.

\medskip

Define the logarithmic potential
\[
\Psi_U(z):=\int_U \log|z-z'|\, f(|z'|)\,dA(z').
\]
We claim that $\Psi_U$ is \emph{locally radial}: for each $z_0\in U$,
there exists a ball $B(z_0)\subset U$ such that $\Psi_U$ agrees on
$B(z_0)$ with a radial function of $|z|$.

To prove this, let $B\supset U$ be a ball centered at $0$, and define
\[
\Phi(z):=\int_B \log|z-z'|\, f(|z'|)\,dA(z').
\]
Since both $B$ and the weight $f(|z'|)$ are radial, $\Phi$ is radial; in
particular, $\nabla \Phi(z)$ is parallel to $z$ for every $z\in \R^2$.
Furthermore, on the interior of $B$,
\[
\Delta \Phi(z)=2\pi\, f(|z|).
\]

Now $\Psi_U$ satisfies \eqref{eq:potential} on the unbounded
component of $\C\setminus(U\cup\{0\})$. Differentiating that identity
in $z$ for $z\not\in \overline{U}\cup\{0\}$ gives
\[
\nabla \Psi_U(z)=c(\lambda,k)\,\frac{z}{|z|^2},
\qquad z\in \C\setminus(\overline{U}\cup\{0\}).
\]
The classical jump relations for the single-layer / volume logarithmic
potential of an $L^\infty$ density on a $C^2$ domain (see, e.g.,
\cite{Gustafsson}) imply that $\nabla\Psi_U$ extends continuously
across $\partial U$ from \emph{both} sides, and the boundary values
agree. Hence on $\partial U$,
\[
\nabla \Psi_U(z)=c(\lambda,k)\,\frac{z}{|z|^2},
\qquad z\in \partial U,
\]
which is, in particular, parallel to $z$. Therefore the difference
\[
v:=\Psi_U-\Phi
\]
satisfies $\Delta v=2\pi f(|z|)\mathbf 1_U-2\pi f(|z|)\mathbf 1_B
=-2\pi f(|z|)\mathbf 1_{B\setminus U}$ in the sense of distributions
on $B$, so $v$ is harmonic in $U$ and in $B\setminus\overline{U}$, and
on $\partial U$ satisfies
\[
\nabla v(x)\ \text{is parallel to}\ x.
\]
The next lemma converts this gradient condition into local radiality of $v$.

\begin{lemma}\label{lem:harmonic-locally-radial}
Let $V\subset \R^2$ be a domain, and suppose there exists a harmonic
function $v\in C^1(\overline{V})$ such that $\nabla v(x)$ is parallel to
$x$ for every $x\in \partial V$. Then $\nabla v(x)$ is parallel to $x$
for every $x\in V$. In particular, on every circle
$S=\{|x|=r\}\subset \R^2$, the function $v$ is constant on every
connected component of $S\cap V$.
\end{lemma}

\begin{proof}
Consider the angular derivative
\[
g(x):=x_1\,\partial_2 v(x)-x_2\,\partial_1 v(x)
=\partial_\theta v(x),
\]
which equals the tangential derivative of $v$ along the circle
through $x$ centered at $0$. A direct computation gives
\[
\Delta g
=x_1\,\Delta(\partial_2 v)-x_2\,\Delta(\partial_1 v)
+2\big(\partial_1\partial_2 v-\partial_2\partial_1 v\big)
=0,
\]
since $v$ is harmonic and partial derivatives of harmonic functions
are themselves harmonic. By the boundary hypothesis, $\nabla v(x)$ is
parallel to $x$ on $\partial V$, so $g\equiv 0$ on $\partial V$.
Together with $g\in C(\overline{V})$, the maximum principle for
harmonic functions yields $g\equiv 0$ in $V$. Hence $\nabla v(x)$ is
parallel to $x$ at every interior point as well, and on any circle
$S=\{|x|=r\}$ the tangential derivative $\partial_\theta v$ vanishes
on $S\cap V$. Therefore $v$ is constant on each connected component
of $S\cap V$.
\end{proof}

Applying Lemma~\ref{lem:harmonic-locally-radial} on small subdomains
$V\subset U$ where $v$ is harmonic and $C^1$ up to the boundary, we
conclude that $v$ is locally radial on $U$. Since $\Phi$ is globally
radial, so is $\Psi_U=\Phi+v$, locally on $U$.

\begin{lemma}\label{lem:local-to-global-radial}
Let \(W\subset \R^2\setminus\{0\}\) be a connected open set, and let
\(u\in C(W)\). Suppose that for every \(x\in W\) there exist \(r_x>0\)
and a scalar function \(\phi_x\) such that
\[
u(y)=\phi_x(|y|)
\qquad\text{for all } y\in B(x,r_x)\subset W.
\]
Then there exists a single scalar function \(\phi\), defined on the
interval
\[
I_W:=\{|y|:y\in W\},
\]
such that
\[
u(y)=\phi(|y|)
\qquad\text{for all } y\in W.
\]
\end{lemma}

\begin{proof}
For each \(x\in W\), choose a ball \(B_x:=B(x,r_x)\subset W\) on which
\(u\) is a radial function of \(|y|\), with profile $\phi_x$. We
first note that whenever $B_x\cap B_{x'}\ne\emptyset$,
\[
\phi_x(|y|)=u(y)=\phi_{x'}(|y|)
\qquad\text{for every }y\in B_x\cap B_{x'}.
\]
Since $B_x\cap B_{x'}$ is open and avoids the origin, the map $y\mapsto
|y|$ is a submersion on it, so the set of radii it attains contains a
nontrivial open interval on which $\phi_x$ and $\phi_{x'}$ agree.

Fix $x_*\in W$ and define a candidate function $\phi$ on $I_W$ as
follows. Given $y\in W$, connect $x_*$ to $y$ by a polygonal path in
$W$, which exists because $W$ is open and connected. Its image is
compact, hence covered by finitely many balls
$B_{x_1},\dots,B_{x_m}$ with successive nonempty overlaps and
$x_1=x_*$, $y\in B_{x_m}$. The previous paragraph implies $\phi_{x_j}$
and $\phi_{x_{j+1}}$ agree on a nontrivial interval of radii, so by
analytic-continuation along the chain we may unambiguously set
$\phi(|y|):=\phi_{x_m}(|y|)$. Independence of the chain follows from
the same overlap argument: any two chains from $x_*$ to $y$ can be
interleaved into a single chain, and successive overlaps force the
final value to coincide. Hence $\phi$ is well-defined on $I_W$ and
satisfies $u(y)=\phi(|y|)$ for every $y\in W$.
\end{proof}

We now conclude the proof. Since \(U\) is simply connected and bounded,
it suffices to rule out the possibility that \(U\) fails to be a disk
centered at the origin. Suppose therefore, for contradiction, that
\(U\) is not such a disk. Then \(\partial U\) is not a single circle
centered at $0$, so there exist radii \(0<r_1<r_2\) such that the
annulus
\[
A:=\{z\in \R^2:\ r_1<|z|<r_2\}
\]
meets both $U$ and $\partial U$ in a non-trivial way; in particular,
we may pick a point
\[
x_0\in \partial U\cap A.
\]
Let
\[
\widetilde{U}_0:=U\cap B_r(x_0),
\]
for $r>0$ sufficiently small that $\widetilde U_0$ is connected and
$0\notin \overline{\widetilde U_0}$. Since \(\Psi_U\) is locally radial
in \(U\), Lemma~\ref{lem:local-to-global-radial} applied to
$W=\widetilde U_0$ yields a scalar function \(\psi_U\) on
$I_{\widetilde U_0}$ such that
\[
\Psi_U(x)=\psi_U(|x|),\qquad x\in \widetilde{U}_0.
\]

On the other hand, by \eqref{eq:potential} and the continuity of the
gradient of a volume logarithmic potential with bounded density across
a $C^2$ boundary (see, e.g., \cite{Gustafsson,Isakov}), the gradient
$\nabla\Psi_U$ extends continuously across $\partial U$, and for
$z\in\partial U$ we have
\[
\nabla \Psi_U(z)=c\,\frac{z}{|z|^2}
\]
for the same constant $c=c(\lambda,k)\in \R$ as in
\eqref{eq:potential}. Combining this with $\Psi_U(x)=\psi_U(|x|)$ on
$\widetilde U_0$ and passing to $z\in \partial \widetilde U_0\cap
\partial U$ from the inside,
\[
\psi_U'(|z|)\,\frac{z}{|z|}=c\,\frac{z}{|z|^2},
\qquad z\in \partial U\cap \overline{B_r(x_0)},
\]
hence
\[
\psi_U'(r)=\frac{c}{r}
\]
for every $r$ in the set of radii attained by points of
$\partial U\cap \overline{B_r(x_0)}$. Since $\partial U$ is a $C^2$
curve, we may, after shrinking $r$ once more, parametrize
$\partial U\cap B_r(x_0)$ by a $C^2$ embedding
$\gamma:(-\varepsilon,\varepsilon)\to \R^2$ with $\gamma(0)=x_0$.
Because $\partial U$ is not a circle centered at $0$, the smooth
function $t\mapsto |\gamma(t)|^2$ is non-constant on the connected
$C^2$ arc $\partial U\cap A$. Hence its derivative does not vanish
identically on any subinterval, and we may choose $x_0\in\partial U\cap A$
so that
\[
\frac{d}{dt}|\gamma(t)|^2\Big|_{t=0}\ne 0,
\]
which is equivalent to $\frac{d}{dt}|\gamma(t)|\big|_{t=0}\ne 0$ since
$|\gamma(0)|=|x_0|>0$. Hence by the open-mapping
property for $C^1$ functions of one variable with nonzero derivative,
the image of $|\gamma|$ contains a nontrivial open interval
\(I\subset (r_1,r_2)\).
Therefore
\[
\psi_U'(r)=\frac{c}{r}
\qquad\text{for every }r\in I,
\]
whence
\[
\psi_U(r)=c\log r+C_0,
\qquad r\in I,
\]
for some real constant $C_0$.

Let
\[
W_I:=\{x\in \widetilde U_0:\ |x|\in I\}.
\]
Then \(W_I\) is a nonempty open subset of \(U\), and on \(W_I\),
\[
\Psi_U(x)=c\log|x|+C_0.
\]
Hence
\[
\Delta \Psi_U\equiv 0 \qquad \text{in } W_I,
\]
since the origin is excluded from $\widetilde U_0$. On the other hand,
since $0\notin \overline{U}$, the logarithmic potential $\Psi_U$ has
distributional Laplacian
\[
\Delta \Psi_U=2\pi\, f(|z|)\,\mathbf{1}_U
\]
in $\C$, with no contribution at the origin. Hence on the open set
$W_I\subset U\setminus\{0\}$ we have
\[
\Delta \Psi_U=2\pi\, f(|z|)=2\pi\, |z|^{2k}e^{-\pi |z|^2}.
\]
The right-hand side vanishes only at the origin, which lies outside
$W_I$, while the left-hand side just shown to be $0$ on $W_I$ — a
contradiction. This rules out the case in which $U$ is not a disk
centered at $0$, completing the proof.
\end{proof}

\begin{remark}
The structure of the argument above is in the spirit of
\cite{Aharonov-Schiffer-Zalcman}, where logarithmic-potential
rigidity statements were derived from Schiffer-style overdetermined
boundary conditions, and of the quadrature-domain literature
\cite{Gustafsson,Gustafsson-Shahgholian,Shahgholian,Karp}, in which
the regularity of free boundaries and the matching of an interior
potential with the Newtonian potential of a radially symmetric
density play a central role. Concretely, identity
\eqref{eq:potential} can be read as saying that $U$ is a
\emph{generalized quadrature domain} for the radial weight $f(|z|)$,
with quadrature data concentrated at the origin. From this perspective,
the conclusion that $U$ is a disk is a (sharp) instance of the
classical principle that radial quadrature data forces a radial
domain. The role of the $C^2$ boundary regularity is concentrated
in the boundary extension of $\nabla\Psi_U$ and in the
open-interval-of-radii argument that ends the proof; both can be
relaxed at the cost of additional technicalities, which we will
revisit in the second proof.
\end{remark}


\subsection{Second proof}

\begin{proof}[Second proof of Theorem \ref{thm:eig-loc-U}]
We keep the notation
\[
f(|z|):=|z|^{2k}e^{-\pi|z|^2},
\qquad
c:=c(\lambda,k),
\]
from the first proof. The argument below is in the spirit of the holomorphic
moving-plane/quadrature-domain methods of \cite{Aharonov-Schiffer-Zalcman,Gustafsson,Karp},
in that we exploit the existence of a holomorphic function on \(U\) with prescribed real
boundary values to recover full radial symmetry. Starting again from \eqref{eq:fourier-delta} and
applying Lemma~\ref{lemma:fourier-to-fboundary}, we obtain a compactly supported distributional
solution of
\begin{equation}\label{eq:potato-second-proof}
\Delta u=f(|z|)\mathbf{1}_U-c\,\delta_0
\qquad\text{in }\mathcal D'(\R^2).
\end{equation}
Using the planar fundamental solution \((2\pi)^{-1}\log|z|\), we may write
\begin{equation}\label{eq:explicit-u-second-proof}
u(z):=\frac{1}{2\pi}\int_U \log|z-z'|\,f(|z'|)\,dA(z')-\frac{c}{2\pi}\log|z|.
\end{equation}
By \eqref{eq:potential}, this function satisfies
\[
u\equiv 0
\qquad\text{on }\R^2\setminus (U\cup\{0\}).
\]

We first show that \(0\in U\). Since the singular term in
\eqref{eq:potato-second-proof} is \(-c\delta_0\), one must have \(0\in\overline U\). If
\(0\in\partial U\), then there exist exterior points \(z_j\to 0\) with
\[
u(z_j)=0.
\]
On the other hand, the first term in \eqref{eq:explicit-u-second-proof} remains finite as
\(z\to 0\), because \(f\in L^\infty(U)\) and \(\log|z-z'|\) is locally integrable, whereas the
second term equals \(-\frac{c}{2\pi}\log|z|\to +\infty\). This contradicts the exterior identity
\(u(z_j)=0\). Hence \(0\notin\partial U\), and therefore
\[
0\in U.
\]

We next show directly from \eqref{eq:explicit-u-second-proof} that
\[
u\in C^1(\R^2\setminus\{0\}).
\]
Indeed, if we write
\[
g(z):=f(|z|)\mathbf{1}_U(z)\in L^\infty_c(\R^2),
\]
then
\[
u(z)=\frac{1}{2\pi}\int_{\R^2}\log|z-z'|\,g(z')\,dA(z')-\frac{c}{2\pi}\log|z|.
\]
Since \(g\in L^\infty_c(\R^2)\), the only contribution to the integral defining the convolution
\((K\ast g)(z)\) below comes from a fixed bounded set, so we may replace \(K\) by its restriction
\(K\mathbf{1}_{B_R}\) for any sufficiently large ball \(B_R\). The kernel
\(K(\zeta)=\zeta/|\zeta|^2\) satisfies \(|K(\zeta)|=|\zeta|^{-1}\), which is integrable on bounded
subsets of \(\R^2\); hence \(K\mathbf{1}_{B_R}\in L^1(\R^2)\). It is a standard fact
that the convolution of an \(L^1\) function with a bounded function is continuous
(this follows from the density of \(C_c(\R^2)\) in \(L^1(\R^2)\) and continuity of translations in
\(L^1\)), so
\[
(K\ast g)(z):=\int_{\R^2}\frac{z-z'}{|z-z'|^2}\,g(z')\,dA(z')
\]
is a continuous function of \(z\in\R^2\). On the other hand, a direct distributional computation
shows
\[
\nabla\!\left(\int_{\R^2}\log|z-z'|\,g(z')\,dA(z')\right)=K\ast g
\qquad\text{in }\mathcal D'(\R^2),
\]
so the gradient of the convolution \(\log|\cdot|\ast g\) coincides almost everywhere with the
continuous function \(K\ast g\). It follows that \(\log|\cdot|\ast g\) is of class \(C^1\) on
\(\R^2\), and hence \(u\in C^1(\R^2\setminus\{0\})\) (the only obstruction to global \(C^1\)
regularity being the explicit logarithmic singularity at the origin coming from the
\(-\frac{c}{2\pi}\log|z|\) term).

Now let \(z_*\in\partial U\). Since \(0\in U\), the boundary point \(z_*\) is bounded away from
\(0\); fix a small ball \(B_\rho(z_*)\Subset \R^2\setminus\{0\}\). On the exterior open set
\(\R^2\setminus \overline U\) we have \(u\equiv 0\), so \(\nabla u\equiv 0\) there. As
\(\partial U\) is contained in the boundary of the exterior, every \(z_*\in\partial U\) is a limit
of exterior points \(z_j\to z_*\) along which \(\nabla u(z_j)=0\). By the \(C^1\) regularity of
\(u\) on \(B_\rho(z_*)\), the gradient \(\nabla u\) is continuous at \(z_*\), and we conclude
\(\nabla u(z_*)=0\). In complex notation \(\partial_z=\tfrac{1}{2}(\partial_x-i\partial_y)\), this
yields
\begin{equation}\label{eq:dz-zero-second-proof}
\partial_z u=0
\qquad\text{on }\partial U.
\end{equation}
We emphasize that \emph{no} boundary regularity of \(\partial U\) has been used at this point, in
contrast with the first proof: the vanishing of \(\nabla u\) on \(\partial U\) is forced solely by
exterior continuity, which in turn is supplied by the explicit
formula~\eqref{eq:explicit-u-second-proof}.

Define
\[
\Phi(t):=\int_0^t s^k e^{-\pi s}\,ds
\qquad\text{and}\qquad
\Psi(z):=\frac{\Phi(|z|^2)}{z},
\quad z\neq 0.
\]
A direct computation, using \(|z|^2=z\bar z\) and \(\partial_{\bar z}(z\bar z)=z\), gives, for
\(z\neq 0\),
\[
\partial_{\bar z}\Psi(z)
=
\frac{\Phi'(|z|^2)\,\partial_{\bar z}(|z|^2)}{z}
=
\frac{|z|^{2k}e^{-\pi|z|^2}\cdot z}{z}
=
|z|^{2k}e^{-\pi|z|^2}
=
f(|z|).
\]
At the origin, the apparent singularity \(1/z\) is integrable against the vanishing factor
\(\Phi(|z|^2)\): since \(k\ge 0\) and \(s^k e^{-\pi s}\le s^k\) for \(s\ge 0\), we have
\(\Phi(t)=\int_0^t s^k e^{-\pi s}\,ds\le \int_0^t s^k\,ds = t^{k+1}/(k+1)=O(t^{k+1})\) as
\(t\to 0\). Hence \(\Psi(z)=O(|z|^{2k+1})\) near \(z=0\), so \(\Psi\in L^\infty_{\mathrm{loc}}\)
(in fact continuous, with \(\Psi(0)=0\)). Consequently, the identity
\(\partial_{\bar z}\Psi=f(|z|)\) persists in the distributional sense on all of \(U\); no extra
\(\delta_0\) contribution arises at the origin.

Recall also the standard identity
\[
\partial_{\bar z}\!\left(\frac{1}{z}\right)=\pi\delta_0
\qquad\text{in }\mathcal D'(\R^2),
\]
which is equivalent to the fact that \(\frac{1}{\pi z}\) is a fundamental solution of
\(\partial_{\bar z}\); equivalently, \(\Delta\log|z|=4\partial_z\partial_{\bar z}\log|z|=2\pi\delta_0\)
combined with \(\partial_z\log|z|=\frac{1}{2z}\). Now set
\[
F(z):=4\partial_z u(z)-\Psi(z)+\frac{c}{\pi z},
\qquad z\in U\setminus\{0\}.
\]
We first check that \(F\) defines a distribution on all of \(U\) (a~priori it is given only on
\(U\setminus\{0\}\)). From \eqref{eq:explicit-u-second-proof} we have, near \(z=0\),
\[
4\partial_z u(z)
=
\frac{2}{\pi}\partial_z\!\left(\int_{\R^2}\log|z-z'|\,g(z')\,dA(z')\right)-\frac{c}{\pi z},
\]
where the first term is locally bounded near \(z=0\) (the Cauchy-type integral
\(\int (z-z')^{-1} g(z')\,dA(z')\) of a bounded compactly supported function is locally
integrable, by the same \(L^1_{\mathrm{loc}}\)-kernel argument as above). Adding \(-\Psi\)
(continuous with \(\Psi(0)=0\)) and \(c/(\pi z)\) cancels the explicit \(1/z\)-singularity, so
\(F\) is locally bounded near the origin and hence defines a distribution on \(U\).

Using \(\Delta=4\partial_z\partial_{\bar z}\), the relation \(\partial_{\bar z}(1/z)=\pi\delta_0\),
the computation of \(\partial_{\bar z}\Psi\) above, and \eqref{eq:potato-second-proof}, we
compute, in \(\mathcal D'(U)\),
\[
\partial_{\bar z}F
=
4\partial_z\partial_{\bar z}u-\partial_{\bar z}\Psi+\frac{c}{\pi}\partial_{\bar z}\!\left(\frac{1}{z}\right)
=
\Delta u-f(|z|)+c\,\delta_0
=
\bigl(f(|z|)-c\delta_0\bigr)-f(|z|)+c\delta_0
=0.
\]
The singular term \(-c\delta_0\) in \eqref{eq:potato-second-proof}, arriving through
\(4\partial_z\partial_{\bar z}u=\Delta u\), is thus exactly cancelled by
\((c/\pi)\partial_{\bar z}(1/z)=c\delta_0\). Hence \(F\) is a distributional solution of the
Cauchy--Riemann equation \(\partial_{\bar z}F=0\) on \(U\), and by Weyl's lemma for
\(\partial_{\bar z}\) (equivalently, hypoellipticity of the Cauchy--Riemann operator) it is
represented by a genuine holomorphic function on \(U\).

Define
\[
W(z):=\pi z\,F(z).
\]
Then \(W\) is holomorphic on \(U\) and \(W(0)=0\). Since \(\partial_z u\) is continuous up to
\(\partial U\), so is \(W\), and by \eqref{eq:dz-zero-second-proof},
\[
W|_{\partial U}
=
\pi z\left(4\partial_z u-\frac{\Phi(|z|^2)}{z}+\frac{c}{\pi z}\right)\Big|_{\partial U}
=
c-\pi\Phi(|z|^2).
\]
In particular, \(W|_{\partial U}\) is real-valued. The imaginary part of \(W\) is therefore
harmonic in \(U\), continuous up to \(\partial U\), and vanishes on \(\partial U\). By the
maximum principle,
\[
\Im W\equiv 0
\qquad\text{in }U.
\]
Since a holomorphic function with zero imaginary part is constant, there exists \(\kappa\in\R\)
such that
\[
W\equiv \kappa
\qquad\text{in }U.
\]

Restricting to the boundary, and using that \(W\) is continuous up to \(\partial U\) with
boundary value \(c-\pi\Phi(|z|^2)\), we obtain
\[
c-\pi\Phi(|z|^2)=\kappa
\qquad\text{for all }z\in\partial U.
\]
The function \(t\mapsto \Phi(t)=\int_0^t s^k e^{-\pi s}\,ds\) has derivative
\(\Phi'(t)=t^k e^{-\pi t}>0\) for \(t>0\), hence is strictly increasing on \([0,\infty)\). The
identity above therefore forces \(|z|^2\) to take a single value on \(\partial U\), i.e.\ \(|z|\)
is constant on \(\partial U\). Consequently \(\partial U\) is contained in a circle centered at
the origin. Let \(R>0\) be such that \(\partial U\subset \{|z|=R\}\). Since \(U\) is open and
\(U\cap\partial U=\varnothing\), we have \(U\subset \{|z|<R\}\cup\{|z|>R\}\); since \(U\) is
connected and the two pieces are disjoint open sets, \(U\) must lie entirely in one of them, and
boundedness of \(U\) forces \(U\subset \{|z|<R\}\). If this inclusion were strict, pick
\(z_1\in\{|z|<R\}\setminus U\) and \(z_0\in U\). The straight segment from \(z_0\) to \(z_1\) is
contained in the convex set \(\{|z|<R\}\), and as it joins a point of \(U\) to a point of its
complement it must cross \(\partial U\) at some \(z_*\); but then \(z_*\in\partial U\cap\{|z|<R\}\),
contradicting \(\partial U\subset \{|z|=R\}\). Therefore
\[
U=\{|z|<R\}
\]
that is, \(U\) is a disk.
\end{proof}

\begin{remark}
The two proofs use the boundary assumptions in rather different ways, and it is instructive to
keep careful track of the regularity bookkeeping.

In the first proof, once the logarithmic potential \(\Psi_U\) is written explicitly, no elliptic
boundary regularity is needed to conclude that \(\Psi_U\in C^1(\R^2\setminus\{0\})\): this follows
directly from convolution with the locally integrable kernel \(\nabla\log|z|\), exactly as in the
\(L^1_{\mathrm{loc}}\)-translation argument used here for \eqref{eq:explicit-u-second-proof}. Thus
the local radiality argument itself only uses continuity of the gradient up to the boundary. The
genuine extra input appears only in the final geometric step, where one needs that, near a suitable
point \(x_0\in \partial U\), the set \(\partial U\cap B_r(x_0)\) contains a nontrivial boundary
arc, so that the radii attained there form an interval. Accordingly, the first proof does not
really need \(C^2\) regularity; it works under any hypothesis guaranteeing such a local arc
structure, for instance when \(U\) is a Jordan domain. In particular, any Lipschitz, \(C^1\), or
\(C^2\) boundary is more than sufficient.

By contrast, the second proof does not use any boundary regularity at all. The explicit formula
\eqref{eq:explicit-u-second-proof} already yields \(u\in C^1(\R^2\setminus\{0\})\), and since
\(u\equiv 0\) on the exterior open set \(\R^2\setminus \overline U\), continuity implies that both
\(u\) and \(\nabla u\) vanish on \(\partial U\) -- precisely the boundary information packaged in
\eqref{eq:dz-zero-second-proof}. This is exactly the input needed to construct the holomorphic
function \(W=\pi z F\) and conclude that \(|z|\) is constant on \(\partial U\). Hence the second
proof recovers the theorem under the original assumptions of Abreu and D\"orfler
\cite{Abreu-Doerfler}, namely for arbitrary bounded simply connected domains, with no boundary
smoothness hypothesis. The mechanism is also conceptually closer to the classical
quadrature-domain rigidity results of \cite{Aharonov-Schiffer-Zalcman,Gustafsson,Karp}, where the
existence of a holomorphic function on \(U\) with real boundary values forces the boundary to lie
on a level set of a real-analytic function, here the elementary radial function
\(z\mapsto |z|^2\).
\end{remark}

% ==============================================================
% Section 5
% ==============================================================

\section{Proof of Theorem \ref{thm:GGRT-sharp}}\label{sec:GGRT-sharp}

We prove that the quantitative stability estimate in Corollary~1.4 of
\cite{Gomez-Guerra-Ramos-Tilli}, which builds on the Faber--Krahn rigidity
result of \cite{Nicola-Tilli} characterizing the disk as the unique optimizer
of the first localization eigenvalue (see also the foundational work of
\cite{Daubechies} on the spectrum of the Daubechies localization operator on
the disk), is sharp. The proof proceeds by constructing a family of domains
which are perturbations of the disk, whose first localization eigenvalue has
the same first-order behaviour as the disk, but whose distance to the class of
disks is of order $\varepsilon$.

\begin{proof}[Proof of Theorem~\ref{thm:GGRT-sharp}]
Fix
\[
\lambda_0:=1-e^{-\pi},
\]
which is the eigenvalue of the localization operator of the unit disk $D_1$
associated with the ground state $h_0$. For $\varepsilon$ small, define
\[
P_\varepsilon(w):=1+\varepsilon w^2.
\]
By Theorem~\ref{thm:perturb}, applied with \(\lambda=\lambda_0\) and
\(P_\varepsilon(w)=1+\varepsilon w^2\), for all sufficiently small
\(|\varepsilon|\) there exists a domain \(U_\varepsilon\) such that
\begin{equation}\label{eq:functiona-P-eps-new}
\int_{U_\varepsilon} P_\varepsilon(z)e^{\pi \overline zw}e^{-\pi|z|^2}\,dA(z)
=
\lambda_0 P_\varepsilon(w),
\qquad w\in\C.
\end{equation}
Moreover, by construction, $U_\varepsilon$ is a $C^2$ perturbation of the unit
disk. Hence, after possibly shrinking the range of admissible $\varepsilon$,
its boundary can be written as
\begin{equation}\label{eq:def-u-eps-new}
\partial U_\varepsilon
=
\{(1+\varphi_\varepsilon(y))y:\ y\in \mathbb S^1\},
\end{equation}
where
\[
\varphi_\varepsilon\in C^2(\mathbb S^1),
\qquad
\|\varphi_\varepsilon\|_{C^2(\mathbb S^1)}\lesssim |\varepsilon|.
\]
The map
\[
\varepsilon\mapsto \varphi_\varepsilon
\]
may in fact be chosen \(C^1\) as a map from a neighborhood of \(0\) into
\(C^1(\mathbb S^1)\). To see this, we briefly retrace the proof of
Theorem~\ref{thm:perturb}. The construction is the disk moment-map
construction of Proposition~\ref{prop:disk-inverse-tailored}. Thus
\(\varphi_\varepsilon\) is obtained from an equation of the form
\(F(\varepsilon,\varphi)=0\), where
\[
F:(-\varepsilon_0,\varepsilon_0)\times \mathcal U\to \mathcal V
\]
is the \(C^1\) moment map between the analytic boundary spaces used there,
with \(F(0,0)=0\) and \(\partial_\varphi F(0,0)\) invertible by the explicit
Fourier diagonalization at the unit disk. Since the polynomial
\(P_\varepsilon(z)=1+\varepsilon z^2\)
depends polynomially on \(\varepsilon\), the dependence of \(F\) on
\(\varepsilon\) is \(C^\infty\). The \(C^1\) version of the implicit function
theorem in Banach spaces then yields a \(C^1\) branch
\(\varepsilon\mapsto \varphi_\varepsilon\) into \(C^1(\mathbb S^1)\). We
therefore write
\[
\dot\varphi_0(y):=
\left.\partial_\varepsilon\right|_{\varepsilon=0}\varphi_\varepsilon(y).
\]

\medskip
\noindent
{\bf Step 1: first variation of the eigenvalue identity.} We differentiate \eqref{eq:functiona-P-eps-new} at $\varepsilon=0$.
Since \(P_\varepsilon\to 1\) in \(C^\infty\) and
\(\varepsilon\mapsto \varphi_\varepsilon\) is \(C^1\) into \(C^1(\mathbb S^1)\),
all terms in \eqref{eq:functiona-P-eps-new} are differentiable at
\(\varepsilon=0\). Since $\partial U_\varepsilon$ is $C^2$ and admits the
explicit normal graph parametrization \eqref{eq:def-u-eps-new} over
$\mathbb S^1$ with $\varphi_\varepsilon$ depending $C^1$-smoothly on
$\varepsilon$, the Reynolds transport identity reduces to a direct application
of the fundamental theorem of calculus in polar coordinates. Concretely, with
\(z=re^{it}\), for any smooth test function \(H\) we have
\[
\int_{U_\varepsilon} H(z)\,dA(z)
=
\int_0^{2\pi}\int_0^{1+\varphi_\varepsilon(e^{it})}
H(re^{it})\,r\,dr\,dt.
\]
Differentiating in \(\varepsilon\) at \(\varepsilon=0\) and using the
fundamental theorem of calculus,
\[
\left.\frac{d}{d\varepsilon}\right|_{\varepsilon=0}
\int_{U_\varepsilon} H(z)\,dA(z)
=
\int_0^{2\pi}\dot\varphi_0(e^{it})H(e^{it})\,dt
=
\int_{\mathbb S^1}\dot\varphi_0(z)H(z)\,d\mathcal H^1(z),
\]
where the factor \(r=1\) on the unit circle has been absorbed into the
arc-length measure. This is precisely Reynolds' transport identity:
\[
\left.\frac{d}{d\varepsilon}\right|_{\varepsilon=0}
\int_{U_\varepsilon} H(z)\,dA(z)
=
\int_{\mathbb S^1} \dot\varphi_0(z)H(z)\,d\mathcal H^1(z).
\]
Applying this to
\[
H(z)=e^{\pi \overline zw}e^{-\pi|z|^2},
\]
and differentiating the factor $P_\varepsilon(z)=1+\varepsilon z^2$, we obtain
\begin{equation}\label{eq:first-variation-identity-new}
\int_{\mathbb S^1}
\dot\varphi_0(z)e^{\pi \overline zw}e^{-\pi|z|^2}\,d\mathcal H^1(z)
+
\int_{D_1} z^2 e^{\pi \overline zw}e^{-\pi|z|^2}\,dA(z)
=
\lambda_0 w^2.
\end{equation}
Since $D_1$ is radial, the monomial $z^2$ is an eigenfunction of the
localization operator of $D_1$; therefore there exists $\lambda_2>0$ such that
\[
\int_{D_1} z^2 e^{\pi \overline zw}e^{-\pi|z|^2}\,dA(z)=\lambda_2 w^2.
\]
Hence \eqref{eq:first-variation-identity-new} becomes
\begin{equation}\label{eq:boundary-identity-new}
\int_{\mathbb S^1}
\dot\varphi_0(z)e^{\pi \overline zw}e^{-\pi|z|^2}\,d\mathcal H^1(z)
=
(\lambda_0-\lambda_2)w^2.
\end{equation}
In fact, \(\lambda_2\) can be computed explicitly. Expanding
\(e^{\pi \overline zw}=\sum_{n\geq 0}\frac{\pi^n}{n!}\overline z^{\,n}w^n\)
and using orthogonality of the monomials \(\overline z^{\,n}\) over the disk,
the only term contributing to \(w^2\) is the one with \(n=2\):
\[
\int_{D_1} z^2 e^{\pi \overline zw}e^{-\pi|z|^2}\,dA(z)
=
\frac{\pi^2}{2}w^2\int_{D_1}z^2\overline z^{\,2}e^{-\pi|z|^2}\,dA(z)
=
\frac{\pi^2}{2}w^2\int_{D_1}|z|^4 e^{-\pi|z|^2}\,dA(z).
\]
Passing to polar coordinates and integrating by parts twice,
\[
\int_{D_1}|z|^4 e^{-\pi|z|^2}\,dA(z)
=
2\pi\int_0^1 r^4 e^{-\pi r^2}\,r\,dr
=
\pi\int_0^1 s^2 e^{-\pi s}\,ds
\]
where we substituted \(s=r^2\). A direct computation, integrating by parts
twice, yields
\[
\pi\int_0^1 s^2 e^{-\pi s}\,ds
=
\frac{2}{\pi^2}\Bigl(1-\Bigl(1+\pi+\frac{\pi^2}{2}\Bigr)e^{-\pi}\Bigr).
\]
Multiplying by \(\frac{\pi^2}{2}\), we obtain
\[
\lambda_2
=
\frac{\pi^2}{2}\int_{D_1}|z|^4e^{-\pi|z|^2}\,dA(z)
=
1-\Bigl(1+\pi+\frac{\pi^2}{2}\Bigr)e^{-\pi}.
\]
Therefore
\[
\lambda_0-\lambda_2
=
(1-e^{-\pi})-\Bigl(1-\Bigl(1+\pi+\frac{\pi^2}{2}\Bigr)e^{-\pi}\Bigr)
=
\pi\Bigl(1+\frac{\pi}{2}\Bigr)e^{-\pi}>0,
\]
so the coefficient of \(w^2\) on the right-hand side is indeed nonzero.
In particular, since $|z|=1$ on $\mathbb S^1$, we may rewrite this as
\begin{equation}\label{eq:boundary-identity-circle-new}
e^{-\pi}
\int_{\mathbb S^1}
\dot\varphi_0(z)e^{\pi \overline zw}\,d\mathcal H^1(z)
=
(\lambda_0-\lambda_2)w^2.
\end{equation}

Setting $w=0$ gives
\begin{equation}\label{eq:cancellation-f-0-new}
\int_{\mathbb S^1}\dot\varphi_0\,d\mathcal H^1=0.
\end{equation}

\medskip
\noindent
{\bf Step 2: identification of the leading mode.} Expanding
\[
e^{\pi \overline zw}
=
\sum_{n\geq 0}\frac{\pi^n}{n!}\overline z^{\,n}w^n,
\]
and matching the coefficient of $w^n$ on both sides of
\eqref{eq:boundary-identity-circle-new}, we find
\[
\int_{\mathbb S^1}\dot\varphi_0(z)\overline z^{\,n}\,d\mathcal H^1(z)=0
\qquad \text{for all } n\neq 2,
\]
whereas the coefficient corresponding to $n=2$ equals
$\frac{2}{\pi^2}(\lambda_0-\lambda_2)e^\pi\neq 0$. Since on $\mathbb S^1$
we have $\overline z=e^{-it}$ when $z=e^{it}$, this says that the Fourier
coefficient $\widehat{\dot\varphi_0}(n)$ vanishes for every $n\neq 2$. Since
$\dot\varphi_0$ is real-valued, $\widehat{\dot\varphi_0}(-n)
=\overline{\widehat{\dot\varphi_0}(n)}$, so the $n=-2$ mode is also nonzero
and conjugate to the $n=2$ mode, while all other modes vanish. We conclude
that
\begin{equation}\label{eq:mode-decomposition-new}
\dot\varphi_0(e^{it})
=
a\cos(2t)+b\sin(2t)
\end{equation}
for some pair $(a,b)\neq (0,0)$. In particular,
\begin{equation}\label{eq:l1-nontrivial-new}
\int_{\mathbb S^1} |\dot\varphi_0|\,d\mathcal H^1>0.
\end{equation}

\medskip
\noindent
{\bf Step 3: normalization of the area.} The area of $U_\varepsilon$ is given by
\begin{equation}\label{eq:area-u-eps-new}
|U_\varepsilon|
=
\frac12\int_{\mathbb S^1} (1+\varphi_\varepsilon(y))^2\,d\mathcal H^1(y).
\end{equation}
Since
\[
\varphi_\varepsilon
=
\varepsilon \dot\varphi_0 + O(\varepsilon^2)
\quad \text{in } C^1(\mathbb S^1),
\]
and \eqref{eq:cancellation-f-0-new} gives the vanishing of the first-order
term, we obtain from \eqref{eq:area-u-eps-new}
\begin{equation}\label{eq:measure-close-u-eps-new}
|U_\varepsilon|=\pi+O(\varepsilon^2).
\end{equation}

Let
\[
r_\varepsilon:=\sqrt{\frac{\pi}{|U_\varepsilon|}},
\qquad
\widetilde U_\varepsilon:=r_\varepsilon U_\varepsilon.
\]
Then $|\widetilde U_\varepsilon|=\pi$, and by
\eqref{eq:measure-close-u-eps-new},
\[
r_\varepsilon=1+O(\varepsilon^2).
\]
Hence $\widetilde U_\varepsilon$ still admits a normal graph representation
\[
\partial \widetilde U_\varepsilon
=
\{(1+\widetilde\varphi_\varepsilon(y))y:\ y\in \mathbb S^1\},
\]
with
\[
\widetilde\varphi_\varepsilon
=
\varepsilon \dot\varphi_0 + O(\varepsilon^2)
\quad \text{in } C^1(\mathbb S^1).
\]
Thus the area normalization does not alter the first-order deformation.

\medskip
\noindent
{\bf Step 4: distance from the centered disk.} Since both $\widetilde U_\varepsilon$ and $D_1$ are star-shaped with respect to
the origin, the symmetric difference may be computed in polar coordinates:
\[
|\widetilde U_\varepsilon \triangle D_1|
=
\frac12\int_{\mathbb S^1}
\left|(1+\widetilde\varphi_\varepsilon(y))^2-1\right|
\,d\mathcal H^1(y).
\]
Since \(\widetilde\varphi_\varepsilon=\varepsilon\dot\varphi_0+O(\varepsilon^2)\)
in \(C^1(\mathbb S^1)\), we have in particular
\(\|\widetilde\varphi_\varepsilon\|_{L^\infty(\mathbb S^1)}\lesssim |\varepsilon|\),
whence \(\widetilde\varphi_\varepsilon^2=O(\varepsilon^2)\) in
\(L^\infty(\mathbb S^1)\). Hence
\[
(1+\widetilde\varphi_\varepsilon)^2-1
=
2\widetilde\varphi_\varepsilon+\widetilde\varphi_\varepsilon^2
=
2\varepsilon \dot\varphi_0 + O(\varepsilon^2)
\quad \text{in } L^\infty(\mathbb S^1),
\]
where the implied constant depends only on the \(C^1\)-norm bound on
\(\widetilde\varphi_\varepsilon\). Therefore
\[
|\widetilde U_\varepsilon \triangle D_1|
=
\frac12\int_{\mathbb S^1}
\left|2\varepsilon \dot\varphi_0 + O(\varepsilon^2)\right|
\,d\mathcal H^1.
\]
Set \(\alpha:=\int_{\mathbb S^1}|\dot\varphi_0|\,d\mathcal H^1\). By
\eqref{eq:l1-nontrivial-new}, \(\alpha>0\). The reverse triangle inequality in
\(L^1\) and the \(O(\varepsilon^2)\) remainder give
\[
\int_{\mathbb S^1}\bigl|2\varepsilon\dot\varphi_0+O(\varepsilon^2)\bigr|\,d\mathcal H^1
\geq
2|\varepsilon|\alpha-O(\varepsilon^2)
\]
and similarly the upper bound
\[
\int_{\mathbb S^1}\bigl|2\varepsilon\dot\varphi_0+O(\varepsilon^2)\bigr|\,d\mathcal H^1
\leq
2|\varepsilon|\alpha+O(\varepsilon^2).
\]
Consequently there exist constants $c,C>0$ such that
\begin{equation}\label{eq:distance-centered-disk-new}
c|\varepsilon|
\leq
|\widetilde U_\varepsilon \triangle D_1|
\leq
C|\varepsilon|
\end{equation}
for all sufficiently small $\varepsilon$.

\medskip
\noindent
{\bf Step 5: comparison with arbitrary disks.} Let $D$ be any disk of area $\pi$. Then $D=D_1+a$ for some $a\in \C$. We first consider disks whose centers are not too close to the origin.
Since
\[
|D_1\triangle (D_1+a)|
\gtrsim \min\{|a|,1\},
\]
the triangle inequality and \eqref{eq:distance-centered-disk-new} imply that
there exists $K\gg 1$ such that, if $|a|\geq K|\varepsilon|$, then
\begin{equation}\label{eq:far-disks-new}
|\widetilde U_\varepsilon \triangle D|
\geq
|D_1\triangle D|-|\widetilde U_\varepsilon\triangle D_1|
\gtrsim
|a|-C|\varepsilon|
\gtrsim
|\varepsilon|.
\end{equation}

It remains to consider disks whose centers satisfy $|a|\leq K|\varepsilon|$.
Write
\[
a=\eta \varepsilon e^{i\theta},
\qquad
|\eta|\leq K,
\quad
\theta\in [0,2\pi).
\]
We claim that, for $\varepsilon$ sufficiently small, the translated domain
\[
\widetilde U_\varepsilon-a
\]
is still a normal graph over $\mathbb S^1$:
\[
\partial(\widetilde U_\varepsilon-a)
=
\{(1+\psi_\varepsilon^{\theta,\eta}(y))y:\ y\in \mathbb S^1\},
\]
with
\[
\psi_\varepsilon^{\theta,\eta}\in C^1(\mathbb S^1),
\qquad
\|\psi_\varepsilon^{\theta,\eta}\|_{C^1}\lesssim |\varepsilon|,
\]
uniformly in $(\theta,\eta)$ with $|\eta|\leq K$. This follows because
$\widetilde U_\varepsilon$ is already an $O(\varepsilon)$-$C^1$ perturbation of
$D_1$, and the translation vector $a$ is also of order $\varepsilon$. We now compute the first-order term of $\psi_\varepsilon^{\theta,\eta}$.
A translation by $a$ contributes, to first order in $\varepsilon$, the normal
displacement
\[
-\eta \Re(e^{-i\theta}y).
\]
Indeed, for \(y\in \mathbb S^1\),
\[
|y-a|
=
\bigl(1-2\Re(\overline y a)+|a|^2\bigr)^{1/2}
=
1-\Re(\overline y a)+O(|a|^2),
\]
uniformly in \(y\), and here \(a=\eta\varepsilon e^{i\theta}\).
Therefore
\begin{equation}\label{eq:psi-expansion-new}
\psi_\varepsilon^{\theta,\eta}(y)
=
\varepsilon\Big(\dot\varphi_0(y)-\eta \Re(e^{-i\theta}y)\Big)
+
O(\varepsilon^2)
\quad \text{in } C^1(\mathbb S^1),
\end{equation}
uniformly for $|\eta|\leq K$.

Define
\[
g^{\theta,\eta}(y)
:=
\dot\varphi_0(y)-\eta \Re(e^{-i\theta}y).
\]
Using \eqref{eq:mode-decomposition-new}, we see that $g^{\theta,\eta}$ belongs
to the finite-dimensional space spanned by the first and second harmonics.
More precisely,
\[
g^{\theta,\eta}(e^{it})
=
a\cos(2t)+b\sin(2t)-\eta \cos(t-\theta).
\]
On the Fourier side this is the key cancellation point: the translation
correction $-\eta\Re(e^{-i\theta}y)$ contributes only to the first harmonics
(its Fourier modes are $\pm 1$), while the perturbation
$\dot\varphi_0$ lives entirely in the second harmonics ($\pm 2$). Hence the
second-harmonic component is invariant under the translation correction, and
in particular the second-harmonic part of $g^{\theta,\eta}$ is independent of
$(\theta,\eta)$ and nonzero, because $(a,b)\neq (0,0)$. To extract a uniform \(L^1\) lower bound, note that
\(g^{\theta,\eta}\) lies in the \(4\)-dimensional space
\[
V:=\mathrm{span}\bigl\{\cos t,\sin t,\cos 2t,\sin 2t\bigr\}\subset L^1(\mathbb S^1).
\]
On this finite-dimensional space, all norms are equivalent; in particular, the
\(L^1(\mathbb S^1)\) and \(L^\infty(\mathbb S^1)\) norms are equivalent, and
convergence in coefficients is equivalent to convergence in any of these norms. The map
\[
(\theta,\eta)\in [0,2\pi)\times [-K,K]\longmapsto g^{\theta,\eta}\in V
\]
is continuous, and its image
\[
\mathcal G_K:=\{g^{\theta,\eta}:\ \theta\in[0,2\pi),\ |\eta|\leq K\}
\]
is the continuous image of a compact set, hence compact in \(V\). Moreover, the
\((\cos 2t,\sin 2t)\)-coefficients of every \(g^{\theta,\eta}\) are exactly
\((a,b)\neq (0,0)\), independent of \((\theta,\eta)\); thus \(0\notin
\mathcal G_K\). The function \(g\mapsto\|g\|_{L^1(\mathbb S^1)}\) is continuous
and strictly positive on \(V\setminus\{0\}\), so it attains a positive minimum
on the compact set \(\mathcal G_K\). It follows that there exists
$\delta_K>0$ such that
\begin{equation}\label{eq:uniform-l1-lower-new}
\int_{\mathbb S^1}|g^{\theta,\eta}(y)|\,d\mathcal H^1(y)\geq \delta_K
\qquad
\text{for all }\theta\in[0,2\pi),\ |\eta|\leq K.
\end{equation}

Now,
\[
|\widetilde U_\varepsilon \triangle D|
=
|(\widetilde U_\varepsilon-a)\triangle D_1|.
\]
Since $\widetilde U_\varepsilon-a$ is star-shaped with radial graph
$1+\psi_\varepsilon^{\theta,\eta}$, the same polar-coordinate computation as in
Step~4 gives
\[
|(\widetilde U_\varepsilon-a)\triangle D_1|
=
\frac12\int_{\mathbb S^1}
\left|(1+\psi_\varepsilon^{\theta,\eta}(y))^2-1\right|
\,d\mathcal H^1(y).
\]
Using \eqref{eq:psi-expansion-new},
\[
(1+\psi_\varepsilon^{\theta,\eta})^2-1
=
2\varepsilon g^{\theta,\eta}+O(\varepsilon^2)
\quad \text{in } L^\infty(\mathbb S^1),
\]
uniformly for $|\eta|\leq K$. Thus, by \eqref{eq:uniform-l1-lower-new},
there exists $c_K>0$ such that
\begin{equation}\label{eq:near-disks-new}
|\widetilde U_\varepsilon \triangle D|
=
|(\widetilde U_\varepsilon-a)\triangle D_1|
\geq
c_K |\varepsilon|
\end{equation}
whenever $|a|\leq K|\varepsilon|$. Combining \eqref{eq:far-disks-new} and \eqref{eq:near-disks-new}, we conclude
that
\[
\inf_{\substack{D \text{ disk}\\ |D|=\pi}}
|\widetilde U_\varepsilon \triangle D|
\gtrsim |\varepsilon|.
\]
Since \(r_\varepsilon=1+O(\varepsilon^2)\), the dilation from \(U_\varepsilon\)
to \(\widetilde U_\varepsilon\) changes the set only by
\[
|\widetilde U_\varepsilon\triangle U_\varepsilon|=O(\varepsilon^2),
\]
and therefore does not affect the first-order distance estimate. Thus the
area-normalized family \(\widetilde U_\varepsilon\) remains at
symmetric-difference distance \(\gtrsim |\varepsilon|\) from the class of
disks.

\medskip
\noindent
{\bf Step 6: square-root sharpness of the deficit estimate.} It remains to translate the distance estimate into the corresponding deficit
estimate. Recall that, by construction via Theorem~\ref{thm:perturb}, the
first localization eigenvalue of \(U_\varepsilon\) equals \(\lambda_0\)
\emph{exactly}, since \eqref{eq:functiona-P-eps-new} states that
\(P_\varepsilon\) is an eigenfunction with eigenvalue \(\lambda_0\) for the
localization operator on \(U_\varepsilon\). Indeed, since
\(P_\varepsilon\to 1\) in \(C^\infty\) on compact sets and \(U_\varepsilon\to D_1\)
in $C^2$, the principal eigenvalue \(\lambda_1(U_\varepsilon)\) and the
corresponding $L^2$-normalized eigenfunction depend continuously on
\(\varepsilon\) by standard analytic perturbation theory; for \(\varepsilon=0\),
the principal eigenpair is precisely \((1-e^{-\pi},1)\) (where the constant
function corresponds to the Hermite ground state $h_0$ in the Bargmann--Fock
realization, cf.~\cite{Daubechies}). Hence for small \(|\varepsilon|\) the
eigenpair \((\lambda_0,P_\varepsilon)\) provided by \eqref{eq:functiona-P-eps-new}
is the principal one. By
\eqref{eq:measure-close-u-eps-new},
\(|U_\varepsilon|=\pi+O(\varepsilon^2)\), and rescaling by
\(r_\varepsilon=1+O(\varepsilon^2)\) to area \(\pi\) changes the principal
eigenvalue at most by \(O(\varepsilon^2)\): since the localization operator
on a dilated set $rU$ depends analytically on $r$, with derivative at $r=1$
proportional to a boundary functional that vanishes on the disk by radial
symmetry, the disk is a critical point of the dilation map, so the
eigenvalue varies quadratically in \(r-1=O(\varepsilon^2)\). It follows that the eigenvalue deficit of
\(\widetilde U_\varepsilon\) with respect to the disk satisfies
\[
\lambda_0-\lambda(\widetilde U_\varepsilon)=O(\varepsilon^2),
\]
where \(\lambda(\widetilde U_\varepsilon)\) denotes the first localization
eigenvalue of \(\widetilde U_\varepsilon\). On the other hand, by Step~5,
the symmetric-difference asymmetry satisfies
\[
\inf_{\substack{D\text{ disk}\\ |D|=\pi}}|\widetilde U_\varepsilon\triangle D|
\gtrsim |\varepsilon|.
\]
Comparing the two,
\[
\inf_{\substack{D\text{ disk}\\ |D|=\pi}}|\widetilde U_\varepsilon\triangle D|
\gtrsim |\varepsilon|
\simeq
\sqrt{\lambda_0-\lambda(\widetilde U_\varepsilon)}.
\]
This shows that the inequality
\[
\inf_{\substack{D\text{ disk}\\ |D|=\pi}}
|\Omega\triangle D|
\lesssim
\sqrt{\lambda_0-\lambda(\Omega)}
\]
of Corollary~1.4 in \cite{Gomez-Guerra-Ramos-Tilli} is square-root-tight along
the family \(\widetilde U_\varepsilon\): the deficit is of order
\(\varepsilon^2\), the asymmetry is of order \(\varepsilon\), and so the
square root of the deficit and the asymmetry are of comparable size. This
exhibits the optimality of the exponent \(1/2\) and concludes the proof of the
sharpness of the estimate in Corollary~1.4 of
\cite{Gomez-Guerra-Ramos-Tilli}, complementing the rigidity result of
\cite{Nicola-Tilli}.

\end{proof}

% ==============================================================
% Section 6
% ==============================================================

\section{Local maximizers for the Faber--Krahn problem in the Gaussian STFT setting}\label{sec:local-max}

The purpose of this section is to establish a local rigidity statement for the
Faber--Krahn problem associated with the Gaussian STFT. The global problem,
solved in \cite{Nicola-Tilli}, asserts that, among measurable subsets of \(\C\)
of fixed measure, the centered disk uniquely maximizes the first eigenvalue of
the Gaussian STFT localization operator. Here we show that this rigidity is
already visible at the \emph{local} level: any set whose first eigenvalue cannot
be increased by small symmetric-difference perturbations must, up to a
translation, be a disk. This complements the quantitative stability result of
\cite{Gomez-Guerra-Ramos-Tilli} and provides the variational counterpart of the
inverse-spectral analysis carried out in the previous sections.

We work throughout in the Bargmann--Fock representation; we refer the reader to
\cite{Bargmann,Folland,Grochenig} for the relevant background. Thus, if
\(F\in \mathcal F^2(\C)\) is the Bargmann transform of \(f\in L^2(\R)\), then
\[
|V_\varphi f(z)|^2 = |F(z)|^2 e^{-\pi |z|^2},
\]
up to the normalization already fixed in the previous sections. Accordingly, for
\(F\in \mathcal F^2(\C)\) we set
\[
u_F(z):=|F(z)|^2 e^{-\pi |z|^2}.
\]

For a measurable set \(\Omega\subset \C\), let
\[
\lambda_{1,\varphi}(\Omega)
:=
\sup_{\|f\|_2=1}\int_\Omega |V_\varphi f(z)|^2\,dz.
\]
We say that \(\Omega\) is a \emph{local maximizer} for \(\lambda_{1,\varphi}\) under the
measure constraint if there exists \(\delta_0>0\) such that, whenever
\(E\subset \C\) is measurable, \(|E|=|\Omega|\), and
\[
\int_\C |\mathbf 1_\Omega-\mathbf 1_E|\,e^{-\pi |z|^2}\,dz<\delta_0,
\]
one has
\[
\lambda_{1,\varphi}(\Omega)\ge \lambda_{1,\varphi}(E).
\]

We now prove the local rigidity statement announced in the introduction. The
proof proceeds in six steps. Steps 1--3 are essentially the analogues, in our
setting, of standard arguments in the Faber--Krahn theory: \(\Sigma\) is forced
to be a level set of \(u_F\), the corresponding eigenvalue is simple, and a
first-variation identity ties \(F\) to a divergence-free vector field built from
holomorphic data. The crucial new ingredient lies in Step 4, where we show that
the \emph{derivative} \(F'\) is again a first eigenfunction. Combined with the
simplicity from Step 2, this forces \(F\) to satisfy a first-order linear ODE,
and hence to be a Gaussian. Once this is in place, Steps 5--6 use the Weyl
invariance of the Bargmann--Fock setting and a one-dimensional rearrangement
argument to identify \(\Sigma\) with a centered disk.

\begin{proof}[Proof of the local-maximizer theorem]
Let \(\Sigma\subset \C\) be a local maximizer for \(\lambda_{1,\varphi}\) under
the measure constraint, and let \(F\in \mathcal F^2(\C)\) be a normalized first
eigenfunction for the localization operator associated with \(\Sigma\).
Equivalently,
\begin{equation}\label{eq:eigenvalue-equation}
\int_\Sigma H(z)\overline{F(z)}\,d\mu(z)
=
\lambda_{1,\varphi}(\Sigma)\int_\C H(z)\overline{F(z)}\,d\mu(z)
\end{equation}
for every \(H\in \mathcal F^2(\C)\), where
\[
d\mu(z):=e^{-\pi |z|^2}\,dA(z).
\]

\medskip
\noindent
{\bf Step 1: \(\Sigma\) is a superlevel set of \(u_F\).}

We claim that there exists \(t_0\ge 0\) such that, up to sets of measure zero,
\begin{equation}\label{eq:superlevel-claim}
\Sigma=\{u_F>t_0\}.
\end{equation}

The strategy is a bathtub-style perturbation argument: if \(\Sigma\) were not a
superlevel set of \(u_F\), one could move a small piece of \(\Sigma\) where
\(u_F\) is small to a region where \(u_F\) is larger, increasing the
localization integral while preserving Lebesgue measure and remaining within
the local-maximizer neighborhood of \(\Sigma\).

Recall that \(F\) is admissible in the variational problem defining
\(\lambda_{1,\varphi}(\Sigma')\) for any measurable set \(\Sigma'\) of finite
measure. In particular,
\[
\lambda_{1,\varphi}(\Sigma')\ge \int_{\Sigma'} u_F\,dA,
\qquad
\lambda_{1,\varphi}(\Sigma)= \int_\Sigma u_F\,dA,
\]
the latter equality holding because \(F\) is the actual eigenfunction for
\(\Sigma\). It therefore suffices to find a competitor \(\Sigma'\) with
\(|\Sigma'|=|\Sigma|\), close to \(\Sigma\) in the sense of weighted symmetric
difference, and with \(\int_{\Sigma'} u_F\,dA>\int_\Sigma u_F\,dA\); that would
violate local maximality.

Suppose, for the sake of contradiction, that \eqref{eq:superlevel-claim} fails.
The function \(u_F\) is real-analytic on \(\C\setminus\{F=0\}\), so for almost
every \(t\ge 0\) the level set \(\{u_F=t\}\) has Lebesgue measure zero;
consequently the distribution function \(t\mapsto |\{u_F>t\}|\) is continuous
and strictly decreasing on its support. Hence there exists a unique
\(t_0\ge 0\) such that
\[
A:=\{u_F>t_0\}\quad\text{satisfies}\quad |A|=|\Sigma|;
\]
since \eqref{eq:superlevel-claim} fails by assumption, \(A\) and \(\Sigma\)
do not coincide up to a null set, so both \(A\setminus\Sigma\) and
\(\Sigma\setminus A\) have positive Lebesgue measure (and in fact positive
weighted measure, since \(e^{-\pi|z|^2}>0\) everywhere).

We now perform the bathtub construction explicitly. The set
\(A=\{u_F>t_0\}\) is open since \(u_F\) is continuous, while
\(\Sigma\setminus A\) is only known to be measurable. By inner regularity of
the Lebesgue measure, choose a compact set \(L\subset A\setminus\Sigma\) with
\(|L|>0\) and a compact set \(K\subset \Sigma\setminus A\) with \(|K|>0\)
(both nonempty by the previous paragraph). Note that \(L\subset A\) is then
contained in the open set \(A\), and we can shrink \(L\) further if needed so
that \(\int_L d\mu, \int_K d\mu<\delta_0/4\). Since the Lebesgue measure
restricted to a compact set has the intermediate-value property, we may
choose subsets \(S\subset L\) and \(T\subset K\) with the common Lebesgue
measure
\[
|S|=|T|=\tfrac12\min(|L|,|K|)>0.
\]
Then \(S\subset A\setminus\Sigma\), \(T\subset \Sigma\setminus A\), and
\[
\int_\C (\mathbf 1_S+\mathbf 1_T)\,e^{-\pi |z|^2}\,dz<\tfrac{\delta_0}{2}<\delta_0.
\]
Define
\[
\Sigma':=(\Sigma\setminus T)\cup S.
\]
Then \(|\Sigma'|=|\Sigma|\) and
\[
\int_\C |\mathbf 1_\Sigma-\mathbf 1_{\Sigma'}|e^{-\pi |z|^2}\,dz
=\int_\C(\mathbf 1_S+\mathbf 1_T)e^{-\pi|z|^2}\,dz<\delta_0.
\]
On the other hand, by construction,
\[
u_F(z)>t_0 \quad \text{on } S\subset A,
\qquad
u_F(z)\le t_0 \quad \text{on } T\subset \Sigma\setminus A.
\]
Moreover, since \(u_F\) is continuous and \(L\subset A=\{u_F>t_0\}\) is
compact, the continuous function \(u_F\) attains its minimum on \(L\); hence
there exists \(\eta>0\) with \(u_F\ge t_0+\eta\) on \(L\), and a fortiori on
\(S\). It then follows that
\[
\int_S u_F\,dA \ge (t_0+\eta)\,|S|,
\qquad
\int_T u_F\,dA \le t_0\,|T|=t_0\,|S|,
\]
and hence
\[
\int_{\Sigma'} u_F\,dA
=
\int_\Sigma u_F\,dA + \int_S u_F\,dA-\int_T u_F\,dA
\ge \int_\Sigma u_F\,dA + \eta\,|S|
>
\int_\Sigma u_F\,dA.
\]
Since \(F\) is admissible for the variational problem defining
\(\lambda_{1,\varphi}(\Sigma')\),
\[
\lambda_{1,\varphi}(\Sigma')
\ge \int_{\Sigma'} u_F\,dA
>
\int_\Sigma u_F\,dA
=
\lambda_{1,\varphi}(\Sigma),
\]
contradicting the local maximality of \(\Sigma\). This proves
\eqref{eq:superlevel-claim}.

In particular, since \(u_F\) is real-analytic outside the (discrete) zero set of
\(F\), \(u_F\) is constant on \(\partial \Sigma\) in the sense that there exists
\(c_\Sigma\ge 0\) with
\begin{equation}\label{eq:boundary-level-set-new}
u_F(z)=c_\Sigma
\qquad \text{for } z\in \partial \Sigma.
\end{equation}
We will see in Step 2 that \(F\) does not vanish on \(\overline\Sigma\), and in
particular \(c_\Sigma>0\).

\medskip
\noindent
{\bf Step 2: simplicity of the first eigenvalue.}

We claim that the first eigenvalue \(\lambda_{1,\varphi}(\Sigma)\) is simple.
The argument uses a modular-deformation principle: any two first eigenfunctions
have the same modulus on \(\partial\Sigma\) up to a constant, and a holomorphic
function with constant modulus on a connected open set must itself be constant.
This kind of argument plays a role also in the stability analysis of
\cite{Gomez-Guerra-Ramos-Tilli}.

Suppose, for contradiction, that \(G\in \mathcal F^2(\C)\) is another first
eigenfunction, linearly independent of \(F\). Step 1 applied to \(G\) shows that
\(\Sigma\) is also a superlevel set of \(u_G\); hence there exists
\(c'_\Sigma\ge 0\) such that
\[
u_G(z)=c'_\Sigma
\qquad\text{for } z\in \partial \Sigma.
\]
Equivalently,
\[
|F(z)|^2 e^{-\pi |z|^2}=c_\Sigma,
\qquad
|G(z)|^2 e^{-\pi |z|^2}=c'_\Sigma,
\qquad z\in \partial \Sigma.
\]

We next observe that first eigenfunctions are zero-free in the interior of
\(\Sigma\). Indeed, by Step 1 we have \(u_F>c_\Sigma\) on the (open) interior of
\(\Sigma\). Since
\[
u_F=|F|^2 e^{-\pi |z|^2}
\]
and the Gaussian factor \(e^{-\pi|z|^2}\) is strictly positive everywhere, this
yields \(|F|>0\) throughout the interior of \(\Sigma\); the same applies to
\(|G|\). In particular, \(c_\Sigma>0\) and \(c'_\Sigma>0\). Moreover, since
\(F,G\) are entire and continuous, \(|F|>0\) on a neighborhood of any compact
subset of \(\Sigma^\circ\); by Step 1, we may take a slight enlargement of
\(\Sigma\) (still a superlevel set of \(u_F\)) on which \(|F|,|G|>0\), so we may
work on a connected open set \(U\Supset \Sigma\) where both \(F\) and \(G\) are
zero-free.

The ratio \(F/G\) is holomorphic and nowhere zero on \(U\), and the function
\[
\log \left|\frac{F}{G}\right|
=
\frac12 \log \frac{|F|^2}{|G|^2}
=
\frac12 \log \frac{u_F}{u_G}
\]
is harmonic on \(U\): it is the real part of any local branch of
\(\log(F/G)\), which is well defined on simply connected subsets, and harmonic
in any case as the logarithm of the modulus of a zero-free holomorphic
function. Note that the Gaussian factors cancel inside the logarithm, so on
\(\partial \Sigma\) by \eqref{eq:boundary-level-set-new} this function equals
the constant \(\frac12\log(c_\Sigma/c'_\Sigma)\). By the maximum principle
applied separately on each connected component of \(\Sigma\), the function
\(\log|F/G|\) is constant on each component. Hence
\[
\left|\frac{F}{G}\right| \equiv \text{constant}
\qquad \text{on each connected component of } \Sigma.
\]
A nonconstant holomorphic function is open, so a holomorphic function with
constant modulus on a connected open set must itself be constant. Therefore
\[
F=\gamma_j G
\qquad \text{on each connected component } \Sigma_j\text{ of } \Sigma.
\]
The relation \(F=\gamma_j G\) is an identity between two entire functions on a
set with an accumulation point; the analytic continuation principle gives
\(F=\gamma_j G\) on all of \(\C\) for any single \(j\). In particular all
\(\gamma_j\) coincide, and \(F\) and \(G\) are linearly dependent, contradicting
our assumption. This proves simplicity.

\medskip
\noindent
{\bf Step 3: a first-variation identity.}

By Step 2, the eigenfunction \(F\) is determined up to a unimodular constant.
We now derive a first-variation identity satisfied by \(F\), based on the fact
that \(\Sigma\) is a level set of \(u_F\) and that holomorphic vector fields are
divergence-free.

The function \(u_F=|F|^2e^{-\pi|z|^2}\) is real-analytic on
\(\C\setminus\{F=0\}\), since \(F\) is entire and the exponential factor is
real-analytic. After replacing \(\Sigma\) by its open superlevel-set
representative \(\{u_F>c_\Sigma\}\) from Step 1 (which is permissible since
local maximality only sees the equivalence class of \(\Sigma\) up to null
sets), we have \(u_F\ge c_\Sigma>0\) on \(\overline\Sigma\); since the
Gaussian factor never vanishes, this forces \(F\ne 0\) on
\(\overline\Sigma\), so \(u_F\) is real-analytic on a neighborhood of
\(\overline\Sigma\). The boundary \(\partial\Sigma\) is therefore a level set
of a non-constant real-analytic function and so is the union of finitely many
real-analytic arcs, plus a possible finite set of singular points where
\(\nabla u_F\) vanishes (see, e.g., the standard local structure of
real-analytic varieties). In particular, \(\partial\Sigma\) is Lipschitz, and
\(C^1\) away from the critical set of \(u_F\); the divergence theorem applies
to \(\Sigma\) for any \(C^1\)-vector field on a neighborhood of
\(\overline\Sigma\), the singular boundary points forming a set of
one-dimensional Hausdorff measure zero.

Let \(\psi\) be holomorphic on an open neighborhood \(V\) of
\(\overline \Sigma\), and write
\[
\psi=u+iv
\]
with \(u,v\) real-valued. Consider the real vector field
\[
X_\psi:=(u,-v).
\]
Since \(\psi\) is holomorphic, the Cauchy--Riemann equations
\[
u_x=v_y,
\qquad
u_y=-v_x,
\]
imply
\[
\operatorname{div}X_\psi=u_x+(-v)_y=u_x-v_y=0
\qquad \text{in }V.
\]
Because \(u_F\equiv c_\Sigma\) on \(\partial\Sigma\), the divergence theorem on
\(\Sigma\) gives
\[
\int_\Sigma \langle \nabla u_F, X_\psi\rangle\,dA
=
\int_\Sigma \operatorname{div}(u_F X_\psi)\,dA
=
\int_{\partial \Sigma} u_F\, X_\psi\cdot \nu\, d\mathcal H^1
=
c_\Sigma\int_{\partial \Sigma} X_\psi\cdot \nu\, d\mathcal H^1
=
0,
\]
where the last equality uses
\(\int_{\partial\Sigma}X_\psi\cdot\nu\,d\mathcal H^1
=\int_\Sigma\operatorname{div}X_\psi\,dA=0\) by Cauchy--Riemann. Since
\(u_F=|F|^2e^{-\pi |z|^2}\) and
\(\nabla(e^{-\pi|z|^2})=-2\pi e^{-\pi|z|^2}(x,y)\), this becomes
\begin{equation}\label{eq:vector-field-identity-real-new}
0=
\int_\Sigma
\left(
\langle \nabla |F|^2,X_\psi\rangle
-
2\pi |F|^2 \langle z,X_\psi\rangle
\right)
\,d\mu.
\end{equation}

We now express \eqref{eq:vector-field-identity-real-new} in complex notation.
A direct computation, using \(\partial_z|F|^2=F'\overline F\) and
\(\partial_{\bar z}|F|^2=F\overline{F'}\), gives
\[
\partial_x |F|^2 = F'\overline F+\overline{F'}F = 2\Re(F'\overline F),
\qquad
\partial_y |F|^2 = i(F'\overline F-\overline{F'}F) = -2\Im(F'\overline F).
\]
Writing \(\psi=u+iv\) and recalling \(X_\psi=(u,-v)\),
\[
\langle \nabla |F|^2,X_\psi\rangle
=
2u\,\Re(F'\overline F)+2v\,\Im(F'\overline F)
=
2\,\Re\!\big(\overline{\psi}\,F'\overline F\big)
=
2\,\Re\!\big(\psi\,\overline{F'}F\big),
\]
where in the last equality we used \(\Re w=\Re\overline w\). Likewise, writing
\(z=x+iy\),
\[
\langle z,X_\psi\rangle
=
xu+y(-v)
=
xu-yv
=
\Re\big((x+iy)(u+iv)\big)
=
\Re(z\psi).
\]
Substituting these expressions into
\eqref{eq:vector-field-identity-real-new}, we obtain
\[
0=
\Re\!\int_\Sigma
\psi(z)\Big(\overline{F'(z)}\,F(z)-\pi z\,|F(z)|^2\Big)\,d\mu(z).
\]
Replacing \(\psi\) by \(i\psi\) (which is also holomorphic on \(V\)) gives the
analogous identity for the imaginary part:
\[
0=
\Im\!\int_\Sigma
\psi(z)\Big(\overline{F'(z)}\,F(z)-\pi z\,|F(z)|^2\Big)\,d\mu(z).
\]
Combining the two, and using \(|F|^2=F\overline F\) to factor out \(F(z)\)
from both terms inside the integrand (the \(F\) from \(\overline{F'}F\) and
one copy of \(F\) from \(|F|^2=F\overline F\)),
\begin{equation}\label{eq:orthogonality-main-new}
0=
\int_\Sigma
\psi(z)\,F(z)\Big(\overline{F'(z)}-\pi z\,\overline{F(z)}\Big)\,d\mu(z),
\end{equation}
which holds for every holomorphic \(\psi\) on a neighborhood of
\(\overline \Sigma\). Equivalently, expanding the product,
\begin{equation*}
0=
\int_\Sigma
\psi(z)\Big(\overline{F'(z)}\,F(z)-\pi z\,|F(z)|^2\Big)\,d\mu(z).
\end{equation*}
The form \eqref{eq:orthogonality-main-new} is the one used in Step 4.

\medskip
\noindent
{\bf Step 4: \(F'\) is again a first eigenfunction.}

This step is the crux of the argument and is the main new ingredient compared
to classical Faber--Krahn rigidity proofs. Rather than working with general
test functions in the eigenvalue equation, we exploit the divergence identity
\eqref{eq:orthogonality-main-new} together with a clever choice of test
function involving the reproducing kernel of the Fock space. The end result is
that \emph{the derivative} \(F'\) of any first eigenfunction is itself a first
eigenfunction; since the eigenspace is one-dimensional by Step 2, this forces
a first-order ODE on \(F\), and hence pins down \(F\) up to translation as a
Gaussian.

For every \(w\in \C\), set
\[
K_w(z):=e^{\pi \overline w z};
\]
this is the reproducing kernel of \(\mathcal F^2(\C)\) at the point \(w\), in
the sense that
\begin{equation}\label{eq:reproducing-kernel}
\int_\C H(z)\,e^{\pi w \overline z}\,d\mu(z)=H(w),
\qquad H\in\mathcal F^2(\C);
\end{equation}
see, e.g., \cite[Ch.~1]{Folland} or \cite{Bargmann}. Equivalently,
\(\langle H,K_w\rangle_{\mathcal F^2}=H(w)\).

By Step 1, after replacing \(\Sigma\) by the corresponding superlevel-set
representative if necessary, we may suppose that
\[
\Sigma=\{u_F>c_\Sigma\}.
\]
In particular, \(u_F\ge c_\Sigma>0\) on \(\overline \Sigma\). Since
\(u_F(z)=|F(z)|^2e^{-\pi |z|^2}\), it follows that \(F\) has no zeros on
\(\overline \Sigma\). By continuity of \(F\) and the compactness of
\(\overline \Sigma\), there exists an open neighborhood
\(V\supset \overline \Sigma\) on which \(F\) has no zeros. Hence, for every
\(w\in \C\), the function
\[
\psi_w(z):=\frac{K_w(z)}{F(z)}=\frac{e^{\pi\overline wz}}{F(z)}
\]
is holomorphic on \(V\). Applying \eqref{eq:orthogonality-main-new} with
\(\psi=\psi_w\), so that
\(\psi_w(z)F(z)=e^{\pi\overline w z}\), we obtain
\[
0=
\int_\Sigma
e^{\pi \overline wz}
\Big(\overline{F'(z)}-\pi z\,\overline{F(z)}\Big)\,d\mu(z),
\]
or, separating the two terms,
\begin{equation}\label{eq:key-kernel-identity-new}
\int_\Sigma
e^{\pi \overline wz}\overline{F'(z)}\,d\mu(z)
=
\int_\Sigma
\pi z\,e^{\pi \overline wz}\overline{F(z)}\,d\mu(z).
\end{equation}

We now relate the right-hand side of \eqref{eq:key-kernel-identity-new} to
\(\overline{F'(w)}\) using the eigenvalue equation
\eqref{eq:eigenvalue-equation} together with the reproducing identity
\eqref{eq:reproducing-kernel}. Choosing
\[
H(z):=z\,e^{\pi\overline w z}\in \mathcal F^2(\C)
\]
(the product of the kernel \(K_w\) by the polynomial \(z\); both factors
preserve membership in \(\mathcal F^2\)) in \eqref{eq:eigenvalue-equation}, we
obtain
\begin{equation}\label{eq:eigvalue-applied}
\int_\Sigma z\, e^{\pi \overline wz}\,\overline{F(z)}\,d\mu(z)
=
\lambda_{1,\varphi}(\Sigma)\int_\C z\, e^{\pi \overline wz}\,\overline{F(z)}\,d\mu(z).
\end{equation}
On the other hand, taking complex conjugates in \eqref{eq:reproducing-kernel}
applied to \(F\) at the point \(w\), one finds
\[
\int_\C e^{\pi \overline wz}\,\overline{F(z)}\,d\mu(z)=\overline{F(w)}.
\]
Both sides depend on \(w\) through the antiholomorphic variable
\(\overline w\); we differentiate with respect to \(\overline w\), exchanging
differentiation and integration on the left. The interchange is legitimate
because, for \(w\) in any compact set \(\{|w|\le R\}\), the integrand
\(\partial_{\overline w}(e^{\pi\overline wz}\overline{F(z)})
=\pi z\,e^{\pi\overline wz}\overline{F(z)}\) is dominated, in absolute
value, by \(\pi |z|\,e^{\pi R|z|}|F(z)|\). The standard pointwise bound for
Fock-space functions gives \(|F(z)|\le\|F\|_{\mathcal F^2}\,e^{\pi |z|^2/2}\),
so this dominator is bounded by a constant times
\(|z|\,e^{\pi R|z|}\,e^{\pi|z|^2/2}\); after multiplication by the Gaussian
weight \(e^{-\pi|z|^2}\) of \(d\mu\), the resulting integrand
\(|z|\,e^{\pi R|z|}\,e^{-\pi |z|^2/2}\) is integrable on \(\C\). Dominated
convergence therefore justifies the interchange. We conclude
\[
\pi\int_\C z\, e^{\pi \overline wz}\,\overline{F(z)}\,d\mu(z)
=
\overline{F'(w)}.
\]
Multiplying \eqref{eq:eigvalue-applied} by \(\pi\) and substituting,
\[
\pi\int_\Sigma z\, e^{\pi \overline wz}\,\overline{F(z)}\,d\mu(z)
=
\lambda_{1,\varphi}(\Sigma)\,\overline{F'(w)}.
\]
Combining this with \eqref{eq:key-kernel-identity-new},
\[
\int_\Sigma e^{\pi \overline wz}\,\overline{F'(z)}\,d\mu(z)
=
\lambda_{1,\varphi}(\Sigma)\,\overline{F'(w)},
\qquad w\in \C.
\]
Taking complex conjugates, this is exactly the eigenvalue equation
\eqref{eq:eigenvalue-equation} applied with \(F\) replaced by \(F'\) (and an
arbitrary \(K_w\) playing the role of the test function). Since the linear
span of \(\{K_w\}_{w\in\C}\) is dense in \(\mathcal F^2(\C)\), the identity
extends to every \(H\in\mathcal F^2(\C)\), and we conclude that \(F'\) is a
first eigenfunction of the localization operator on \(\Sigma\), with eigenvalue
\(\lambda_{1,\varphi}(\Sigma)\).

By the simplicity established in Step 2,
\[
F'=\theta F
\]
for some constant \(\theta\in \C\). Solving this ODE gives
\[
F(z)=C e^{\theta z},
\]
for some nonzero constant \(C\).

\medskip
\noindent
{\bf Step 5: reduction to the constant eigenfunction.}

We now use the Weyl invariance of the Bargmann--Fock setting to reduce the
analysis to the case \(F\equiv 1\). Recall from \eqref{eq:Ta} that, for
\(a\in\C\), the Weyl translation
\[
(T_a F)(z)= e^{\pi a z-\frac{\pi}{2}|a|^2}\,F(z-a)
\]
is unitary on \(\mathcal F^2(\C)\) and satisfies
\(u_{T_aF}(z)=u_F(z-a)\). In particular, \(T_a\) maps eigenfunctions of the
localization operator on \(\Sigma\) to eigenfunctions of the localization
operator on the translated set \(\Sigma+a\), with the same eigenvalue, and the
quantity \(\int_\Sigma u_F\,dA\) is invariant under the simultaneous
substitutions \(F\mapsto T_aF\) and \(\Sigma\mapsto \Sigma+a\).

We claim that we may choose \(a\in\C\) so that \(T_{-a}F\) is a constant
function. Indeed, by Step 4, \(F(z)=Ce^{\theta z}\) for some \(C,\theta\in\C\)
with \(C\ne 0\). Compute, using the identity
\(u_{T_{-a}F}(z)=u_F(z+a)\):
\[
u_{T_{-a}F}(z)=|F(z+a)|^2 e^{-\pi|z+a|^2}
=|C|^2 e^{2\Re(\theta(z+a))}\,e^{-\pi|z+a|^2}.
\]
Expanding,
\[
2\Re(\theta(z+a))-\pi|z+a|^2
=
2\Re(\theta z)+2\Re(\theta a)-\pi|z|^2-2\pi\Re(\overline a z)-\pi|a|^2.
\]
The linear terms in \(z\) on the right cancel if and only if
\(2\Re(\theta z)=2\pi\Re(\overline a z)\) for every \(z\in\C\), which is the
condition \(\overline a=\theta/\pi\), i.e.
\[
a:=\frac{\overline\theta}{\pi}.
\]
With this choice,
\[
u_{T_{-a}F}(z)=|C'|^2 e^{-\pi|z|^2},
\qquad
|C'|^2:=|C|^2 e^{2\Re(\theta a)-\pi|a|^2}=|C|^2 e^{|\theta|^2/\pi}.
\]
Since \(T_{-a}F\) is entire and \(|T_{-a}F(z)|^2 e^{-\pi|z|^2}=|C'|^2
e^{-\pi|z|^2}\), it follows that \(|T_{-a}F|\equiv|C'|\) is constant; a
holomorphic function with constant modulus is constant, hence
\(T_{-a}F\equiv C''\) for some \(C''\in\C\) with \(|C''|=|C'|\). Normalizing,
we may absorb the unimodular constant into the eigenfunction and assume
\[
F\equiv 1
\]
after the change \(\Sigma\rightsquigarrow \Sigma-a\) (and the corresponding
unitary substitution \(F\rightsquigarrow T_{-a}F\)).

In this normalization the ground state \(h_0\) (corresponding to
\(F\equiv 1\) in Bargmann--Fock coordinates) is an eigenfunction of the
localization operator of \(\Sigma\). By Step 1 (applied with the new
normalization) we have, up to a null set,
\[
\Sigma=\{z\in\C:\ |C'|^2 e^{-\pi |z|^2}>c_\Sigma\}.
\]
Since \(r\mapsto e^{-\pi r^2}\) is \emph{strictly} decreasing, this
superlevel set is automatically a centered open disk; it cannot be an
annulus. Thus \(\Sigma\) is, up to a null set, the centered disk \(D_R\)
with \(|D_R|=|\Sigma|\), and Step 6 below records the corresponding
quantitative rearrangement statement, which is an instance of the global
Faber--Krahn theorem of \cite{Nicola-Tilli}.

\medskip
\noindent
{\bf Step 6: exclusion of annuli and conclusion.}

It remains to verify that, among radial sets of prescribed measure, the only
local maximizer is the centered disk. Although this is essentially automatic
after Step 5, we record the standard rearrangement argument since it makes
explicit \emph{why} no annular local maximizer can exist; this provides the
endpoint matching the global Faber--Krahn theorem of \cite{Nicola-Tilli}.

Suppose, more generally, that \(\Sigma\) is any radial set of finite measure
which arises as a superlevel set of a radial real-analytic non-constant
function on \(\C\); write
\[
\widetilde \Sigma:=\{r>0:\ re^{it}\in \Sigma \text{ for some } t\in [0,2\pi)\}.
\]
By Step 1, \(\Sigma\) admits an open representative; combined with the
real-analytic structure established in Step 3, the boundary
\(\partial\Sigma\) is a finite union of real-analytic arcs, hence
\(\widetilde \Sigma\subset (0,\infty)\) is open and is a finite union of
disjoint open intervals:
\[
\widetilde \Sigma=\bigcup_{j=1}^N (a_j,b_j),
\qquad
0\le a_1<b_1\le a_2<b_2\le \cdots\le a_N<b_N.
\]
Then, using \(\int re^{-\pi r^2}\,dr=-\frac{1}{2\pi}e^{-\pi r^2}\),
\[
\int_\Sigma e^{-\pi |z|^2}\,dA(z)
=
2\pi \int_{\widetilde \Sigma} r e^{-\pi r^2}\,dr
=
\sum_{j=1}^N \big(e^{-\pi a_j^2}-e^{-\pi b_j^2}\big),
\]
under the constraint
\[
|\Sigma|
=
2\pi \int_{\widetilde \Sigma} r\,dr
=
\pi \sum_{j=1}^N (b_j^2-a_j^2).
\]
We are therefore led to the optimization problem
\begin{equation}\label{eq:radial-opt}
\text{maximize}\qquad
\sum_{j=1}^N \big(e^{-\pi a_j^2}-e^{-\pi b_j^2}\big)
\qquad \text{subject to}\qquad
\sum_{j=1}^N (b_j^2-a_j^2)=c,
\end{equation}
where \(c=|\Sigma|/\pi\) is fixed. Substituting \(s_j:=a_j^2\),
\(t_j:=b_j^2\) (with \(0\le s_1<t_1\le s_2<\cdots\)), and writing
\(\Phi(x):=e^{-\pi x}\), \eqref{eq:radial-opt} becomes
\[
\text{maximize}\qquad
\sum_{j=1}^N \big(\Phi(s_j)-\Phi(t_j)\big)
\qquad \text{subject to}\qquad
\sum_{j=1}^N (t_j-s_j)=c.
\]
The function \(\Phi\) is strictly convex and strictly decreasing on
\([0,\infty)\); its derivative \(-\Phi'(\sigma)=\pi e^{-\pi\sigma}\) is strictly
decreasing in \(\sigma\). For any \(t>s\ge 0\),
\[
\Phi(s)-\Phi(t)
=
\int_s^t (-\Phi'(\sigma))\,d\sigma,
\]
so
\[
\sum_{j=1}^N\big(\Phi(s_j)-\Phi(t_j)\big)
=
\int_{\widetilde\Sigma_2}(-\Phi'(\sigma))\,d\sigma,
\qquad
\widetilde\Sigma_2:=\bigcup_{j=1}^N (s_j,t_j).
\]
The set \(\widetilde\Sigma_2\subset[0,\infty)\) has Lebesgue measure equal to
\(c\), and the integrand \(-\Phi'\) is strictly decreasing in \(\sigma\). The
bathtub principle (or, equivalently, a direct rearrangement argument) shows
that, among all measurable subsets of \([0,\infty)\) of total Lebesgue measure
\(c\), the integral \(\int_{\widetilde\Sigma_2}(-\Phi'(\sigma))\,d\sigma\) is
strictly maximized by the leftmost interval \((0,c)\). This corresponds to
\(N=1\), \(s_1=0\), \(t_1=c\), i.e.
\[
\widetilde \Sigma=(0,R),
\qquad R:=\sqrt c,
\]
and equivalently \(\Sigma\) is the centered disk \(D_R\) with
\(|D_R|=|\Sigma|\). Any annular configuration (\(N\ge 2\), or \(N=1\) with
\(s_1>0\)) produces a strictly smaller value of \(\int_\Sigma u_F\,dA\), and is
therefore not a local maximizer.

We conclude that every local maximizer of \(\lambda_{1,\varphi}\) under the
measure constraint is, up to translation, a disk.
\end{proof}


% ==============================================================
% Section 7
% ==============================================================

\section{Global maximizers for the Faber--Krahn problem in the Gaussian STFT setting}\label{sec:global-max}

In this section we prove existence of global maximizers for the Gaussian
time-frequency concentration problem under a measure constraint, and then
combine that existence result with the local rigidity theorem from the previous
section to recover the disk as the unique maximizer, up to translation. The
strategy is by now classical in time-frequency analysis: a profile decomposition
in the spirit of P.-L.~Lions' concentration-compactness, transplanted to the
Fock space via the Bargmann transform, dichotomizes the loss of compactness of a
maximizing sequence into the action of the Weyl translation group. A vanishing
remainder is then ruled out by an iterative extraction whose maximality forces
all the energy to remain captured by the extracted profiles. Closely related
arguments have appeared in the time-frequency setting in
\cite{Nicola-Romero-Trapasso, Marceca-Romero-Speckbacher, Fulsche-Hagger,
Samuelsen}, and the disk-rigidity result we recover is the celebrated theorem
of Nicola--Tilli \cite{Nicola-Tilli, Nicola-Tilli-2}.

Since the Bargmann transform is unitary, the variational problem may be written
on the Fock space side as
\[
M(s):=\sup_{|\Omega|=s}\lambda_{1,\varphi}(\Omega)
=
\sup_{\|F\|_{\mathcal F^2}=1} I_F(s).
\]
Here, and throughout the section, we exploit the translation invariance
\eqref{eq:I-translation-global} together with the pointwise identity
\eqref{eq:u-translation-global} (cf.\ also \eqref{eq:uF-translate}), the
basic regularity properties recorded in Lemma~\ref{lem:basic-properties-global},
and the elementary properties of the Fock space gathered in
Lemma~\ref{lem:Fock-basic}.

We begin with a Weyl-orthogonality lemma that captures the asymptotic
decoupling of two functions whose translation parameters drift apart. This
plays the role of \emph{Pythagorean orthogonality at infinity} on the Fock
space.

\begin{lemma}\label{lem:weyl-orthogonality-global}
Let \(F,G\in \mathcal F^2(\C)\), and let \(\{a_k\}\subset \C\) satisfy
\(|a_k|\to \infty\). Then
\[
T_{a_k}G \rightharpoonup 0
\qquad\text{weakly in }\mathcal F^2(\C),
\]
and in particular
\[
\langle F,T_{a_k}G\rangle_{\mathcal F^2}\to 0.
\]
Consequently,
\[
\|F+T_{a_k}G\|_{\mathcal F^2}^2
=
\|F\|_{\mathcal F^2}^2+\|G\|_{\mathcal F^2}^2+o(1).
\]
More generally, if \(|a_k-b_k|\to\infty\), then
\[
\langle T_{a_k}F,T_{b_k}G\rangle_{\mathcal F^2}\to 0.
\]
\end{lemma}

\begin{proof}
Using unitarity of the Weyl operators \(T_a\) on \(\mathcal F^2(\C)\) (cf.\
\eqref{eq:Ta}),
\[
\langle F,T_{a_k}G\rangle_{\mathcal F^2}
=
\langle T_{-a_k}F,G\rangle_{\mathcal F^2}.
\]
Now \(T_{-a_k}F\) is bounded in \(\mathcal F^2\) by unitarity, and the
identity \eqref{eq:u-translation-global} (equivalently
\eqref{eq:uF-translate}) yields
\[
u_{T_{-a_k}F}(z)=u_F(z+a_k)\qquad (z\in\C),
\]
so that the densities \(u_{T_{-a_k}F}\) drift uniformly to infinity.
Equivalently, since \(u_F\in C_0(\C)\) by Lemma~\ref{lem:Fock-basic}(iii), for
every compact \(K\subset \C\),
\[
\sup_{z\in K}|T_{-a_k}F(z)|e^{-\pi |z|^2/2}
=
\sup_{z\in K}u_{T_{-a_k}F}(z)^{1/2}
=
\sup_{z\in K}u_F(z+a_k)^{1/2}\to 0.
\]
Hence \(T_{-a_k}F\to 0\) locally uniformly. Combining this with the
boundedness of \(\{T_{-a_k}F\}\) in \(\mathcal F^2(\C)\), and using the fact
that the linear span of the reproducing kernels \(\{K_w\}_{w\in\C}\) is dense
in \(\mathcal F^2\), we conclude that \(T_{-a_k}F\rightharpoonup 0\) weakly in
\(\mathcal F^2\); in particular \(\langle T_{-a_k}F,G\rangle\to 0\). This
proves the first claim.

For the norm decoupling, expand the square:
\[
\|F+T_{a_k}G\|_{\mathcal F^2}^2
=
\|F\|_{\mathcal F^2}^2+\|T_{a_k}G\|_{\mathcal F^2}^2
+2\Re\langle F,T_{a_k}G\rangle_{\mathcal F^2}
=
\|F\|_{\mathcal F^2}^2+\|G\|_{\mathcal F^2}^2+o(1),
\]
where we used unitarity of \(T_{a_k}\) and the cross-term decay just proved.

Finally, the general case reduces to the first by writing
\[
\langle T_{a_k}F,T_{b_k}G\rangle_{\mathcal F^2}
=
\langle F,T_{-a_k}T_{b_k}G\rangle_{\mathcal F^2}
=
e^{i\theta_k}\langle F,T_{b_k-a_k}G\rangle_{\mathcal F^2},
\]
for some real phase \(\theta_k\) coming from the cocycle of the Weyl
representation; since this phase factor has modulus one, the conclusion
follows from \(|b_k-a_k|\to\infty\).
\end{proof}

The next lemma compares the density of a finite sum of widely separated
translates with the density of each individual translate. The precise statement
is what makes the profile decomposition compatible with the variational
quantities \(I_F(s)\), since these depend only on the densities and not on
the phases of the underlying functions.

\begin{lemma}\label{lem:pointwise-decoupling-global}
Fix \(J\ge 1\), let \(G^{(1)},\dots,G^{(J)}\in \mathcal F^2(\C)\), and let
\(\{a_k^{(j)}\}_{k\ge 1}\subset \C\) for \(1\le j\le J\) satisfy
\[
|a_k^{(j)}-a_k^{(\ell)}|\to\infty
\qquad\text{whenever }j\neq \ell.
\]
For each \(k\), set
\[
S_k:=\sum_{j=1}^J T_{a_k^{(j)}}G^{(j)}.
\]
Then the following hold.
\begin{enumerate}
\item[(i)] For every \(R>0\) and every \(1\le j\le J\),
\[
\sup_{z\in D_R(a_k^{(j)})}
\big|u_{S_k}(z)-u_{T_{a_k^{(j)}}G^{(j)}}(z)\big|
\to 0
\qquad (k\to\infty).
\]

\item[(ii)] For every \(R>0\), if
\[
E_{k,R}:=\C\setminus \bigcup_{j=1}^J D_R(a_k^{(j)}),
\]
then
\[
\limsup_{k\to\infty}\sup_{z\in E_{k,R}}u_{S_k}(z)
\le
\left(
\sum_{j=1}^J
\sup_{\zeta\in \C\setminus D_R} u_{G^{(j)}}(\zeta)^{1/2}
\right)^2.
\]
In particular, given \(\varepsilon>0\), there exists \(R_\varepsilon>0\) such that
for every \(R\ge R_\varepsilon\),
\[
\sup_{z\in E_{k,R}}u_{S_k}(z)\le \varepsilon
\]
for all sufficiently large \(k\).
\end{enumerate}
\end{lemma}

\begin{proof}
Throughout the proof we use repeatedly the identity
\[
u_{T_a F}(z)=u_F(z-a),
\qquad z,a\in \C,
\]
which is \eqref{eq:u-translation-global} (equivalently
\eqref{eq:uF-translate}); together with the elementary subadditivity
\[
u_{F+G}(z)^{1/2}\le u_F(z)^{1/2}+u_G(z)^{1/2},
\qquad z\in\C,
\]
which follows from the definition \(u_F=|F|^2 e^{-\pi |z|^2}\) and the
triangle inequality applied to \(|F+G|\le |F|+|G|\).

\medskip
\noindent\emph{Proof of (i).}
Fix \(R>0\) and \(j\in\{1,\dots,J\}\). Write
\[
S_k=T_{a_k^{(j)}}G^{(j)}+H_{k,j},
\qquad
H_{k,j}:=\sum_{\ell\neq j} T_{a_k^{(\ell)}}G^{(\ell)}.
\]
For every \(z\in \C\), the elementary inequality
\(\big||a+b|^2-|a|^2\big|\le 2|a|\,|b|+|b|^2\) yields
\begin{align*}
\big|u_{S_k}(z)-u_{T_{a_k^{(j)}}G^{(j)}}(z)\big|
&=
e^{-\pi |z|^2}
\big||T_{a_k^{(j)}}G^{(j)}(z)+H_{k,j}(z)|^2
-|T_{a_k^{(j)}}G^{(j)}(z)|^2\big| \\
&\le
2\,u_{T_{a_k^{(j)}}G^{(j)}}(z)^{1/2}u_{H_{k,j}}(z)^{1/2}
+u_{H_{k,j}}(z).
\end{align*}
By Lemma~\ref{lem:Fock-basic}(ii),
\[
u_{T_{a_k^{(j)}}G^{(j)}}(z)\le \|G^{(j)}\|_{\mathcal F^2}^2
\qquad \text{for every }z\in\C\text{ and every }k.
\]
Thus, to prove (i) it is enough to show that
\[
\sup_{z\in D_R(a_k^{(j)})} u_{H_{k,j}}(z)\to 0.
\]

For \(z\in D_R(a_k^{(j)})\) and \(\ell\neq j\),
\[
z-a_k^{(\ell)}\in D_R(a_k^{(j)}-a_k^{(\ell)}).
\]
Since \(|a_k^{(j)}-a_k^{(\ell)}|\to\infty\),
\[
\operatorname{dist}\big(D_R(a_k^{(j)}-a_k^{(\ell)}),0\big)
\ge |a_k^{(j)}-a_k^{(\ell)}|-R\to\infty.
\]
Because \(u_{G^{(\ell)}}\in C_0(\C)\) by Lemma~\ref{lem:Fock-basic}(iii),
\[
\sup_{z\in D_R(a_k^{(j)})} u_{T_{a_k^{(\ell)}}G^{(\ell)}}(z)
=
\sup_{\zeta\in D_R(a_k^{(j)}-a_k^{(\ell)})} u_{G^{(\ell)}}(\zeta)
\xrightarrow[k\to\infty]{} 0.
\]
The subadditivity of \(u_F^{1/2}\) yields
\[
u_{H_{k,j}}(z)^{1/2}
\le
\sum_{\ell\neq j} u_{T_{a_k^{(\ell)}}G^{(\ell)}}(z)^{1/2},
\]
so the preceding convergence implies
\[
\sup_{z\in D_R(a_k^{(j)})} u_{H_{k,j}}(z)^{1/2}\to 0.
\]
This proves (i).

\medskip
\noindent\emph{Proof of (ii).}
Fix \(R>0\) and \(z\in E_{k,R}\). By definition, \(z\notin D_R(a_k^{(j)})\)
for every \(j\), and hence \(|z-a_k^{(j)}|\ge R\). Therefore
\[
u_{T_{a_k^{(j)}}G^{(j)}}(z)
=
u_{G^{(j)}}(z-a_k^{(j)})
\le
\sup_{\zeta\in \C\setminus D_R}u_{G^{(j)}}(\zeta).
\]
Using again the subadditivity of \(u_F^{1/2}\),
\[
u_{S_k}(z)^{1/2}
\le
\sum_{j=1}^J u_{T_{a_k^{(j)}}G^{(j)}}(z)^{1/2}
\le
\sum_{j=1}^J \sup_{\zeta\in \C\setminus D_R}u_{G^{(j)}}(\zeta)^{1/2}.
\]
Squaring yields
\[
u_{S_k}(z)
\le
\left(
\sum_{j=1}^J \sup_{\zeta\in \C\setminus D_R}u_{G^{(j)}}(\zeta)^{1/2}
\right)^2.
\]
Taking the supremum over \(z\in E_{k,R}\) proves the displayed estimate in
(ii). Since each \(u_{G^{(j)}}\in C_0(\C)\), the right-hand side tends to
\(0\) as \(R\to\infty\); the final claim follows by choosing
\(R_\varepsilon\) so that the displayed bracket is at most
\(\sqrt{\varepsilon}\).
\end{proof}

The next lemma is a robust uniform-density perturbation estimate. Concretely,
it says that adding to a bounded sequence in \(\mathcal F^2(\C)\) a residue
whose density vanishes uniformly leaves the density essentially unchanged in
\(L^\infty\). This will be applied with \(S_k\) the truncated profile sum and
\(R_k\) the residual of the profile decomposition.

\begin{lemma}\label{lem:density-perturbation-global}
Let \(\{S_k\}\subset \mathcal F^2(\C)\) be bounded and let
\(\{R_k\}\subset \mathcal F^2(\C)\) satisfy
\[
\|u_{R_k}\|_{L^\infty(\C)}\to 0.
\]
Then
\[
\|u_{S_k+R_k}-u_{S_k}\|_{L^\infty(\C)}\to 0.
\]
More precisely, if \(\sup_k \|S_k\|_{\mathcal F^2}\le M\), then
\[
\|u_{S_k+R_k}-u_{S_k}\|_{L^\infty(\C)}
\le
2M\,\|u_{R_k}\|_{L^\infty(\C)}^{1/2}
+\|u_{R_k}\|_{L^\infty(\C)}.
\]
\end{lemma}

\begin{proof}
For every \(z\in \C\),
\begin{align*}
\big|u_{S_k+R_k}(z)-u_{S_k}(z)\big|
&=
e^{-\pi |z|^2}\big||S_k(z)+R_k(z)|^2-|S_k(z)|^2\big| \\
&\le
2\,u_{S_k}(z)^{1/2}u_{R_k}(z)^{1/2}
+u_{R_k}(z),
\end{align*}
again by the elementary inequality
\(\big||a+b|^2-|a|^2\big|\le 2|a|\,|b|+|b|^2\). By
Lemma~\ref{lem:Fock-basic}(ii),
\[
u_{S_k}(z)\le \|S_k\|_{\mathcal F^2}^2\le M^2
\qquad (z\in\C).
\]
Hence
\[
\big|u_{S_k+R_k}(z)-u_{S_k}(z)\big|
\le
2M\,\|u_{R_k}\|_{L^\infty}^{1/2}
+\|u_{R_k}\|_{L^\infty}.
\]
Taking the supremum over \(z\) yields the claim.
\end{proof}

The next lemma supplies the soft input on the function \(s\mapsto I_F(s)\)
that will be used at the end of the variational argument. It is essentially a
consequence of the bathtub principle and the basic regularity of \(u_F\),
already noted in Lemma~\ref{lem:basic-properties-global}.

\begin{lemma}\label{lem:I-continuity-global}
Let \(F\in \mathcal F^2(\C)\). Then the function
\[
s\mapsto I_F(s)
\]
is finite, nondecreasing, and continuous on \([0,\infty)\).
\end{lemma}

\begin{proof}
By Lemma~\ref{lem:Fock-basic}(ii)--(iii) together with the identity
\(\int_\C u_F\,dA=\|F\|_{\mathcal F^2}^2\), we have
\(u_F\in C_0(\C)\cap L^1(\C)\), so \(I_F(s)\) is finite for every \(s\ge 0\)
and bounded above by \(\|F\|_{\mathcal F^2}^2\). Monotonicity is immediate
from the definition of \(I_F(s)\) as a supremum over \(|\Omega|=s\): if
\(s_1\le s_2\), every admissible set for \(s_1\) can be enlarged to a set of
measure \(s_2\) without decreasing the integral, since \(u_F\ge 0\).

To prove continuity, write \(u:=u_F\) and let \(u^*\) denote the decreasing
rearrangement of \(u\) on \([0,\infty)\). Since \(u\in C_0(\C)\cap L^1(\C)\),
\(u^*\in L^1([0,\infty))\) and is equimeasurable with \(u\). The bathtub
principle yields
\[
I_F(s)=\int_0^s u^*(\tau)\,d\tau,\qquad s\ge 0.
\]
The right-hand side is the integral of an
\(L^1_{\mathrm{loc}}([0,\infty))\) function over \([0,s]\) and is therefore
absolutely continuous, hence continuous, in \(s\).
\end{proof}

The next result is the profile decomposition adapted to the present setting,
in the spirit of P.-L.~Lions' concentration-compactness method. Although the
overall scheme is by now standard, we spell out the points used in the
sequel; the version stated here is a Fock-space variant of profile
decompositions used in time-frequency analysis, see e.g.\
\cite{Nicola-Romero-Trapasso, Marceca-Romero-Speckbacher, Fulsche-Hagger,
Samuelsen}.

\begin{proposition}[Profile decomposition in the Fock space]\label{prop:profile-global}
Let \(\{F_k\}_{k\ge 1}\subset \mathcal F^2(\C)\) be a bounded sequence. Then,
after passing to a subsequence, there exist profiles
\[
G^{(j)}\in \mathcal F^2(\C),
\qquad j\ge 1,
\]
and sequences of points
\[
a_k^{(j)}\in \C,
\qquad j\ge 1,\ k\ge 1,
\]
such that:
\begin{enumerate}
    \item[(i)] for each \(j\neq \ell\),
    \[
    |a_k^{(j)}-a_k^{(\ell)}|\to \infty
    \qquad (k\to \infty);
    \]
    \item[(ii)] for every \(J\ge 1\), if
    \[
    R_k^{(J)}
    :=
    F_k-\sum_{j=1}^J T_{a_k^{(j)}}G^{(j)},
    \]
    then
    \[
    T_{-a_k^{(j)}}R_k^{(J)}\to 0
    \qquad\text{locally uniformly on }\C
    \]
    for every \(1\le j\le J\);
    \item[(iii)] for every \(J\ge 1\),
    \[
    \|F_k\|_{\mathcal F^2}^2
    =
    \sum_{j=1}^J \|G^{(j)}\|_{\mathcal F^2}^2
    +\|R_k^{(J)}\|_{\mathcal F^2}^2
    +o_k(1);
    \]
    \item[(iv)]
    \[
    \lim_{J\to\infty}\left(\limsup_{k\to\infty}
    \|u_{R_k^{(J)}}\|_{L^\infty(\C)}\right)=0.
    \]
\end{enumerate}
\end{proposition}

\begin{proof}
We argue by the standard iterative extraction scheme, but spell out the points
used in the rest of the section. Throughout, we say that a function in
\(\mathcal F^2(\C)\) is \emph{nontrivial} if it is not identically zero.

\medskip
\noindent\emph{Extraction of the first profile.}
If
\[
\limsup_{k\to\infty}\|u_{F_k}\|_{L^\infty(\C)}=0,
\]
there is nothing to prove: take all profiles equal to zero, all centers
arbitrary (say zero), and observe that (i)--(iii) hold trivially while (iv)
holds by hypothesis. Otherwise, after passing to a subsequence, choose
\(a_k^{(1)}\in \C\) such that
\[
u_{F_k}(a_k^{(1)})\ge \tfrac12\|u_{F_k}\|_{L^\infty(\C)}.
\]
By \eqref{eq:I-translation-global} and the density-translation identity
\eqref{eq:u-translation-global}, the sequence \(T_{-a_k^{(1)}}F_k\) has the
same Fock norm and the corresponding density translated; relabelling, we may
assume \(a_k^{(1)}=0\). Since \(\{F_k\}\) is bounded in \(\mathcal F^2(\C)\),
the reproducing-kernel bound \(|F_k(z)|\le \|F_k\|_{\mathcal F^2}\,
e^{\pi |z|^2/2}\) (Lemma~\ref{lem:growth}(ii)) makes \(\{F_k\}\) a normal
family of entire functions on every compact subset of \(\C\). Passing to a
subsequence and applying Montel's theorem, \(F_k\to G^{(1)}\) locally
uniformly on \(\C\), with \(G^{(1)}\in\mathcal F^2(\C)\) by Fatou's lemma. By
the choice of \(a_k^{(1)}=0\),
\[
u_{G^{(1)}}(0)=\lim_{k\to\infty}u_{F_k}(0)
\ge \tfrac12 \limsup_{k\to\infty}\|u_{F_k}\|_{L^\infty(\C)}>0,
\]
so \(G^{(1)}\) is nontrivial. Set
\[
R_k^{(1)}:=F_k-G^{(1)}.
\]

\medskip
\noindent\emph{Iteration.}
Suppose profiles \(G^{(1)},\dots,G^{(J-1)}\) and centers
\(a_k^{(1)},\dots,a_k^{(J-1)}\) have been extracted, with remainder
\(R_k^{(J-1)}\). If
\[
\limsup_{k\to\infty}\|u_{R_k^{(J-1)}}\|_{L^\infty(\C)}=0,
\]
we stop, set \(G^{(j)}=0\) and \(a_k^{(j)}=0\) for all \(j\ge J\), and
properties (i)--(iv) are easily seen to hold. Otherwise choose
\(a_k^{(J)}\in\C\) so that
\[
u_{R_k^{(J-1)}}(a_k^{(J)})
\ge \tfrac12\|u_{R_k^{(J-1)}}\|_{L^\infty(\C)}.
\]
Property (i) at this stage (proved below) forces
\(|a_k^{(J)}-a_k^{(m)}|\to\infty\) for every \(m<J\). Granted this, the
sequence \(\{T_{-a_k^{(J)}}R_k^{(J-1)}\}\) is bounded in \(\mathcal F^2\) (by
unitarity of \(T_{-a_k^{(J)}}\) together with (iii) at stage \(J-1\)), so
after passing to a further subsequence and applying Montel's theorem, it
converges locally uniformly to some \(G^{(J)}\in\mathcal F^2(\C)\). By the
choice of \(a_k^{(J)}\),
\[
u_{G^{(J)}}(0)
\ge \tfrac12 \limsup_{k\to\infty}\|u_{R_k^{(J-1)}}\|_{L^\infty(\C)}>0,
\]
so \(G^{(J)}\) is nontrivial. Set
\[
R_k^{(J)}:=R_k^{(J-1)}-T_{a_k^{(J)}}G^{(J)}.
\]
Iterating either terminates at some finite step (in which case (iv) is
immediate) or produces an infinite sequence of nontrivial profiles. By a
standard diagonal extraction we obtain a single subsequence (still denoted
\(\{F_k\}\)) for which all the partial decompositions are valid.

\medskip
\noindent\emph{Proof of (i).}
Fix \(j<\ell\), and suppose for contradiction that
\(|a_k^{(\ell)}-a_k^{(j)}|\) does not tend to infinity. After passing to a
subsequence, \(|a_k^{(\ell)}-a_k^{(j)}|\le C\) for some constant \(C\). At
the \(\ell\)-th step, by construction
\[
T_{-a_k^{(\ell)}}R_k^{(\ell-1)}\to G^{(\ell)}
\qquad\text{locally uniformly on }\C,
\]
with \(G^{(\ell)}\) nontrivial. On the other hand, property (ii) verified up
to stage \(\ell-1\) (proved below) gives
\[
T_{-a_k^{(j)}}R_k^{(\ell-1)}\to 0
\qquad\text{locally uniformly on }\C.
\]
The Weyl identity
\(T_{-a_k^{(\ell)}}=e^{i\theta_k}T_{-a_k^{(j)}}T_{a_k^{(j)}-a_k^{(\ell)}}\)
(for an appropriate real phase \(\theta_k\) coming from the Weyl co-cycle, of
modulus one) shows that
\[
T_{-a_k^{(\ell)}}R_k^{(\ell-1)}(z)
=
e^{i\theta_k}\big(T_{-a_k^{(j)}}R_k^{(\ell-1)}\big)\!
\big(z-(a_k^{(\ell)}-a_k^{(j)})\big).
\]
Since \(|a_k^{(\ell)}-a_k^{(j)}|\le C\), the points
\(z-(a_k^{(\ell)}-a_k^{(j)})\) remain in a fixed compact set as \(z\)
varies over a compact set; so the local uniform convergence of
\(T_{-a_k^{(j)}}R_k^{(\ell-1)}\) to \(0\) forces
\(T_{-a_k^{(\ell)}}R_k^{(\ell-1)}\to 0\) locally uniformly as well.
This contradicts the nontriviality of \(G^{(\ell)}\). Hence
\[
|a_k^{(j)}-a_k^{(\ell)}|\to\infty
\qquad\text{whenever }j\neq \ell.
\]

\medskip
\noindent\emph{Proof of (ii).}
Fix \(J\ge 1\) and \(1\le j\le J\). By construction, at the \(j\)-th
extraction step,
\[
T_{-a_k^{(j)}}R_k^{(j-1)}\to G^{(j)}
\qquad\text{locally uniformly on }\C.
\]
Subtracting the \(j\)-th extracted profile,
\[
T_{-a_k^{(j)}}R_k^{(j)}
=
T_{-a_k^{(j)}}R_k^{(j-1)}-G^{(j)}
\to 0
\qquad\text{locally uniformly on }\C.
\]
For \(J>j\), since
\[
R_k^{(J)}
=
R_k^{(j)}-\sum_{\ell=j+1}^J T_{a_k^{(\ell)}}G^{(\ell)},
\]
the linearity of the Weyl operators yields
\[
T_{-a_k^{(j)}}R_k^{(J)}
=
T_{-a_k^{(j)}}R_k^{(j)}
-\sum_{\ell=j+1}^J e^{i\theta_{k,\ell}}\,T_{a_k^{(\ell)}-a_k^{(j)}}G^{(\ell)}
\]
for some real phases \(\theta_{k,\ell}\) coming from the Weyl co-cycle. The
first term tends to \(0\) locally uniformly by the previous paragraph. For
each fixed \(\ell>j\), part (i) gives \(|a_k^{(\ell)}-a_k^{(j)}|\to\infty\),
and the identity \(u_{T_b G}(z)=u_G(z-b)\) (i.e.\
\eqref{eq:u-translation-global}) implies that
\(T_{a_k^{(\ell)}-a_k^{(j)}}G^{(\ell)}\to 0\) locally uniformly in modulus
(its density is \(u_{G^{(\ell)}}(z-(a_k^{(\ell)}-a_k^{(j)}))\to 0\) for
\(z\) in a compact set, since \(u_{G^{(\ell)}}\in C_0(\C)\)). Thus
\[
T_{-a_k^{(j)}}R_k^{(J)}\to 0
\qquad\text{locally uniformly on }\C,
\]
which is exactly (ii).

\medskip
\noindent\emph{Proof of (iii).}
For fixed \(J\ge 1\),
\[
F_k=\sum_{j=1}^J T_{a_k^{(j)}}G^{(j)}+R_k^{(J)}.
\]
Expanding the squared \(\mathcal F^2\)-norm and using
Lemma~\ref{lem:weyl-orthogonality-global} together with the pairwise
divergence of the translation parameters from (i),
\[
\Big\|\sum_{j=1}^J T_{a_k^{(j)}}G^{(j)}\Big\|_{\mathcal F^2}^2
=
\sum_{j=1}^J \|G^{(j)}\|_{\mathcal F^2}^2+o_k(1).
\]
Moreover, since
\[
T_{-a_k^{(m)}}R_k^{(J)}\to 0
\qquad\text{locally uniformly on }\C
\]
for every \(1\le m\le J\) by (ii), and the sequence
\(\{T_{-a_k^{(m)}}R_k^{(J)}\}_k\) is bounded in \(\mathcal F^2\), the same
weak-convergence argument as in
Lemma~\ref{lem:weyl-orthogonality-global} yields
\[
\langle T_{a_k^{(m)}}G^{(m)},R_k^{(J)}\rangle_{\mathcal F^2}
=
\langle G^{(m)},T_{-a_k^{(m)}}R_k^{(J)}\rangle_{\mathcal F^2}\to 0,
\]
for every fixed \(m\). Expanding \(\|F_k\|_{\mathcal F^2}^2\) and collecting
terms,
\[
\|F_k\|_{\mathcal F^2}^2
=
\sum_{j=1}^J \|G^{(j)}\|_{\mathcal F^2}^2
+\|R_k^{(J)}\|_{\mathcal F^2}^2
+o_k(1).
\]
In particular, since \(\|F_k\|_{\mathcal F^2}\) is bounded and the terms on
the right are nonnegative, summing the profile-norm contributions over \(j\)
yields the energy bound
\begin{equation}\label{eq:energy-bound-profiles}
\sum_{j\ge 1}\|G^{(j)}\|_{\mathcal F^2}^2
\le \limsup_{k\to\infty}\|F_k\|_{\mathcal F^2}^2.
\end{equation}
This is the key compactness ingredient that controls the iteration and
ensures (iv).

\medskip
\noindent\emph{Proof of (iv).}
Set
\[
\eta_J:=\limsup_{k\to\infty}\|u_{R_k^{(J)}}\|_{L^\infty(\C)},
\]
which is monotone nonincreasing in \(J\) by construction (since
\(R_k^{(J+1)}\) is obtained from \(R_k^{(J)}\) by extracting one further
nontrivial profile). Let \(\eta_\infty:=\lim_{J\to\infty}\eta_J\); we must
show \(\eta_\infty=0\). Suppose, for contradiction, that
\(\eta_\infty\ge \delta>0\).

By the iteration, at each stage \(J\) such that \(\eta_J\ge \delta\), the
extracted profile \(G^{(J+1)}\) satisfies \(u_{G^{(J+1)}}(0)\ge \delta/2\) by
construction. By Lemma~\ref{lem:Fock-basic}(ii),
\[
\|G^{(J+1)}\|_{\mathcal F^2}^2 \ge u_{G^{(J+1)}}(0)\ge \delta/2.
\]
If the iteration produced infinitely many such stages (which is the case
under the assumption \(\eta_\infty\ge \delta\)), then
\[
\sum_{j\ge 1}\|G^{(j)}\|_{\mathcal F^2}^2=\infty,
\]
contradicting \eqref{eq:energy-bound-profiles}. Hence \(\eta_\infty=0\),
proving (iv).
\end{proof}

We now turn to the global problem. The proof of existence consists of
combining the profile decomposition with the elementary continuity properties
of \(I_F(s)\), and using the homogeneity \(I_{cF}(s)=|c|^2 I_F(s)\) recorded
in Lemma~\ref{lem:basic-properties-global}(ii) to redistribute mass to a
single nontrivial profile.

\begin{theorem}\label{thm:existence-global-maximizer}
For every \(s>0\), there exists \(F_*\in \mathcal F^2(\C)\) with
\(\|F_*\|_{\mathcal F^2}=1\) such that
\[
I_{F_*}(s)=M(s).
\]
Equivalently, there exists a measurable set \(\Omega_*\subset \C\) with
\(|\Omega_*|=s\) such that
\[
\lambda_{1,\varphi}(\Omega_*)=M(s).
\]
\end{theorem}

\begin{proof}
Let \(\{F_k\}\subset \mathcal F^2(\C)\) be a maximizing sequence:
\[
\|F_k\|_{\mathcal F^2}=1,
\qquad
I_{F_k}(s)\to M(s).
\]
By replacing \(F_k\) with a Weyl translate if necessary, and using the
translation invariance \eqref{eq:I-translation-global} together with
\eqref{eq:u-translation-global}, we may suppose that
\[
u_{F_k}(0)=\|u_{F_k}\|_{L^\infty(\C)}.
\]
Note that \(M(s)>0\): indeed, taking \(F\equiv 1\), one has
\(F\in \mathcal F^2(\C)\) with \(\|F\|_{\mathcal F^2}^2=1\) and density
\(u_F(z)=e^{-\pi |z|^2}>0\) everywhere, so for any \(s>0\), \(I_F(s)>0\).
Consequently, \(\|u_{F_k}\|_{L^\infty(\C)}\) is bounded below by a positive
constant uniformly in \(k\); otherwise, for any \(\eta>0\) we would have
\(\|u_{F_k}\|_{L^\infty}<\eta\) eventually, hence \(I_{F_k}(s)\le \eta s\),
contradicting \(I_{F_k}(s)\to M(s)>0\).

Applying Proposition~\ref{prop:profile-global}, after passing to a
subsequence we may write
\[
F_k=\sum_{j=1}^J T_{a_k^{(j)}}G^{(j)}+R_k^{(J)}
\]
for every fixed \(J\ge 1\), with the conclusions of that proposition.

\medskip
\noindent\emph{Step 1: an upper bound for the limit.}
Fix \(J\ge 1\) and let \(\Omega_k^{(J)}\) be a measurable set of measure
\(s\) such that
\[
I_{F_k}(s)=\int_{\Omega_k^{(J)}} u_{F_k}\,dA;
\]
by Lemma~\ref{lem:basic-properties-global}(i) we may take
\(\Omega_k^{(J)}\) to be a superlevel set of \(u_{F_k}\). Let
\(\varepsilon>0\). Since \(u_{G^{(j)}}\in L^1(\C)\) for every \(j\) (with
\(\int_\C u_{G^{(j)}}\,dA=\|G^{(j)}\|_{\mathcal F^2}^2\)), choose \(R>0\) so
large that
\[
\int_{\C\setminus D_R} u_{G^{(j)}}\,dA<\varepsilon
\qquad\text{for all } 1\le j\le J,
\]
and large enough that the conclusion of
Lemma~\ref{lem:pointwise-decoupling-global}(ii) holds with this choice of
\(R\). Since the translation parameters are pairwise divergent by
Proposition~\ref{prop:profile-global}(i), for \(k\) large enough the disks
\[
D_R(a_k^{(1)}),\dots,D_R(a_k^{(J)})
\]
are pairwise disjoint. Write
\[
B_{k,j}:=D_R(a_k^{(j)}),
\qquad
E_{k,j}:=\Omega_k^{(J)}\cap B_{k,j},
\qquad
E_{k,0}:=\Omega_k^{(J)}\setminus \bigcup_{j=1}^J B_{k,j}.
\]
Set
\[
s_{k,j}:=|E_{k,j}|,
\qquad 0\le j\le J.
\]
Then, since the \(\{B_{k,j}\}_{j=1}^J\) are disjoint and
\(\Omega_k^{(J)}\) has measure \(s\),
\[
\sum_{j=0}^J s_{k,j}=s.
\]

Set
\[
S_k^{(J)}:=\sum_{j=1}^J T_{a_k^{(j)}}G^{(j)}.
\]
By Proposition~\ref{prop:profile-global}(iii), the sequence
\(\{S_k^{(J)}\}_k\) is bounded in \(\mathcal F^2(\C)\), with bound
\(\le 1+o_k(1)\). By Proposition~\ref{prop:profile-global}(iv), if \(J\) is
taken sufficiently large then
\[
\limsup_{k\to\infty}\|u_{R_k^{(J)}}\|_{L^\infty(\C)}<\varepsilon^2.
\]
Lemma~\ref{lem:density-perturbation-global} therefore gives
\[
\|u_{F_k}-u_{S_k^{(J)}}\|_{L^\infty(\C)}\to 0
\qquad (k\to\infty),
\]
and in particular, after taking \(k\) large enough,
\begin{equation}\label{eq:density-approximation-global}
\|u_{F_k}-u_{S_k^{(J)}}\|_{L^\infty(\C)}<\varepsilon.
\end{equation}

Now
\[
E_{k,0}
=
\Omega_k^{(J)}\setminus \bigcup_{j=1}^J B_{k,j}
\subset
\C\setminus \bigcup_{j=1}^J D_R(a_k^{(j)}).
\]
Lemma~\ref{lem:pointwise-decoupling-global}(i) gives, for each fixed
\(j\in\{1,\dots,J\}\),
\[
\sup_{z\in B_{k,j}}
\big|u_{S_k^{(J)}}(z)-u_{T_{a_k^{(j)}}G^{(j)}}(z)\big|\to 0
\qquad (k\to\infty),
\]
while Lemma~\ref{lem:pointwise-decoupling-global}(ii) yields, for our
choice of \(R\) and for \(k\) sufficiently large,
\[
\sup_{z\in E_{k,0}}u_{S_k^{(J)}}(z)\le \varepsilon.
\]
Combining these bounds with \eqref{eq:density-approximation-global}, we
conclude that for all sufficiently large \(k\),
\begin{enumerate}
\item[(a)] on each \(B_{k,j}\),
\[
u_{F_k}(z)\le u_{T_{a_k^{(j)}}G^{(j)}}(z)+2\varepsilon;
\]
\item[(b)] on \(E_{k,0}\),
\[
u_{F_k}(z)\le 2\varepsilon.
\]
\end{enumerate}

Therefore, decomposing the integral defining \(I_{F_k}(s)\) according to the
partition \(\Omega_k^{(J)}=E_{k,0}\sqcup\bigsqcup_{j=1}^J E_{k,j}\),
\begin{align*}
I_{F_k}(s)
&=
\int_{\Omega_k^{(J)}} u_{F_k}\,dA \\
&\le
\sum_{j=1}^J \int_{E_{k,j}} u_{T_{a_k^{(j)}}G^{(j)}}\,dA
+2\varepsilon \sum_{j=0}^J |E_{k,j}| \\
&=
\sum_{j=1}^J \int_{E_{k,j}-a_k^{(j)}} u_{G^{(j)}}\,dA
+2\varepsilon s \\
&\le
\sum_{j=1}^J I_{G^{(j)}}(s_{k,j})
+2\varepsilon s,
\end{align*}
where in the next-to-last step we used the change of variables
\(z\mapsto z-a_k^{(j)}\) together with \eqref{eq:u-translation-global}, and
in the last step the definition of \(I_{G^{(j)}}(s_{k,j})\) as the supremum
of \(\int_E u_{G^{(j)}}\,dA\) over sets of measure
\(s_{k,j}=|E_{k,j}|=|E_{k,j}-a_k^{(j)}|\).

Passing to a further subsequence if necessary, we may suppose
\[
s_{k,j}\to \sigma_j\in [0,s]
\qquad (k\to\infty)
\]
for each \(1\le j\le J\). By the continuity of \(I_{G^{(j)}}\) in the
measure parameter, given by Lemma~\ref{lem:I-continuity-global},
\[
M(s)=\lim_{k\to\infty}I_{F_k}(s)
\le
\sum_{j=1}^J I_{G^{(j)}}(\sigma_j)+2\varepsilon s.
\]
Since \(\varepsilon>0\) was arbitrary,
\begin{equation}\label{eq:M-upper-profiles-global}
M(s)\le \sum_{j=1}^J I_{G^{(j)}}(\sigma_j).
\end{equation}

\medskip
\noindent\emph{Step 2: redistribution of mass and saturation of the bound.}
We now use Lemma~\ref{lem:basic-properties-global}(ii). For each \(j\) with
\(G^{(j)}\neq 0\), set
\[
\widehat G^{(j)}:=\frac{G^{(j)}}{\|G^{(j)}\|_{\mathcal F^2}},
\qquad
\|\widehat G^{(j)}\|_{\mathcal F^2}=1,
\]
and use the homogeneity \(I_{cF}(\sigma)=|c|^2 I_F(\sigma)\) (with
\(c=\|G^{(j)}\|_{\mathcal F^2}\)) to obtain
\[
I_{G^{(j)}}(\sigma_j)
=
\|G^{(j)}\|_{\mathcal F^2}^2
\,I_{\widehat G^{(j)}}(\sigma_j).
\]
Since \(\widehat G^{(j)}\) has unit \(\mathcal F^2\)-norm, by definition of
\(M\) and the monotonicity of \(I_{\widehat G^{(j)}}\) given by
Lemma~\ref{lem:I-continuity-global} (using \(\sigma_j\le s\)),
\[
I_{\widehat G^{(j)}}(\sigma_j)\le I_{\widehat G^{(j)}}(s)\le M(s).
\]
Hence
\[
I_{G^{(j)}}(\sigma_j)\le \|G^{(j)}\|_{\mathcal F^2}^2\, M(s).
\]
For \(j\) with \(G^{(j)}=0\) the inequality is trivial. Summing over
\(1\le j\le J\) and using Proposition~\ref{prop:profile-global}(iii)
together with the normalization \(\|F_k\|_{\mathcal F^2}^2=1\),
\[
\sum_{j=1}^J I_{G^{(j)}}(\sigma_j)
\le
M(s)\sum_{j=1}^J \|G^{(j)}\|_{\mathcal F^2}^2
\le M(s).
\]
Combined with \eqref{eq:M-upper-profiles-global}, this yields
\[
M(s)\le \sum_{j=1}^J I_{G^{(j)}}(\sigma_j)\le M(s)
\]
for every \(J\). Therefore equality holds throughout for every \(J\), which
forces
\[
\sum_{j\ge 1}\|G^{(j)}\|_{\mathcal F^2}^2=1
\]
(so the residual \(R_k^{(J)}\) carries no \(\mathcal F^2\)-mass in the
limit), and moreover, for every \(j\) with \(G^{(j)}\neq 0\),
\[
I_{\widehat G^{(j)}}(\sigma_j)=M(s).
\]
By monotonicity of \(I_{\widehat G^{(j)}}\),
\(I_{\widehat G^{(j)}}(s)\ge I_{\widehat G^{(j)}}(\sigma_j)=M(s)\); since
\(M(s)\) is the supremum over unit-norm functions, we conclude
\[
I_{\widehat G^{(j)}}(s)=M(s).
\]
Thus every nonzero normalized profile is a global maximizer.

\medskip
\noindent\emph{Step 3: producing a maximizing set.}
Choose any \(j_0\) with \(G^{(j_0)}\neq 0\) (such \(j_0\) exists since
\(\sum_{j\ge 1}\|G^{(j)}\|_{\mathcal F^2}^2=1\)) and set
\(F_*:=\widehat G^{(j_0)}\), so that
\[
\|F_*\|_{\mathcal F^2}=1,
\qquad
I_{F_*}(s)=M(s).
\]
By Lemma~\ref{lem:basic-properties-global}(i), there exists a superlevel
set \(\Omega_*\) of \(u_{F_*}\) with \(|\Omega_*|=s\) and
\[
\int_{\Omega_*}u_{F_*}\,dA=I_{F_*}(s)=M(s).
\]
Then
\[
\lambda_{1,\varphi}(\Omega_*)\ge \int_{\Omega_*}u_{F_*}\,dA=M(s),
\]
while by definition of \(M(s)\) we have
\(\lambda_{1,\varphi}(\Omega_*)\le M(s)\). Hence
\[
\lambda_{1,\varphi}(\Omega_*)=M(s),
\]
and \(\Omega_*\) is a global maximizer of measure \(s\).
\end{proof}

We may now identify the global maximizers, by combining the existence theorem
above with the local-rigidity theorem of the previous section.

\begin{theorem}\label{thm:global-maximizers-are-disks}
For every \(s>0\), every global maximizer for
\[
\sup_{|\Omega|=s}\lambda_{1,\varphi}(\Omega)
\]
is, up to translation, a disk.
\end{theorem}

\begin{proof}
Let \(\Omega_*\) be a global maximizer of measure \(s\) (which exists by
Theorem~\ref{thm:existence-global-maximizer}). Then \(\Omega_*\) is, in
particular, a local maximizer of \(\lambda_{1,\varphi}\) under the same
measure constraint: every sufficiently small smooth deformation of
\(\Omega_*\) preserving the measure \(|\Omega|=s\) yields a competitor
whose value of \(\lambda_{1,\varphi}\) is at most
\(M(s)=\lambda_{1,\varphi}(\Omega_*)\). In particular, every competitor
\(E\) with \(|E|=|\Omega_*|\) and Gaussian-weighted symmetric-difference
distance \(\int_\C |\mathbf 1_{\Omega_*}-\mathbf 1_E|e^{-\pi|z|^2}\,dz<\delta_0\)
to \(\Omega_*\) (the gauge used in the local-maximizer definition of
Section~\ref{sec:local-max}) satisfies
\(\lambda_{1,\varphi}(E)\le \lambda_{1,\varphi}(\Omega_*)\), because
\(\lambda_{1,\varphi}(E)\le M(s)=\lambda_{1,\varphi}(\Omega_*)\) by the
defining property of \(M(s)\). The local-rigidity theorem proved in
the previous section therefore applies: every such local maximizer is, up
to translation, a disk. Since the problem is invariant under Weyl
translations, this characterization is sharp.
\end{proof}

As an immediate consequence, we recover the classical Nicola--Tilli theorem
\cite{Nicola-Tilli, Nicola-Tilli-2}.

\begin{corollary}
Among all measurable sets of prescribed measure, the disk maximizes Gaussian
time-frequency concentration:
\[
\lambda_{1,\varphi}(\Omega)\le \lambda_{1,\varphi}(D)
\]
whenever \(|\Omega|=|D|\) and \(D\) is a disk.
\end{corollary}

\begin{proof}
Existence of a global maximizer is given by
Theorem~\ref{thm:existence-global-maximizer}, and
Theorem~\ref{thm:global-maximizers-are-disks} shows that every such
maximizer is a disk, up to translation. Since translation by an element of
\(\C\) preserves both Lebesgue measure and the value of
\(\lambda_{1,\varphi}\), the claim follows.
\end{proof}



\bibliography{biblio}
\bibliographystyle{amsplain}

\end{document}
